\documentclass{article}
\usepackage{iclr2023_conference,times}

\usepackage{hyperref}
\usepackage{url}
\usepackage{tikz}
\usepackage{algorithm}
\usepackage{algorithmic}
\usepackage{subfig}
\usepackage{graphicx}
\usepackage{amsmath}
\usepackage{amssymb}
\usepackage{mathtools}
\usepackage{amsthm}
\usepackage{bbm}
\usepackage{booktabs}
\usepackage{multirow}

\usepackage{listings}
\usepackage{xcolor}
\usepackage{color}
\definecolor{dkgreen}{rgb}{0,0.6,0}
\definecolor{gray}{rgb}{0.5,0.5,0.5}
\definecolor{mauve}{rgb}{0.58,0,0.82}

\usepackage[utf8]{inputenc}
\usepackage[T1]{fontenc}
\usepackage{amsfonts}
\usepackage{nicefrac}
\usepackage{microtype}

\newtheorem{theorem}{Theorem}[section]
\newtheorem{proposition}[theorem]{Proposition}
\newtheorem{lemma}[theorem]{Lemma}
\newtheorem{corollary}[theorem]{Corollary}
\newtheorem{definition}[theorem]{Definition}
\newtheorem{assumption}[theorem]{Assumption}
\newtheorem{remark}[theorem]{Remark}
\newtheorem{example}[theorem]{Example}

\newcommand{\revise}[1]{{\color{magenta}[#1]}}

% Silence hyperref warnings about \textsuperscript and math in section titles / captions.
% hyperref builds PDF bookmarks from these and cannot handle formatting/math commands;
% we tell it to render \textsuperscript{+} as a literal "+" and to ignore math-mode tokens.
\pdfstringdefDisableCommands{%
  \def\textsuperscript#1{#1}%
  \def\mathsurround{}%
  \def\z@{}%
}

\title{Heterogeneous-Type Markov Games with Bayesian Coordination:\\
A Posterior Decoupling Approach}

\author{Anonymous Authors\\
Affiliation\\
\texttt{email}}

\newcommand{\fix}{\marginpar{FIX}}
\newcommand{\new}{\marginpar{NEW}}

\begin{document}

\maketitle

\begin{abstract}
We study online learning in heterogeneous-type Markov games where $n$ agents have latent reward types $\theta_i \in \Theta_i$ and the joint transition kernel is unknown. We propose Hierarchical Posterior Sampling for Markov Games (H-PSMG), a centralized algorithm that samples per-agent type marginals and the transition kernel from a factored Bayesian posterior at each episode and executes the coarse correlated equilibrium of the sampled world.

Our central theoretical contribution is a posterior decoupling theorem: under type-independent transitions and own-type reward locality, the joint Bayesian posterior over $(P, \boldsymbol{\theta})$ factorizes across the transition kernel and the per-agent type marginals at every episode. This factorization simultaneously yields (i) exponential-in-$n$ storage savings, from $O(|\Theta|^n)$ to $O(n|\Theta_i|)$, and (ii) exponential-in-$n$ tightening of the regret bound, from the generic Bayesian-RL guarantee $\tilde O(\sqrt{|\Theta|^n K})$ to $\tilde O(H \sqrt{|\Theta_i| K})$.

Building on the decoupling theorem we establish five further results. \textbf{Regret upper bound:} H-PSMG achieves $\tilde O(H^{3/2}\sqrt{SA_{\mathrm{joint}}K} + H\sqrt{|\Theta_i| K})$ Bayesian regret per agent. \textbf{Rate separation from heuristics:} on a constructed exploration-trap HT-MG, H-PSMG attains $\tilde O(\sqrt{K})$ while MAP-type estimation and type-agnostic posterior sampling suffer $\Omega(K)$ linear regret. \textbf{Exponential burn-in improvement:} we propose H-PSMG\textsuperscript{+}, which augments H-PSMG with a type-discrimination bonus computed from the per-agent factored marginals. The bonus requires the decoupling theorem to be definable and, on a deep-trap HT-MG family, reduces expected burn-in regret from $\Omega(|\Theta_i|^n H)$ to $O(H)$, an exponential-in-$n$ separation against Joint-PSRL. \textbf{Matching lower bound:} any algorithm has $\Omega(\sqrt{H^3 S A_{\mathrm{joint}} K} + \sqrt{|\Theta_i| HK})$ regret on some HT-MG, matching our upper bound modulo a $\sqrt{H}$ gap and log factors. \textbf{Exponential improvement over naive Bayesian RL:} without prior factorization, any algorithm has $\Omega(\sqrt{|\Theta|^n K})$ regret, exponentially worse than ours.

Empirical experiments on LLM substrates validate every prediction of the theory. On the Heterogeneous-Persona Sequential Public Goods Game (HP-SPGG) we construct with $|\Theta_i| = 4$ persona archetypes and four LLM backbones (DeepSeek-V3.2, GPT-5.4-nano, Kimi-K2.6, Llama-4-Maverick), the H-PSMG family separates from the best published LLM-coordination baselines (ECON \citep{yang2025econ}, A-ToM \citep{xu2026atom}) by up to $12\times$ and from type-agnostic baselines by over $25\times$. A $\beta$-sweep produces three regimes (self-decay, self-decay edge, burn-in jump) predicted by Theorem~\ref{thm:hpsmg-plus}; we then directly measure each backbone's posterior concentration speed and confirm the regimes correspond exactly to concentration time, supplying a mechanistic test of the theorem's trigger conditions. Two further theory-driven experiments confirm the predicted scaling: H-PSMG\textsuperscript{+}'s burn-in benefit over vanilla H-PSMG grows monotonically with type-space size $|\Theta_i| \in \{2, \ldots, 6\}$ and with agent count $n \in \{3, 4, 5\}$, the latter together with the analytical $O(|\Theta_i|^n / n)$ storage advantage. On the independently-designed DeepMind Concordia \citep{vezhnevets2023concordia} substrates from the NeurIPS 2024 Cooperative AI Contest \citep{concordia_contest2025}, H-PSMG\textsuperscript{+} achieves within $1\%$ of oracle on all $7$ Pub Coordination configurations and wins or ties oracle on all $9$ Haggling configurations. On SOTOPIA-Hard \citep{zhou2024sotopia}, a benchmark where HT-MG-BC's prior-factorisation assumption breaks, a substrate-adapted variant of H-PSMG\textsuperscript{+} still attains the highest evaluator score on every backbone, demonstrating that the theory transfers even where its formal preconditions weaken.
\end{abstract}

\section{Introduction}
\label{sec:intro}

Multi-agent reinforcement learning in environments where agents differ in their preferences, capabilities, or roles arises naturally in cooperative AI, mechanism design, and economic agent simulation \citep{zhang2021multi,leibo2017multi,perolat2017multi}. The standard Markov game (MG) model \citep{shapley1953stochastic,littman1994markov} captures multi-agent sequential decision making but treats agents as exchangeable in their type structure. In practice agents possess latent attributes, including reward functions, that are private and unknown to others. Reasoning under this kind of uncertainty is the defining feature of Bayesian games \citep{harsanyi1968games} but has remained under-developed in the sequential MG setting with explicit regret guarantees.

A parallel line of work on posterior sampling for reinforcement learning (PSRL) \citep{strens2000bayesian,osband2013more,agrawal2017optimistic,osband2017posterior} has established that Bayesian methods can achieve near-optimal regret in single-agent MDPs. Extensions to two-player zero-sum stochastic games have been carried out by \citet{jafarnia2021learning} and \citet{xiong2022self}, with the most recent contribution by \citet{yueh2026bayesian} addressing episodic finite-horizon zero-sum games. For cooperative MARL, \citet{hsu2024randomized} provided the first randomized exploration analysis in parallel MDPs. Yet a fundamental gap remains. None of these works addresses the case where agents are heterogeneous in reward type and the learner must simultaneously infer (i) the transition dynamics, and (ii) the latent type profile of all participants. This combination is the natural model for cooperative AI with hidden agent preferences.

This paper develops the foundational result for this regime. We introduce the HT-MG-BC framework, formalize the structural assumptions under which posterior reasoning remains tractable, and prove the posterior decoupling theorem. Our contributions are as follows.

\textbf{(1) Framework.} We formalize the HT-MG-BC setting, which couples a Markov game with a Bayesian prior over both the transition kernel and the joint type profile. We define the type-conditional Bayesian regret as the appropriate performance criterion.

\textbf{(2) Algorithm.} We propose two posterior-sampling algorithms: H-PSMG (Algorithm~\ref{alg:hpsmg}), which samples world model and type profile independently and executes the CCE of the sampled world; and H-PSMG\textsuperscript{+} (Algorithm~\ref{alg:hpsmg-plus}), a type-discriminative refinement of H-PSMG enabled by access to per-agent factored marginals.

\textbf{(3) Decoupling theorem.} We prove that under type-independent transitions and own-type reward locality, the joint posterior over the transition kernel and the type profile exactly factorizes at every episode (Theorem~\ref{thm:decoupling}). The factorization holds further across agents, yielding storage and computational savings exponential in the number of agents.

\textbf{(4) Numerical illustration and tightness.} We work through a complete two-agent two-type matrix game example, verify the decoupling numerically, and prove tightness by exhibiting two explicit counter-examples that show both structural assumptions are necessary.

\textbf{(5) Regret upper bound for H-PSMG.} We prove a Bayesian regret bound of $\tilde O(H^{3/2}\sqrt{SA_{\mathrm{joint}}K} + H\sqrt{|\Theta_i| K})$ for H-PSMG (Theorem~\ref{thm:main-regret-informal}), via an additive decomposition into a Dirichlet--Bernstein transition component and a Russo--Van Roy information-ratio type component \citep{russo2014learning,russo2016information}.

\textbf{(6) Rate separation from heuristics.} We construct an exploration-trap HT-MG on which H-PSMG retains its $\tilde O(\sqrt{K})$ rate while MAP-type-greedy and type-agnostic posterior sampling both incur $\Omega(K)$ linear regret (Theorem~\ref{thm:rate-separation}). The trap exposes the necessity of Bayesian exploration over types.

\textbf{(7) Exponential-in-$|\Theta|^n$ burn-in improvement of H-PSMG\textsuperscript{+}.} The type-discrimination bonus is structurally unavailable to Joint-PSRL without first invoking the decoupling theorem. We prove (Theorem~\ref{thm:hpsmg-plus}) that H-PSMG\textsuperscript{+} matches the asymptotic regret bound of H-PSMG and, on a parametric family $\mathcal{G}_m^{\mathrm{trap}}$ of deep-trap HT-MGs, achieves $O(H)$ expected burn-in regret while vanilla H-PSMG (and hence Joint-PSRL under matched RNG) incurs $\Omega(|\Theta_i|^n H)$ burn-in. This is an exponential-in-$n$ separation in burn-in cost. We give the full formal proof in Appendix~\ref{app:hpsmg-plus} and verify the prediction empirically on an LLM substrate: the separation grows monotonically with type-space size $|\Theta_i| \in \{2, \ldots, 6\}$ and with agent count $n \in \{3, 4, 5\}$ (Section~\ref{sec:experiments-llm}).

\textbf{(8) Minimax lower bound.} For any algorithm there exists an HT-MG instance with Bayesian regret $\Omega(\sqrt{H^3 S A_{\mathrm{joint}} K} + \sqrt{|\Theta_i| HK})$ (Theorem~\ref{thm:lower-matching}), matching our upper bound modulo a $\sqrt{H}$ gap on the type-inference term and logarithmic factors.

\textbf{(9) Exponential improvement over generic Bayesian RL.} Without prior factorization (Assumption~\ref{ass:PF}), any algorithm has regret $\Omega(\sqrt{|\Theta|^n K})$ in the worst case (Theorem~\ref{thm:lower-no-pf}). Our analysis therefore provides an exponential-in-$n$ improvement: the regret bound for any Bayesian-RL algorithm on HT-MG, including Joint-PSRL, drops from $\tilde O(\sqrt{|\Theta|^n K})$ to $\tilde O(\sqrt{|\Theta_i| K})$ through the decoupling theorem.

\textbf{(10) Empirical validation on LLM substrates.} We evaluate on three benchmarks. On the Heterogeneous-Persona Sequential Public Goods Game (HP-SPGG), a controlled LLM-MAS benchmark we construct, the H-PSMG family achieves $6$ to $12\times$ lower cumulative regret than published LLM-coordination baselines (ECON, A-ToM family) across four backbones, and a $\beta$-sweep exhibits three regimes (extreme self-decay, self-decay edge, burn-in jump) of Theorem~\ref{thm:hpsmg-plus} consistent with each backbone's posterior concentration speed. On DeepMind's independently-designed Concordia substrates from the NeurIPS 2024 Cooperative AI Contest \citep{vezhnevets2023concordia, concordia_contest2025}, H-PSMG\textsuperscript{+} achieves the highest non-oracle focal score on $16/16$ configurations ($7$ Pub Coordination, $9$ Haggling). On SOTOPIA-Hard \citep{zhou2024sotopia}, a dyadic social-negotiation benchmark whose action space is open-ended natural language and whose persona structure violates HT-MG-BC's prior-factorisation assumption, the substrate-adapted H-PSMG\textsuperscript{+} variant achieves the highest evaluator score across all four backbones, with the largest margin on Kimi-K2.6 ($+0.36$). Theory-driven experiments on HP-SPGG further measure backbone posterior concentration speed and verify the predicted scaling of the bonus benefit with type-space size and agent count.

The rest of the paper is organized as follows. Section~\ref{sec:related} surveys related work. Section~\ref{sec:setting} formalizes HT-MG-BC and the structural assumptions. Section~\ref{sec:algorithm} presents H-PSMG and H-PSMG\textsuperscript{+}. Section~\ref{sec:decoupling} states and proves the posterior decoupling theorem. Section~\ref{sec:example} works through a two-agent numerical illustration. Section~\ref{sec:necessity} establishes necessity of the assumptions. Section~\ref{sec:regret-roadmap} derives the regret bounds: §\ref{sec:rate-separation}--\ref{sec:hpsmg-plus} cover rate separation and the H-PSMG\textsuperscript{+} burn-in result; §\ref{sec:exponential-improvement}--\ref{sec:lower-bound-matching} cover lower bounds. Section~\ref{sec:experiments-llm} reports the empirical study across HP-SPGG, Concordia, and SOTOPIA-Hard, including theory-driven experiments measuring posterior concentration speed and the predicted scaling of the bonus benefit. Section~\ref{sec:discussion} discusses extensions including function approximation.

\section{Related Work}
\label{sec:related}

\paragraph{Posterior sampling in single-agent RL.}
The PSRL framework was first proposed by \citet{strens2000bayesian} and received a Bayesian regret analysis by \citet{osband2013more}, who obtained $\tilde{O}(\tau S\sqrt{AT})$ for finite-horizon MDPs. \citet{agrawal2017optimistic} achieved frequentist worst-case regret $\tilde{O}(D\sqrt{SAT})$ in communicating MDPs through anti-concentration of the Dirichlet distribution, matching the lower bound of \citet{jaksch2010near} in $S$. \citet{ouyang2017learning} addressed non-episodic settings via dynamic episode lengths, and \citet{russo2014learning} developed the information-theoretic analysis framework that we will rely on for the type-inference component.

\paragraph{Posterior sampling in Markov games.}
\citet{jafarnia2021learning} introduced PSRL-ZSG for infinite-horizon average-reward zero-sum stochastic games, obtaining Bayesian regret $\tilde{O}(HS\sqrt{AT})$ against arbitrary history-dependent opponents. \citet{xiong2022self} extended posterior sampling to zero-sum MGs with general function approximation, introducing the multi-agent decoupling coefficient. The most recent result, due to \citet{yueh2026bayesian}, analyzes episodic finite-horizon zero-sum MGs with two-sided and one-sided posterior sampling.

All three results assume zero-sum payoffs and known type structure. Our work moves to general-sum games with unknown type structure, which is qualitatively different.

\paragraph{Multi-agent randomized exploration.}
\citet{hsu2024randomized} presented the first randomized exploration framework for cooperative MARL through CoopTS-PHE and CoopTS-LMC. Their setting is parallel MDPs, where agents share the same MDP and the goal is collaborative exploration of a common environment. Our setting is fundamentally different in that agents have private reward types and the game is not cooperative.

\paragraph{Latent-type and heterogeneous learning.}
\citet{kwon2021rl} established sample complexity and a lower bound for learning latent MDPs, showing that without identifiability assumptions the worst-case regret is exponential in the number of types. \citet{tian2024learning} extended hardness barriers to MGs with adaptive adversaries. On the other hand, \citet{tuynman2022transfer} considered transfer with heterogeneous reward functions in a shared MDP. None of these works combines posterior sampling with type heterogeneity in a game.

\paragraph{Structured priors and Bayesian RL.}
\citet{mutti2024exploiting} introduced Causal-PSRL, a hierarchical Bayesian procedure exploiting causal graph priors in factored MDPs. Their hierarchical structure is conceptually similar to ours in that the prior decomposes into a model component and a structural component, but they treat the structural prior as factoring the dynamics rather than partitioning agents. \citet{arumugam2026toward} recently revived PSRL as an algorithmic primitive for LLM-based agent design.

\paragraph{General-sum equilibrium learning.}
A line of work on general-sum MGs has established sample complexity for learning coarse correlated equilibria via V-learning and its variants \citep{jin2021v,song2022when,mao2023decentralized}. \citet{erez2023regret} provided regret guarantees in this setting. Most recently, \citet{cui2025incentivize} proposed VMG, a model-based optimistic algorithm for both zero-sum NE and general-sum CCE. These works do not employ posterior sampling and do not handle type uncertainty.

\paragraph{Policy-Space Response Oracles and population-based MARL.}
A parallel and highly influential family of multi-agent RL methods follows the Policy-Space Response Oracle (PSRO) framework of \citet{lanctot2017unified}, which generalises the Double Oracle algorithm \citep{mcmahan2003planning} by maintaining a growing policy population, training a best-response oracle (typically a deep RL learner) against a meta-strategy mixture each iteration, and updating the meta-strategy by solving an empirical game over the current population. Variants improve scalability, convergence, and generality. Pipeline PSRO (P2SRO) \citep{mcaleer2020pipeline} parallelises the oracle training across a hierarchical pipeline and is guaranteed to converge to an approximate Nash equilibrium in two-player zero-sum games. $\alpha$-PSRO \citep{muller2019generalized} replaces the Nash meta-solver with $\alpha$-Rank, extending PSRO to general-sum, many-player settings with uniqueness guarantees. JPSRO \citep{marris2021multi} works with mixed joint policies and a Maximum-Gini (Coarse) Correlated Equilibrium meta-solver, achieving convergence to a normal-form (C)CE in $n$-player general-sum extensive-form games. XDO \citep{mcaleer2021xdo} extends PSRO to extensive-form games at the information-set level. Anytime PSRO \citep{mcaleer2022anytime} guarantees monotonically decreasing exploitability across iterations. Online and regret-minimising variants such as Online Double Oracle \citep{dinh2022online} and Efficient PSRO \citep{zhou2022efficient} integrate no-regret online learning at the meta-strategy level. Diversity-promoting variants such as PSD-PSRO \citep{yao2023policy} and DPP-PSRO regularise the population to span the policy space more effectively. Most recent extensions include sample-efficient joint-experience best response \citep{smith2026sample} and LLM-based code-space oracles \citep{hennes2026code}.

PSRO and its descendants pursue a distinct goal from this paper. PSRO targets convergence of the meta-strategy to an equilibrium (Nash, $\alpha$-Rank, CCE) of an empirical game whose payoffs are computed by repeated simulation of pairs of population members. The latent quantity that PSRO infers iteration-by-iteration is each agent's \emph{best-response policy} against the current opponent mixture; the game's payoff structure (rewards, transitions, types) is treated as a black-box simulator. By contrast, our HT-MG-BC framework places the uncertainty on the \emph{game's parameters} -- transition kernel $P$ and per-agent type $\theta_i$ -- and maintains an explicit Bayesian posterior over them. We do not iteratively grow a policy population; we sample one world from the posterior per episode and execute the CCE of the sampled world. Consequently the performance criteria differ: PSRO is evaluated by NashConv / exploitability of the empirical equilibrium (an asymptotic, population-level quantity), whereas our $\tilde O(\sqrt K)$ Bayesian regret bound (Theorem~\ref{thm:main-regret-informal}) is a finite-time, online guarantee against the type-aware oracle policy. JPSRO is the closest variant in that it also uses CCE as its solution concept for $n$-player general-sum coordination, but it assumes the payoff structure is known and addresses the dual problem of finding policies that realise the equilibrium, rather than inferring the latent payoff structure from sequential observations. A possible integration -- using PSRO-style population learning as the inner CCE oracle inside H-PSMG -- is consistent with our oracle abstraction (Section~\ref{sec:algorithm}) and is a natural direction for empirical work at scale; the decoupling theorem (Theorem~\ref{thm:decoupling}) and regret bounds would still apply provided the inner oracle returns an $\epsilon_{\mathrm{CCE}}$-accurate solution.

\paragraph{Ad-hoc coordination with type-based opponent modelling.}
Closer in spirit to our setting is the line on ad-hoc coordination with Bayesian opponent type inference, in particular Harsanyi-Bellman Ad-hoc Coordination (HBA) \citep{albrecht2013game,albrecht2019convergence}, which maintains a Bayesian posterior over a finite set of pre-specified opponent policy types and plans against the posterior via a Bayesian Nash--Bellman backup. Conceptually our $\mu^{\Theta, i}$ generalises HBA's posterior to multiple agents and combines it with a posterior over transitions, while HBA assumes the environment dynamics are known and the type set is provided by a domain expert. HBA has empirical evaluations but no finite-time regret guarantee, and the convergence analysis of \citet{albrecht2019convergence} addresses asymptotic best-response correctness rather than online regret. The decoupling theorem and $\tilde O(\sqrt K)$ regret bound of this paper are therefore the natural finite-sample complement to HBA's asymptotic analysis.

\section{Problem Setting}
\label{sec:setting}

\subsection{Heterogeneous-Type Markov Game}

We consider a heterogeneous-type Markov game (HT-MG) specified by the tuple
\begin{equation}
\mathcal{G} = \langle \mathcal{S}, \{\mathcal{A}_i\}_{i \in [n]}, P^*, \{\Theta_i\}_{i \in [n]}, \{r_i\}_{i \in [n]}, P_0, H \rangle
\end{equation}
where $\mathcal{S}$ is a finite shared state space of size $S$, $\mathcal{A}_i$ is the finite action space of agent $i$, $\mathcal{A} := \prod_{i \in [n]} \mathcal{A}_i$ is the joint action space of size $A := \prod_i |\mathcal{A}_i|$, $P^* : \mathcal{S} \times \mathcal{A} \times \Theta \to \Delta(\mathcal{S})$ is the joint transition kernel with $\Theta := \prod_i \Theta_i$, $r_i : \mathcal{S} \times \mathcal{A} \times \Theta_i \to [0,1]$ is agent $i$'s type-dependent reward function, $P_0 \in \Delta(\Theta) \times \Delta(\mathrm{Kernels}(\mathcal{S} \times \mathcal{A}))$ is the joint prior over the type profile and transition kernel, and $H$ is the episode length.

At the start of each episode, a type profile $\boldsymbol{\theta}^* = (\theta_1^*, \ldots, \theta_n^*) \in \Theta$ is sampled once from the marginal type prior and held fixed across the episode. The transition kernel $P^*$ is fixed across all episodes but unknown. Agent $i$ knows its own type $\theta_i^*$ but not the types of other agents. In the centralized formulation analyzed in this paper, a single learner manages decisions for all agents and maintains a single posterior over $(P^*, \boldsymbol{\theta}^*)$.

For an episode $k$, the centralized learner observes
\begin{equation}
\tau_k = (s_1^k, a_1^k, r_1^k, s_2^k, a_2^k, r_2^k, \ldots, s_H^k, a_H^k, r_H^k, s_{H+1}^k)
\end{equation}
where $a_h^k = (a_h^{k,1}, \ldots, a_h^{k,n}) \in \mathcal{A}$ is the joint action at step $h$ and $r_h^k = (r_h^{k,1}, \ldots, r_h^{k,n})$ is the joint reward vector. The history at the start of episode $k$ is $\mathcal{H}_{k-1} = (\tau_1, \ldots, \tau_{k-1})$.

\subsection{Structural Assumptions}

The decoupling theorem rests on two structural assumptions.

\begin{assumption}[Type-Independent Transitions, TI]
\label{ass:TI}
The transition kernel does not depend on the type profile: $P^*(s' \mid s, a, \boldsymbol{\theta}) = P^*(s' \mid s, a)$ for all $s, s' \in \mathcal{S}, a \in \mathcal{A}, \boldsymbol{\theta} \in \Theta$.
\end{assumption}

\begin{assumption}[Reward Locality, RL]
\label{ass:RL}
Each agent's reward depends only on its own type: $r_i(s, a, \boldsymbol{\theta}) = r_i(s, a, \theta_i)$ for all $i \in [n], s, a, \boldsymbol{\theta}$. Equivalently, the reward likelihood satisfies $P(r_i \mid s, a, \boldsymbol{\theta}) = P(r_i \mid s, a, \theta_i)$.
\end{assumption}

\begin{remark}
Assumption~\ref{ass:RL} permits $r_i$ to depend on the actions of other agents $a_{-i}$. Such dependence is natural in cooperative and social-dilemma settings, where own reward responds to other agents' cooperation. The locality is on the latent type only, not on the action profile.
\end{remark}

\begin{assumption}[Prior Factorization, PF]
\label{ass:PF}
The joint prior factorizes as $P_0(P, \boldsymbol{\theta}) = P_0^P(P) \prod_{i \in [n]} P_0^{\Theta,i}(\theta_i)$.
\end{assumption}

These three assumptions are mild in cooperative AI settings where agents share a physical environment but possess private utility functions. They fail in adversarial settings where opponents may sabotage the dynamics or in social games where one agent's utility directly depends on another's hidden type.

The next three assumptions are used in the regret analysis (Section~\ref{sec:regret-roadmap} and Appendix~\ref{app:regret-decomp}) but not in the decoupling theorem itself.

\begin{assumption}[Type Identifiability, TID]
\label{ass:TID}
There exists $\rho > 0$ such that for every $\theta_i \neq \theta_i' \in \Theta_i$ and every $i \in [n]$, the per-step squared Hellinger distance between reward distributions is uniformly lower-bounded:
\begin{equation}
\inf_{(s,a) \in \mathcal{S} \times \mathcal{A}} D_{\mathrm{H}}^2\big(P(\cdot \mid s, a, \theta_i),\, P(\cdot \mid s, a, \theta_i')\big) \geq \rho.
\end{equation}
\end{assumption}

TID guarantees that any policy that visits a positive-probability $(s, a)$ pair produces a reward signal distinguishing the candidate types. It is the minimal condition under which categorical type posteriors concentrate at a polynomial rate. Without TID, the type-inference regret may scale exponentially in $n$ \citep{kwon2021rl}.

\begin{assumption}[Initial-State Independence, ISI]
\label{ass:ISI}
The initial state distribution $\beta \in \Delta(\mathcal{S})$ from which $s_1^k$ is drawn at the start of each episode is fixed across episodes and independent of the latent $(P^*, \boldsymbol{\theta}^*)$.
\end{assumption}

ISI is standard in episodic RL and is needed for the posterior-sampling equivalence (Lemma~\ref{lem:ps-equivalence} in Appendix~\ref{app:regret-decomp}) to hold conditionally on the initial state.

\begin{assumption}[CCE Selection Rule, CSR]
\label{ass:CSR}
The CCE oracle $\mathrm{CCE}: (P, \boldsymbol{\theta}) \mapsto \pi$ is a measurable, deterministic mapping. When multiple CCE policies exist, ties are broken according to a fixed rule (for instance, lexicographic ordering on policy parameters, or the utilitarian-maximizing CCE).
\end{assumption}

CSR is necessary in general-sum MGs because CCE values are non-unique. Treating the oracle as a deterministic mapping ensures the posterior-sampling argument $\mathbb{E}[f(\hat{P}_k, \hat{\boldsymbol{\theta}}_k)] = \mathbb{E}[f(P^*, \boldsymbol{\theta}^*)]$ applies to functions like $f(P, \boldsymbol{\theta}) = V_i^{\mathrm{CCE}(P, \boldsymbol{\theta})}(s; P, \theta_i)$.

\subsection{Performance Criterion}

For each type profile $\boldsymbol{\theta} \in \Theta$ and transition kernel $P$, let $\pi^{\mathrm{CCE}}(P, \boldsymbol{\theta}) \in \Delta(\mathcal{A})^{\mathcal{S}}$ denote a coarse correlated equilibrium policy of the induced finite-horizon MG. We define the type-conditional Bayesian regret over $K$ episodes as
\begin{equation}
\label{eq:regret-def}
\mathrm{Reg}_i^B(K) := \mathbb{E}_{(P^*, \boldsymbol{\theta}^*) \sim P_0} \left[ \sum_{k=1}^K \left( V_i^{\pi^{\mathrm{CCE}}(P^*, \boldsymbol{\theta}^*)}(s_1^k; \theta_i^*) - V_i^{\pi_k}(s_1^k; \theta_i^*) \right) \right]
\end{equation}
where $V_i^\pi(s; \theta_i)$ is agent $i$'s expected $H$-step return when policy $\pi$ is executed from state $s$ under reward $r_i(\cdot, \cdot, \theta_i)$, and $\pi_k$ is the policy executed by the algorithm in episode $k$. Expectation is taken over the prior, transition randomness, and algorithmic randomness.

\section{Algorithm: Hierarchical Posterior Sampling for Markov Games}
\label{sec:algorithm}

We propose two posterior-sampling algorithms for HT-MG-BC: H-PSMG, the basic factored posterior sampler (Section~\ref{sec:hpsmg-vanilla}), and H-PSMG\textsuperscript{+}, a type-discriminative refinement that adds an exploration bonus computed from the factored marginals (Section~\ref{sec:hpsmg-plus-algo}). Both reduce posterior storage from exponential to linear in the number of agents (Section~\ref{sec:decoupling}); H-PSMG\textsuperscript{+} additionally enjoys an exponential-in-$|\Theta|^n$ burn-in regret improvement over H-PSMG (Section~\ref{sec:hpsmg-plus}).

\subsection{H-PSMG: posterior sampling with factored marginals}
\label{sec:hpsmg-vanilla}

H-PSMG maintains $n+1$ independent posterior modules -- one transition posterior $\mu^P$ and $n$ per-agent type posteriors $\mu^{\Theta, i}$ -- and at each episode samples a world from each posterior, computes the CCE of the sampled world, and executes it for one rollout.

\begin{algorithm}[ht]
\caption{Hierarchical Posterior Sampling for Markov Games (H-PSMG)}
\label{alg:hpsmg}
\textbf{Input}: Priors $P_0^P, \{P_0^{\Theta,i}\}_{i \in [n]}$; horizon $H$; episodes $K$; CCE oracle $\mathrm{CCE}(\cdot, \cdot)$.\\
\textbf{Output}: Policies $\{\pi_k\}_{k=1}^K$.
\begin{algorithmic}[1]
\STATE Initialize $\mu_1^P \gets P_0^P$ and $\mu_1^{\Theta,i} \gets P_0^{\Theta,i}$ for all $i \in [n]$.
\FOR{$k = 1, 2, \ldots, K$}
    \STATE Sample world: $\hat{P}_k \sim \mu_k^P$ and $\hat{\theta}_i^k \sim \mu_k^{\Theta,i}$ independently for $i \in [n]$. Set $\hat{\boldsymbol{\theta}}_k = (\hat{\theta}_1^k, \ldots, \hat{\theta}_n^k)$.
    \STATE Compute equilibrium: $\pi_k \gets \mathrm{CCE}(\hat{P}_k, \hat{\boldsymbol{\theta}}_k)$.
    \STATE Execute $\pi_k$ for $H$ steps from $s_1^k$. Observe $\tau_k$.
    \STATE Update marginals via Bayes:
    $\mu_{k+1}^P(P) \propto \mu_k^P(P) \cdot L_P(\tau_k; P)$,
    $\mu_{k+1}^{\Theta,i}(\theta_i) \propto \mu_k^{\Theta,i}(\theta_i) \cdot L_{\Theta,i}(\tau_k; \theta_i)$,
    where $L_P, L_{\Theta,i}$ are defined in Section~\ref{sec:decoupling}.
\ENDFOR
\end{algorithmic}
\end{algorithm}

The algorithm stores $n$ separate type posteriors of size $|\Theta_i|$ each instead of a joint posterior over $|\Theta|^n = \prod_i |\Theta_i|$ configurations. Theorem~\ref{thm:decoupling} below justifies this factorisation as exact, not an approximation. The CCE oracle is assumed to return an $\epsilon_{\mathrm{CCE}}$-approximate CCE in $\mathrm{poly}(|\mathcal{S}|, A, H, 1/\epsilon_{\mathrm{CCE}})$ time; constructions are available via \citet{jin2021v} and \citet{mao2023decentralized}.

\subsection{H-PSMG\textsuperscript{+}: type-discriminative refinement}
\label{sec:hpsmg-plus-algo}

A failure mode of vanilla posterior sampling on HT-MGs arises when a sampled type profile $\hat{\boldsymbol{\theta}}$ yields uniform predicted reward $r_i(s, a, \hat\theta_i)$ across $(s, a)$, so the sampled-type-conditioned $Q^\star(\cdot; \hat{\boldsymbol{\theta}})$ admits ties at the initial state. The CCE oracle then resolves the tie arbitrarily, and an unlucky tie pattern leaves the rollout stuck in low-information states. The exploration-trap game of Section~\ref{sec:rate-separation} is the canonical instance.

H-PSMG\textsuperscript{+} (Algorithm~\ref{alg:hpsmg-plus}) augments H-PSMG with a per-state-action \emph{type-discrimination bonus}
\begin{equation}
\label{eq:bonus}
D_k(s, a) \;\coloneqq\; \sum_{i=1}^n \mathbb{E}_{\theta, \theta' \stackrel{\mathrm{iid}}{\sim} \mu_k^{\Theta, i}}\!\Big[\,\big|r_i(s, a, \theta) - r_i(s, a, \theta')\big|\,\Big],
\end{equation}
computed from the current factored marginals. During the inner value iteration, the bonus is added to the $Q$-function with weight $\beta > 0$: $Q^+_h(s, a) \mathrel{+}= \beta\, D_k(s, a)$. Backward induction propagates the bonus, so the CCE policy is biased toward trajectories that visit state-actions where competing types disagree in reward prediction.

\begin{algorithm}[ht]
\caption{H-PSMG\textsuperscript{+}: H-PSMG with type-discrimination bonus}
\label{alg:hpsmg-plus}
\textbf{Input}: Priors $P_0^P, \{P_0^{\Theta,i}\}$; horizon $H$; episodes $K$; CCE oracle; bonus weight $\beta > 0$.
\begin{algorithmic}[1]
\STATE Initialise $\mu_1^P \gets P_0^P$, $\mu_1^{\Theta, i} \gets P_0^{\Theta, i}$ for all $i$.
\FOR{$k = 1, 2, \ldots, K$}
  \STATE Sample $\hat P_k \sim \mu_k^P$ and $\hat\theta_{i,k} \sim \mu_k^{\Theta, i}$ for each $i$.
  \STATE Compute $D_k(s, a)$ from current marginals via Equation~\eqref{eq:bonus}.
  \STATE \textbf{Bonus-augmented planning}: solve CCE with $Q^+_h(s, a) \mathrel{+}= \beta\, D_k(s, a)$ in backward induction.
  \STATE Execute the CCE policy, observe $\tau_k = (s_h, a_h, r_h^i)_{h=1}^H$.
  \STATE Update $\mu_{k+1}^P$ and $\mu_{k+1}^{\Theta, i}$ as in Algorithm~\ref{alg:hpsmg}.
\ENDFOR
\end{algorithmic}
\end{algorithm}

The bonus has three properties that make it both elegant and theoretically tractable.

\textbf{(i) Decoupling-driven.} $D_k$ depends explicitly on per-agent marginals $\mu_k^{\Theta, i}$, which H-PSMG stores natively. Joint-PSRL maintains only the joint distribution $\mu_k^\Theta$ on $\Theta = \prod_i \Theta_i$; to compute $D_k$ it must marginalise the joint at cost $O(|\Theta|^n)$ per query, and by Theorem~\ref{thm:decoupling} the resulting marginals are exactly those stored by H-PSMG. The bonus is therefore an algorithmic primitive natural to the factored representation, not to the joint one.

\textbf{(ii) Self-decaying.} As $\mu_k^{\Theta, i}$ concentrates on the true type, $\mu_k^{\Theta, i}(\theta)\, \mu_k^{\Theta, i}(\theta') \to 0$ for $\theta \neq \theta'$, hence $D_k(s, a) \to 0$ pathwise along any trajectory. The bonus does not contaminate the asymptotic regret rate; it only operates during the burn-in phase.

\textbf{(iii) Tie-resolving.} On instances with value-function ties at sampled types (e.g., the exploration-trap family), $D_k > 0$ at type-discriminative states and propagates backward through value iteration to break the tie, eliminating the trap. The exponential-in-$|\Theta|^n$ burn-in improvement this enables is formalised in Section~\ref{sec:hpsmg-plus}.

The two algorithms produce \emph{bit-identical} trajectories under matched RNG when $\beta = 0$; H-PSMG\textsuperscript{+} departs from H-PSMG only via the bonus term. Joint-PSRL is bit-identical to H-PSMG (Theorem~\ref{thm:decoupling}) but cannot replicate H-PSMG\textsuperscript{+} without first invoking the decoupling theorem to extract per-agent marginals from the joint posterior. The decoupling theorem is thus consumed both analytically (Section~\ref{sec:exponential-improvement}) and algorithmically (this section).

\section{Posterior Decoupling Theorem}
\label{sec:decoupling}

This section states and proves the main structural result.

\subsection{Likelihood Factorization}

We first derive the conditional likelihood of an observed trajectory given the latent world. Under H-PSMG, the policy $\pi_k$ at episode $k$ is a function only of the history $\mathcal{H}_{k-1}$ and the algorithmic posterior samples $(\hat{P}_k, \hat{\boldsymbol{\theta}}_k)$, which are themselves measurable with respect to $\mathcal{H}_{k-1}$. Crucially, $\pi_k$ does not depend on the true latent $(P^*, \boldsymbol{\theta}^*)$.

\begin{lemma}[Joint Likelihood Factorization]
\label{lem:likelihood}
Under Assumptions~\ref{ass:TI} and \ref{ass:RL}, the trajectory likelihood at episode $k$ factorizes as
\begin{equation}
\label{eq:lik-factor}
P(\tau_k \mid P, \boldsymbol{\theta}, \mathcal{H}_{k-1}) = c_k(\tau_k) \cdot L_P(\tau_k; P) \cdot \prod_{i \in [n]} L_{\Theta,i}(\tau_k; \theta_i)
\end{equation}
where
\begin{align}
c_k(\tau_k) &:= \mathbb{E}_{\pi_k \mid \mathcal{H}_{k-1}}\left[\prod_{h=1}^H \pi_k(a_h^k \mid s_h^k)\right], \label{eq:ck} \\
L_P(\tau_k; P) &:= \prod_{h=1}^H P(s_{h+1}^k \mid s_h^k, a_h^k), \label{eq:lp} \\
L_{\Theta,i}(\tau_k; \theta_i) &:= \prod_{h=1}^H P(r_h^{k,i} \mid s_h^k, a_h^k, \theta_i). \label{eq:ltheta}
\end{align}
The constant $c_k(\tau_k)$ depends on $\tau_k$ and $\mathcal{H}_{k-1}$ but not on $(P, \boldsymbol{\theta})$.
\end{lemma}

\begin{proof}
Conditional on $\pi_k$, the trajectory likelihood unrolls along the time axis as
\begin{equation}
P(\tau_k \mid \pi_k, P, \boldsymbol{\theta}) = \prod_{h=1}^H \pi_k(a_h^k \mid s_h^k) \cdot P(s_{h+1}^k \mid s_h^k, a_h^k, \boldsymbol{\theta}) \cdot \prod_{i \in [n]} P(r_h^{k,i} \mid s_h^k, a_h^k, \boldsymbol{\theta}).
\end{equation}
Apply Assumption~\ref{ass:TI} to drop $\boldsymbol{\theta}$ from the transition factor, and Assumption~\ref{ass:RL} to localize the reward factor to $\theta_i$:
\begin{equation}
P(\tau_k \mid \pi_k, P, \boldsymbol{\theta}) = \prod_{h=1}^H \pi_k(a_h^k \mid s_h^k) \cdot \prod_{h=1}^H P(s_{h+1}^k \mid s_h^k, a_h^k) \cdot \prod_{h=1}^H \prod_{i \in [n]} P(r_h^{k,i} \mid s_h^k, a_h^k, \theta_i).
\end{equation}
Marginalize over $\pi_k$ with respect to $P(\pi_k \mid \mathcal{H}_{k-1})$, which is independent of $(P, \boldsymbol{\theta})$:
\begin{align}
P(\tau_k \mid P, \boldsymbol{\theta}, \mathcal{H}_{k-1}) &= \int P(\pi_k \mid \mathcal{H}_{k-1}) \cdot P(\tau_k \mid \pi_k, P, \boldsymbol{\theta}) \, d\pi_k \\
&= \mathbb{E}_{\pi_k \mid \mathcal{H}_{k-1}}\left[\prod_{h=1}^H \pi_k(a_h^k \mid s_h^k)\right] \cdot L_P(\tau_k; P) \cdot \prod_{i \in [n]} L_{\Theta,i}(\tau_k; \theta_i).
\end{align}
The factors $L_P$ and $L_{\Theta,i}$ depend respectively on $P$ alone and $\theta_i$ alone, and the leading expectation depends on $\tau_k$ and $\mathcal{H}_{k-1}$ alone.
\end{proof}

\subsection{Main Theorem}

\begin{theorem}[Posterior Decoupling]
\label{thm:decoupling}
Under Assumptions~\ref{ass:TI}, \ref{ass:RL}, and \ref{ass:PF}, the posterior distribution over the latent world factorizes at every episode:
\begin{equation}
\label{eq:decoupling}
P(P, \boldsymbol{\theta} \mid \mathcal{H}_k) = \mu_{k+1}^P(P) \cdot \prod_{i \in [n]} \mu_{k+1}^{\Theta,i}(\theta_i), \quad \forall k \in \{0, 1, \ldots, K\}
\end{equation}
where the marginals admit the recursive updates
\begin{equation}
\mu_{k+1}^P(P) \propto \mu_k^P(P) \cdot L_P(\tau_k; P), \quad \mu_{k+1}^{\Theta,i}(\theta_i) \propto \mu_k^{\Theta,i}(\theta_i) \cdot L_{\Theta,i}(\tau_k; \theta_i)
\end{equation}
with initial conditions $\mu_1^P = P_0^P$ and $\mu_1^{\Theta,i} = P_0^{\Theta,i}$.
\end{theorem}

\begin{proof}
We proceed by induction on $k$.

\textit{Base case ($k=0$).} By Assumption~\ref{ass:PF}, $P(P, \boldsymbol{\theta} \mid \mathcal{H}_0) = P_0(P, \boldsymbol{\theta}) = P_0^P(P) \prod_{i \in [n]} P_0^{\Theta,i}(\theta_i)$, which has the required factored form.

\textit{Inductive step.} Suppose $P(P, \boldsymbol{\theta} \mid \mathcal{H}_{k-1}) = \mu_k^P(P) \prod_{i} \mu_k^{\Theta,i}(\theta_i)$. By Bayes' rule,
\begin{equation}
P(P, \boldsymbol{\theta} \mid \mathcal{H}_k) \propto P(P, \boldsymbol{\theta} \mid \mathcal{H}_{k-1}) \cdot P(\tau_k \mid P, \boldsymbol{\theta}, \mathcal{H}_{k-1}).
\end{equation}
Apply Lemma~\ref{lem:likelihood}:
\begin{align}
P(P, \boldsymbol{\theta} \mid \mathcal{H}_k) &\propto \mu_k^P(P) \prod_{i} \mu_k^{\Theta,i}(\theta_i) \cdot c_k(\tau_k) \cdot L_P(\tau_k; P) \cdot \prod_{i} L_{\Theta,i}(\tau_k; \theta_i) \\
&\propto \left[\mu_k^P(P) L_P(\tau_k; P)\right] \cdot \prod_{i} \left[\mu_k^{\Theta,i}(\theta_i) L_{\Theta,i}(\tau_k; \theta_i)\right].
\end{align}
The factor $c_k(\tau_k)$ does not depend on $(P, \boldsymbol{\theta})$ and is absorbed into the normalization constant. The remaining product factorizes into a $P$-only term and $n$ separate $\theta_i$-only terms. Normalizing each factor independently yields the recursive forms stated in the theorem and establishes the desired posterior factorization.
\end{proof}

\begin{corollary}[Storage and Update Complexity]
\label{cor:storage}
Maintaining the joint posterior $P(P, \boldsymbol{\theta} \mid \mathcal{H}_k)$ under H-PSMG requires $O\left(\dim(P_0^P) + n |\Theta|\right)$ storage, as opposed to the naive $O\left(\dim(P_0^P) \cdot |\Theta|^n\right)$ that a non-factored posterior would require. Each Bayesian update costs $O(H)$ likelihood evaluations per module.
\end{corollary}

\begin{proposition}[Trajectory-Level Equivalence with Joint-PSRL]
\label{prop:bit-identical}
Consider an idealised Joint-PSRL algorithm that maintains the joint posterior $\mu_k^\Theta$ on $\Theta = \prod_i \Theta_i$ and at each episode samples $\hat{\boldsymbol\theta}_k \sim \mu_k^\Theta$. Under the assumptions of Theorem~\ref{thm:decoupling} and \emph{coupled RNG} (formalised below), the trajectories produced by H-PSMG and Joint-PSRL are pathwise identical: $\tau_k^{\mathrm{H\text{-}PSMG}} = \tau_k^{\mathrm{Joint\text{-}PSRL}}$ for every $k \geq 1$ almost surely. In particular the cumulative regrets agree exactly:
\begin{equation}
\mathrm{Reg}_i^B\big[\textnormal{H-PSMG}\big](K) = \mathrm{Reg}_i^B\big[\textnormal{Joint-PSRL}\big](K) \quad \forall i, K, \text{ a.s.}
\end{equation}

The coupling is constructed as follows. At episode $k$, both algorithms draw a single shared seed $u_k = (u_k^P, u_k^{\Theta, 1}, \ldots, u_k^{\Theta, n}) \sim \mathrm{Unif}([0,1]^{n+1})$. H-PSMG uses $u_k^P$ to sample $\hat P_k = F_{\mu_k^P}^{-1}(u_k^P)$ and $u_k^{\Theta, i}$ to sample $\hat\theta_i^k = F_{\mu_k^{\Theta, i}}^{-1}(u_k^{\Theta, i})$ via the inverse-CDF transform. Joint-PSRL uses $u_k^P$ identically for $\hat P_k$ and the joint draw $\hat{\boldsymbol\theta}_k = F_{\mu_k^\Theta}^{-1}\big(u_k^{\Theta, 1}, \ldots, u_k^{\Theta, n}\big)$ via the per-coordinate inverse-CDF of the product factorisation.
\end{proposition}

\begin{proof}
By Theorem~\ref{thm:decoupling}, $\mu_k^\Theta(\theta_1, \ldots, \theta_n) = \prod_i \mu_k^{\Theta, i}(\theta_i)$ at every $k$. The inverse-CDF of a product distribution applied coordinate-wise yields independent samples from the marginals. Hence the two algorithms produce identical $\hat P_k$ and identical $(\hat\theta_1^k, \ldots, \hat\theta_n^k)$, hence identical CCE policies $\pi_k = \mathrm{CCE}(\hat P_k, \hat{\boldsymbol\theta}_k)$ (Assumption~\ref{ass:CSR}: deterministic oracle), hence identical trajectories $\tau_k$ under shared environment randomness. Identical trajectories imply identical posteriors at $k+1$ by Bayes' rule, and the argument repeats inductively.
\end{proof}

Proposition~\ref{prop:bit-identical} formalises that H-PSMG and Joint-PSRL produce bit-identical trajectories under matched RNG on any HT-MG-BC instance. The proposition shows that H-PSMG is not merely a heuristic approximation of Joint-PSRL but its exact statistical equivalent at $O(n|\Theta_i|)$ rather than $O(|\Theta|^n)$ memory. The qualitative implication is that all regret guarantees and lower bounds derived for H-PSMG transfer verbatim to Joint-PSRL on HT-MG-BC instances.

\begin{remark}
The decoupling is more than a notational convenience. It separates two distinct sources of uncertainty into modular inference problems. Model uncertainty is captured by $\mu^P$, which behaves exactly as in single-agent PSRL and inherits all known concentration tools \citep{agrawal2017optimistic}. Type uncertainty is captured by $\mu^{\Theta,i}$, which is a categorical Bayesian estimation problem amenable to information-theoretic analysis \citep{russo2014learning}. This separation is the basis of the regret decomposition outlined in Section~\ref{sec:regret-roadmap}.
\end{remark}

\section{Numerical Illustration: Two-Agent Bayesian Prisoner's Dilemma}
\label{sec:example}

To make the decoupling concrete we work through a minimal instance numerically. The example is a stateless two-agent two-type matrix game played for $K$ rounds.

\subsection{Setup}

Set $|\mathcal{S}| = 1$, $\mathcal{A}_1 = \mathcal{A}_2 = \{C, D\}$, and $\Theta_i = \{P, S\}$ for $i \in \{1, 2\}$, where $P$ denotes pro-social and $S$ denotes selfish. The reward of agent $i$ as a function of its type is
\begin{equation}
r_i^S(a_i, a_{-i}) = \mathbbm{1}[a_i = D], \quad r_i^P(a_i, a_{-i}) = \mathbbm{1}[a_i = C \wedge a_{-i} = C].
\end{equation}
A selfish agent prefers defection. A pro-social agent rewards mutual cooperation. The horizon is $H = 1$ and there is no transition uncertainty, so the model component $\mu^P$ is trivial and the example isolates the type-inference component.

The prior is uniform: $\mu_1^{\Theta,1}(P) = \mu_1^{\Theta,1}(S) = 0.5$ and similarly for agent $2$, so $\mu_1^\Theta(\boldsymbol{\theta}) = 0.25$ for each $\boldsymbol{\theta} \in \Theta$. The true type profile is $\boldsymbol{\theta}^* = (P, S)$.

\subsection{Episode 1}

The algorithm draws $\hat{\boldsymbol{\theta}}_1 \sim \mu_1^\Theta$ uniformly. Suppose the realization is $\hat{\boldsymbol{\theta}}_1 = (P, P)$. The CCE of the all-pro-social game is the pure strategy $(C, C)$. The algorithm executes $(C, C)$, and rewards are realized in the true world: $r_1 = r_1^P(C, C) = 1$ and $r_2 = r_2^S(C, \cdot) = 0$.

The reward likelihoods are
\begin{equation}
L_{\Theta,1}(\tau_1; \theta_1) = P(r_1 = 1 \mid a = (C,C), \theta_1) = \begin{cases} 1 & \theta_1 = P \\ 0 & \theta_1 = S \end{cases}
\end{equation}
\begin{equation}
L_{\Theta,2}(\tau_1; \theta_2) = P(r_2 = 0 \mid a = (C,C), \theta_2) = \begin{cases} 0 & \theta_2 = P \\ 1 & \theta_2 = S \end{cases}.
\end{equation}

Table~\ref{tab:posterior-update} tabulates the posterior update.

\begin{table}[ht]
\caption{Posterior update after episode 1. Joint posterior is computed by Bayes' rule, then verified against the product of marginals.}
\label{tab:posterior-update}
\begin{center}
\begin{tabular}{cccccc}
\toprule
$\theta_1$ & $\theta_2$ & $\mu_1^\Theta(\boldsymbol{\theta})$ & $L_{\Theta,1}$ & $L_{\Theta,2}$ & Unnormalized posterior \\
\midrule
$P$ & $P$ & $0.25$ & $1$ & $0$ & $0$ \\
$P$ & $S$ & $0.25$ & $1$ & $1$ & $0.25$ \\
$S$ & $P$ & $0.25$ & $0$ & $0$ & $0$ \\
$S$ & $S$ & $0.25$ & $0$ & $1$ & $0$ \\
\midrule
\multicolumn{3}{r}{Normalized} & & & \\
$P$ & $S$ & & & & $1.00$ \\
others & & & & & $0.00$ \\
\bottomrule
\end{tabular}
\end{center}
\end{table}

The joint posterior puts probability $1$ on the true type profile $(P, S)$. We verify the decoupling by computing the marginals:
\begin{align}
\mu_2^{\Theta,1}(P) &= \mu_1^{\Theta,1}(P) \cdot L_{\Theta,1}(\tau_1; P) / Z_1 = 0.5 \cdot 1 / 0.5 = 1, \\
\mu_2^{\Theta,1}(S) &= 0, \\
\mu_2^{\Theta,2}(P) &= \mu_1^{\Theta,2}(P) \cdot L_{\Theta,2}(\tau_1; P) / Z_2 = 0.5 \cdot 0 / 0.5 = 0, \\
\mu_2^{\Theta,2}(S) &= 1.
\end{align}
The product $\mu_2^{\Theta,1} \cdot \mu_2^{\Theta,2}$ equals the unnormalized joint $\delta_{(P,S)}$, confirming the decoupling in Theorem~\ref{thm:decoupling}.

\subsection{Noisy Observation Variant}

To show that decoupling is not an artifact of deterministic rewards, suppose now the reward signal flips with probability $\epsilon = 0.1$. The likelihoods become
\begin{equation}
L_{\Theta,1}(\tau_1; \theta_1) = \begin{cases} 0.9 & \theta_1 = P \\ 0.1 & \theta_1 = S \end{cases}, \quad
L_{\Theta,2}(\tau_1; \theta_2) = \begin{cases} 0.1 & \theta_2 = P \\ 0.9 & \theta_2 = S \end{cases}.
\end{equation}
The unnormalized joint posterior is the outer product $\mu_1^\Theta \cdot L_{\Theta,1} \cdot L_{\Theta,2}$ which still factorizes as $\mu_2^{\Theta,1} \cdot \mu_2^{\Theta,2}$ with $\mu_2^{\Theta,1}(P) = 0.9, \mu_2^{\Theta,1}(S) = 0.1$ and $\mu_2^{\Theta,2}(S) = 0.9, \mu_2^{\Theta,2}(P) = 0.1$. Concentration is partial after one observation but the factorization is exact.

\section{Necessity of the Structural Assumptions}
\label{sec:necessity}

We show that both Assumptions~\ref{ass:TI} and \ref{ass:RL} are necessary in the sense that violation of either causes the decoupling to break.

\subsection{Counter-Example I: Type-Dependent Transitions}

Suppose Assumption~\ref{ass:TI} fails. Concretely, consider a two-state HT-MG where $P(s' \mid s, a, \boldsymbol{\theta})$ varies with $\boldsymbol{\theta}$. The trajectory likelihood becomes
\begin{equation}
P(\tau_k \mid P, \boldsymbol{\theta}, \mathcal{H}_{k-1}) = c_k(\tau_k) \cdot \prod_{h=1}^H P(s_{h+1}^k \mid s_h^k, a_h^k, \boldsymbol{\theta}) \cdot \prod_{h, i} P(r_h^{k,i} \mid s_h^k, a_h^k, \theta_i).
\end{equation}
The transition factor is now jointly a function of $P$ and $\boldsymbol{\theta}$ and does not split. Posterior $P(P, \boldsymbol{\theta} \mid \mathcal{H}_k)$ is no longer factorable into a $P$-marginal and a $\boldsymbol{\theta}$-marginal. Empirically, observing a transition $(s, a) \to s'$ simultaneously updates beliefs about both the kernel and the types, with no way to attribute the change to one source alone.

\subsection{Counter-Example II: Non-Local Reward}

Suppose Assumption~\ref{ass:RL} fails. Consider a social-comparison reward $r_1(s, a, \boldsymbol{\theta}) = r_1^{\theta_1}(s, a) - \alpha \cdot r_2^{\theta_2}(s, a)$ for some $\alpha > 0$, so that agent $1$'s utility includes an envy term proportional to agent $2$'s realized return under its own type. Then
\begin{equation}
L_{\Theta,1}(\tau_k; \boldsymbol{\theta}) = \prod_{h=1}^H P\left(r_h^{k,1} - \alpha r_h^{k,2,\theta_2} \mid s_h^k, a_h^k, \theta_1\right)
\end{equation}
contains $\theta_2$ inside the agent-$1$ likelihood factor. The likelihood factorization in Lemma~\ref{lem:likelihood} fails. The posterior couples $\theta_1$ and $\theta_2$ at every update step.

These two counter-examples confirm that the decoupling theorem is tight in its assumptions and provide guidance on when the centralized H-PSMG framework is appropriate.

\section{From Decoupling to Regret: A Roadmap}
\label{sec:regret-roadmap}

The decoupling theorem is the structural foundation for a regret bound on H-PSMG. We outline the derivation strategy in this section and provide the rigorous proof in Appendix~\ref{app:regret-decomp}.

\subsection{Per-Episode Regret Decomposition}

For episode $k$, the per-episode regret of agent $i$ is
\begin{equation}
\Delta_i^k := V_{i,1}^{\pi^*}(s_1^k; P^*, \theta_i^*) - V_{i,1}^{\pi_k}(s_1^k; P^*, \theta_i^*),
\end{equation}
where $\pi^* = \mathrm{CCE}(P^*, \boldsymbol{\theta}^*)$ and $\pi_k = \mathrm{CCE}(\hat{P}_k, \hat{\boldsymbol{\theta}}_k)$. Two technical steps reduce this to a tractable form. First, the posterior sampling equivalence (Lemma~\ref{lem:ps-equivalence} in Appendix~\ref{app:regret-decomp}) implies $\mathbb{E}[V_{i,1}^{\pi^*}(P^*, \theta_i^*)] = \mathbb{E}[V_{i,1}^{\pi_k}(\hat{P}_k, \hat{\theta}_i^k)]$, so the expected regret reduces to a comparison between two evaluations of the same policy $\pi_k$:
\begin{equation}
\mathbb{E}[\Delta_i^k \mid \mathcal{H}_{k-1}] = \mathbb{E}\big[V_{i,1}^{\pi_k}(s_1^k; \hat{P}_k, \hat{\theta}_i^k) - V_{i,1}^{\pi_k}(s_1^k; P^*, \theta_i^*) \,\big|\, \mathcal{H}_{k-1}\big].
\end{equation}
Second, applying the value-difference lemma for HT-MG (Lemma~\ref{lem:value-diff}) telescopes this gap into a transition-misspec term and a reward-misspec term:
\begin{equation}
\mathbb{E}[\Delta_i^k \mid \mathcal{H}_{k-1}] = \underbrace{\mathbb{E}\Big[\sum_{h=1}^H \mathbb{E}_{d_h^{\pi_k, \hat{P}_k}}\big[r_i(\cdot, \hat{\theta}_i^k) - r_i(\cdot, \theta_i^*)\big]\Big]}_{\text{reward-misspec } \Delta_i^{k,R}} + \underbrace{\mathbb{E}\Big[\sum_{h=1}^H \mathbb{E}_{d_h^{\pi_k, \hat{P}_k}}\big\langle \hat{P}_k - P^*, V_{i,h+1}^{\pi_k}(P^*, \theta_i^*)\big\rangle\Big]}_{\text{transition-misspec } \Delta_i^{k,P}}.
\end{equation}
Theorem~\ref{thm:decoupling} ensures the two sources of error are governed by independent posteriors $\mu^{\Theta,i}$ and $\mu^P$, enabling separate analysis.

\subsection{Bounding the Transition Component}

By Theorem~\ref{thm:decoupling}, the transition posterior $\mu^P$ behaves exactly as in single-agent PSRL operating on the joint MDP induced by the joint action space $\mathcal{A} = \prod_i \mathcal{A}_i$. We obtain two complementary bounds. A direct Hölder + Dirichlet $\ell_1$-concentration argument gives
\begin{equation}
\sum_{k=1}^K \mathbb{E}\big[|\Delta_i^{k,P}|\big] \leq \tilde{O}\Big(H^{3/2} \cdot S \cdot \sqrt{A_{\mathrm{joint}} K}\Big), \quad A_{\mathrm{joint}} := \prod_{i=1}^n |\mathcal{A}_i|,
\end{equation}
matching \citet{osband2013more}'s Hoeffding-level Bayesian PSRL bound. Replacing the Hölder step with the variance-aware Bernstein inner-product bound \citep{azar2017minimax,lu2019information} and exploiting the law of total variance trims one factor of $\sqrt{S}$:
\begin{equation}
\sum_{k=1}^K \mathbb{E}\big[|\Delta_i^{k,P}|\big] \leq \tilde{O}\Big(H^{3/2}\sqrt{S \cdot A_{\mathrm{joint}} \cdot K}\Big).
\end{equation}
This matches the frequentist minimax-optimal $\tilde{O}(\sqrt{H^3SAK})$ rate of \citet{azar2017minimax} for single-agent finite-horizon RL. Proposition~\ref{prop:trans-hoeffding} (Hoeffding) and Proposition~\ref{prop:trans-bernstein} (Bernstein-refined) in Appendix~\ref{app:regret-decomp} provide the detailed arguments.

\subsection{Bounding the Type Component}

The reward-misspec term is the cumulative gap between the posterior-sampled type and the true type. Since rewards lie in $[0, 1]$ and $r_i(\cdot, \hat{\theta}_i^k) = r_i(\cdot, \theta_i^*)$ whenever $\hat{\theta}_i^k = \theta_i^*$ (Assumption~\ref{ass:RL}), we have $|\Delta_i^{k,R}| \leq H \cdot \mathbbm{1}[\hat{\theta}_i^k \neq \theta_i^*]$. Applying the information-ratio framework of \citet{russo2016information} to the categorical type-inference problem yields
\begin{equation}
\sum_{k=1}^K \mathbb{E}[\Delta_i^{k,R}] \leq \sqrt{\bar{\Gamma} \cdot \log|\Theta_i| \cdot K},
\end{equation}
where $\bar{\Gamma}$ is a uniform bound on the per-episode information ratio. For the finite-arm case, $\bar{\Gamma} = O(H^2 |\Theta_i|)$, giving $O(H\sqrt{|\Theta_i| \log|\Theta_i| \cdot K})$. Under the stronger information-rate condition $I_k(\theta_i^*; \tau_k) \geq \rho > 0$, the cumulative reward-misspec regret is $O(H \log|\Theta_i| / \rho)$, constant in $K$. Proposition~\ref{prop:type-bound} in Appendix~\ref{app:regret-decomp} contains the proof.

\subsection{Main Regret Bound}

Combining the two components and the CCE oracle approximation gives the formal regret bound, which we state here and prove rigorously in Appendix~\ref{app:regret-decomp}.

\begin{theorem}[Bayesian Regret of H-PSMG, informal]
\label{thm:main-regret-informal}
Under Assumptions~\ref{ass:TI}, \ref{ass:RL}, \ref{ass:PF}, \ref{ass:ISI}, \ref{ass:CSR}, a Dirichlet prior on $P^*$, and a bounded type information ratio $\bar{\Gamma} = O(H^2 |\Theta_i|)$, the Bayesian regret of agent $i$ over $K$ episodes satisfies
\begin{equation}
\label{eq:main-regret-sketch}
\mathbb{E}\big[\mathrm{Reg}_i^B(K)\big] \leq \tilde{O}\Big(H^{3/2}\sqrt{S \cdot A_{\mathrm{joint}} \cdot K}\Big) + O\Big(H \sqrt{|\Theta_i| \log|\Theta_i| \cdot K}\Big) + K \cdot \epsilon_{\mathrm{CCE}},
\end{equation}
where $A_{\mathrm{joint}} = \prod_i |\mathcal{A}_i|$ is the joint action space size and $\epsilon_{\mathrm{CCE}}$ is the CCE oracle approximation error.
\end{theorem}

The bound separates the cost of learning the transition kernel from the cost of identifying the type profile. The state-space size $S$ enters only the model-learning component, reflecting the agent-level decoupling of Theorem~\ref{thm:decoupling}: each agent's type inference depends only on its own type space $\Theta_i$ rather than on the full joint type space $\Theta$. Appendix~\ref{app:regret-decomp} provides the full proof, including the value-difference lemma for HT-MG, the information-ratio analysis of type inference, and the Dirichlet concentration argument for the transition component.

\begin{corollary}[Gap-Dependent Rate under Identifiability]
\label{cor:gap-dependent}
If in addition Assumption~\ref{ass:TID} holds with parameter $\rho > 0$, then the type-inference component collapses to a constant in $K$:
\begin{equation}
\sum_{k=1}^K \mathbb{E}[\Delta_i^{k,R}] \leq O\bigg(\frac{H \log|\Theta_i|}{\rho}\bigg).
\end{equation}
Consequently, the Bayesian regret simplifies to
\begin{equation}
\mathbb{E}\big[\mathrm{Reg}_i^B(K)\big] \leq \tilde{O}\Big(H^{3/2}\sqrt{S A_{\mathrm{joint}} K}\Big) + O\bigg(\frac{H \log|\Theta_i|}{\rho}\bigg) + K \cdot \epsilon_{\mathrm{CCE}}.
\end{equation}
In the asymptotic regime $K \to \infty$ with $\epsilon_{\mathrm{CCE}} = O(K^{-1/2})$, the bound is dominated by transition learning at rate $\tilde{O}(\sqrt{K})$. The type-inference cost is paid only during a short burn-in period of length $O(\log|\Theta_i| / \rho)$.
\end{corollary}

The corollary captures an important practical regime: when types are sufficiently distinct (TID with $\rho$ bounded away from zero), the cost of heterogeneity is essentially a one-time identification cost rather than a recurring exploration penalty. Proof appears in Appendix~\ref{app:regret-decomp}, immediately following Proposition~\ref{prop:type-bound}.

\subsection{Rate Separation from Heuristic and Type-Agnostic Baselines}
\label{sec:rate-separation}

The $\tilde{O}(\sqrt{K})$ rate of Theorem~\ref{thm:main-regret-informal} is not just a constant-factor improvement over heuristic baselines: the gap is in the asymptotic rate. We make this precise.

\begin{theorem}[Rate Separation]
\label{thm:rate-separation}
There exists an HT-MG instance $\mathcal{G}^\star$ satisfying Assumptions~\ref{ass:TI}, \ref{ass:RL}, \ref{ass:PF}, \ref{ass:ISI}, \ref{ass:CSR}, and a uniform prior over types $\mu_0^{\Theta,i} = \mathrm{Unif}(\Theta_i)$, on which the cumulative Bayesian regret of three algorithms satisfies:
\begin{enumerate}
\item[(i)] \textbf{H-PSMG (Algorithm~\ref{alg:hpsmg}):} $\mathbb{E}[\mathrm{Reg}_i^B(K)] = \tilde{O}(\sqrt{K})$ for all $K \geq 1$.
\item[(ii)] \textbf{MAP-Type-Greedy} (substitute the MAP estimate $\hat{\theta}_i^k = \arg\max_\theta \mu_k^{\Theta,i}(\theta)$ for the posterior sample in Algorithm~\ref{alg:hpsmg}): on a positive-probability event $\mathcal{E}_{\mathrm{trap}}$ (over initial randomness), $\mathbb{E}[\mathrm{Reg}_i^B(K) \mid \mathcal{E}_{\mathrm{trap}}] = \Omega(K)$. In particular, $\mathbb{E}[\mathrm{Reg}_i^B(K)] = \Omega(K)$.
\item[(iii)] \textbf{PSRL-NoType} (sample $\hat{\boldsymbol{\theta}}_k \sim \mathrm{Unif}(\Theta)$ every episode, never update the type posterior): $\mathbb{E}[\mathrm{Reg}_i^B(K)] = \Omega(K)$.
\end{enumerate}
\end{theorem}

\begin{proof}[Proof Sketch]
We construct $\mathcal{G}^\star$ explicitly as the \emph{exploration-trap game}: 5 states $\{0, 1, 2, 3, 4\}$, $n = 2$ agents with action set $\{L, R\}$, types $\Theta_i = \{A, B\}$, horizon $H$, deterministic majority-vote transitions ($(R, R) \mapsto s+1$, $(L, L) \mapsto s-1$, mixed $\mapsto s$, clipped at boundaries), initial state $s_1 = 0$, and rewards
\begin{align}
r_i(s, a, A) &= 0.5 \quad \text{for all } s, \\
r_i(s, a, B) &= \begin{cases} 0.5 & \text{if } s \in \{0, 1, 2, 3\}, \\ 1.0 & \text{if } s = 4, \end{cases}
\end{align}
with small Gaussian observation noise $\sigma \to 0$. The true types are $\boldsymbol{\theta}^\star = (B, B)$.

\textbf{(i) H-PSMG.} This is the direct application of Theorem~\ref{thm:main-regret-informal} to $\mathcal{G}^\star$.

\textbf{(ii) MAP-Type-Greedy.} Let $\mathcal{E}_{\mathrm{trap}}$ be the event that the initial MAP estimate of agent 1 is $A$. Under a uniform prior with deterministic tie-breaking (e.g., $\arg\max$ returning the lowest-index maximizer), this event holds with probability $1$; with random uniform tie-breaking it has probability $\geq (1/|\Theta_1|)^n = 1/4$.

On $\mathcal{E}_{\mathrm{trap}}$, by induction on $k$: (1) $\hat{\boldsymbol{\theta}}_{\mathrm{MAP}}^k = (A, A)$; (2) under type $A$, $V_h^{\pi}(s; A) = 0.5(H-h+1)$ is independent of state and policy (uniform reward), so all joint actions yield identical $Q$-values and the CCE-with-Nash-selection picks $\pi(L, L)$ at every $(s, h)$ by tie-breaking; (3) rollout from $s_1 = 0$ under deterministic transitions stays at $s = 0$ for all $H$ steps; (4) at state $0$ both candidate types predict $r_h^i = 0.5$, so the log-likelihood ratio $\log \mu_{k+1}^{\Theta,i}(A) / \mu_{k+1}^{\Theta,i}(B) = \log \mu_k^{\Theta,i}(A) / \mu_k^{\Theta,i}(B) + O(\sigma)$, which converges in the limit $\sigma \to 0$; (5) the MAP estimate remains $A$.

The per-episode value of MAP-Type-Greedy on $\mathcal{E}_{\mathrm{trap}}$ is $V^{\mathrm{MAP}} = 0.5 H$ per agent. The oracle value satisfies $V^\star \geq H - 4 + 1 \cdot (H - 4) = H - 4 + (H - 4)$ approximately (the optimal $(B, B)$ policy reaches state $4$ in $4$ steps and earns $r = 1$ for the remaining $H - 4$ steps, with $r = 0.5$ during the traverse), giving $V^\star \geq 0.5 \cdot 4 + 1.0 \cdot (H - 4) = H - 2$ per agent for $H \geq 4$. The per-episode regret is $V^\star - V^{\mathrm{MAP}} \geq H - 2 - 0.5 H = 0.5 H - 2 > 0$ for $H > 4$. Summed over $K$ episodes, $\mathbb{E}[\mathrm{Reg}_i^B(K) \mid \mathcal{E}_{\mathrm{trap}}] \geq (0.5 H - 2) K = \Omega(K)$, completing the claim.

\textbf{(iii) PSRL-NoType.} The algorithm samples $\hat{\boldsymbol{\theta}}_k \sim \mathrm{Unif}(\Theta_1 \times \Theta_2)$ every episode, independent of history. Let $V(\boldsymbol{\theta}) := \mathbb{E}^{\pi(\boldsymbol{\theta}), P^\star}[\sum_h r_i(s_h, a_h, \theta_i^\star)]$ be the value of the CCE for sampled types $\boldsymbol{\theta}$ evaluated under the true types $\boldsymbol{\theta}^\star = (B, B)$. Then the expected per-episode value is $\bar V = \mathbb{E}_{\boldsymbol{\theta} \sim \mathrm{Unif}}[V(\boldsymbol{\theta})]$. By a direct case analysis, $\bar V < V^\star$ for the trap game, since the three non-$(B,B)$ joint types ($\frac{3}{4}$ of samples) induce policies that fail to reach state $4$. The per-episode regret is the positive constant $V^\star - \bar V$, giving $\mathbb{E}[\mathrm{Reg}_i^B(K)] = (V^\star - \bar V) K = \Omega(K)$.
\end{proof}

The theorem makes precise the qualitative claim that Bayesian posterior sampling on types provides exploration that heuristic alternatives lack. The proof construction is the exploration-trap game $\mathcal{G}^\star$ used to drive the deep-trap family $\mathcal{G}_m^{\mathrm{trap}}$ analysed in Section~\ref{sec:hpsmg-plus}; the empirical counterpart of this separation, in the LLM-tier setting where action spaces and reward observations come from prompted backbones rather than analytic kernels, is reported in Section~\ref{sec:experiments-llm}.

\subsection{Regret guarantees of H-PSMG\textsuperscript{+}}
\label{sec:hpsmg-plus}

The algorithm H-PSMG\textsuperscript{+} (Algorithm~\ref{alg:hpsmg-plus}) was introduced in Section~\ref{sec:hpsmg-plus-algo}. This subsection establishes its two regret guarantees: asymptotic match with H-PSMG (worst-case rate unchanged) and exponential-in-$|\Theta|^n$ burn-in improvement on the deep-trap family. Both are direct corollaries of the type-discrimination bonus structure and the posterior decoupling theorem.

\begin{theorem}[Regret of H-PSMG\textsuperscript{+}]
\label{thm:hpsmg-plus}
Let $\beta > 0$. Under the assumptions of Theorem~\ref{thm:main-regret-informal}:

\textbf{(i) Asymptotic match.} On any HT-MG, the Bayesian regret of H-PSMG\textsuperscript{+} satisfies
\begin{equation}
\mathbb{E}\big[\mathrm{Reg}_i^B(K)\big] \leq \tilde{O}\big(H^{3/2}\sqrt{S A_{\mathrm{joint}} K} + H\sqrt{|\Theta_i| K}\big),
\end{equation}
identical to the H-PSMG bound of Theorem~\ref{thm:main-regret-informal}.

\textbf{(ii) Exponential-in-$|\Theta|^n$ burn-in improvement.} Let $\mathcal{G}_m^{\mathrm{trap}}$ denote the family of deep-trap HT-MGs parameterised by per-agent type-space size $|\Theta_i| = m$ (Appendix~\ref{app:hpsmg-plus} for the precise definition: $n$ agents, $S = m+2$ states, deterministic unanimity-vote transitions, horizon $H$, single jackpot type $\theta_{\mathrm{jack}}$ with $r_i(s^\star, a, \theta_{\mathrm{jack}}) = 1$ at the rightmost state and $r_i = 0.5$ otherwise, uniform prior, true types $\theta_i^\star = \theta_{\mathrm{jack}}$). Under Assumption~\ref{ass:cce-tiebreak} (welfare-maximising CCE oracle, lowest-index tie-break), the expected burn-in regrets are:
\begin{align}
\textnormal{H-PSMG / Joint-PSRL:} \quad &\mathbb{E}\big[\mathrm{Reg}_i^B(\tau_{\textnormal{H-PSMG}})\big] = \Theta(m^n \cdot H), \\
\textnormal{H-PSMG\textsuperscript{+} (any } \beta > 0\textnormal{):} \quad &\mathbb{E}\big[\mathrm{Reg}_i^B(\tau_{\textnormal{H-PSMG}^+})\big] = O(H).
\end{align}
The burn-in separation is $\Omega(m^n)$, exponential in the number of agents.
\end{theorem}

\begin{proof}[Proof Sketch]
\textbf{(i) Asymptotic match.} By a CCE-optimality decomposition (Lemma~\ref{lem:bonus-decomp}, App~\ref{app:hpsmg-plus-asymptotic-proof}), the per-episode regret of H-PSMG\textsuperscript{+} relative to H-PSMG is bounded by the expected cumulative bonus collected on the rollout, $\beta H \cdot \mathbb{E}[\max_{s,a} D_k(s, a)]$. A collision-probability argument (Lemma~\ref{lem:collision-bound}) gives $D_k(s, a) \leq 2 \sum_i (1 - \mu_k^{\Theta, i}(\theta_i^\star))$, reducing the cumulative bonus to a sum of posterior-error probabilities. Under Assumption~\ref{ass:TID} (Hellinger gap $\rho > 0$), standard Bayesian posterior concentration on finite type spaces \citep{ghosal2007convergence} gives an exponential decay $\mathbb{E}[1 - \mu_k^{\Theta, i}(\theta_i^\star)] \leq |\Theta_i| \exp(-\rho H k / 2)$, so the cumulative bonus regret is bounded by a $K$-independent constant $O(\beta n |\Theta_i| / \rho)$ (Proposition~\ref{prop:hpsmg-plus-asymptotic}). This is absorbed into the leading $\tilde O(H\sqrt{|\Theta_i| K})$ type-inference term of Theorem~\ref{thm:main-regret-informal}.

\textbf{(ii) Exponential burn-in.} For H-PSMG: under uniform prior, $\Pr[\hat\theta_k = (\theta_{\mathrm{jack}})^n] = m^{-n}$ per episode, independent across $k$. Before any successful draw, every visited state $s \in \{0, 1, \ldots, S-2\}$ has identical predicted reward $0.5$ across all types, so observations carry no information about types (Lemma~\ref{lem:posterior-stuck}); the marginal stays uniform. The expected wait for the first informative draw is $\mathrm{Geom}(m^{-n})$ with mean $m^n$; each stuck episode contributes per-episode regret $\Theta(H)$, totalling $\Theta(m^n H)$ (Proposition~\ref{prop:vanilla-burnin}).

For H-PSMG\textsuperscript{+}: at uniform prior, the bonus evaluates to $D_k(s^\star, a) = n(m-1)/m^2 > 0$ at the jackpot state and $0$ elsewhere (Equation~\ref{eq:Dk-trap}). Backward induction with this bonus and deterministic unanimity transitions gives social-welfare gap
\begin{equation*}
W(a^R) - W(a^L) = n(H - S + 1)\big(0.5\alpha + \beta n p_k(1-p_k)\big) > 0
\end{equation*}
regardless of the sampled $\hat\theta$ (Lemma~\ref{lem:cce-selects-aR}), where $\alpha$ is the fraction of jackpot-sampled agents. Under Assumption~\ref{ass:cce-tiebreak}, the CCE oracle returns $a^R$; the trajectory reaches $s^\star$ deterministically; the marginal concentrates on $\theta_{\mathrm{jack}}$ within one episode (Proposition~\ref{prop:hpsmg-plus-burnin}). Burn-in regret is $O(H)$, independent of $m$ and $n$.

The full formal proof is in Appendix~\ref{app:hpsmg-plus}.
\end{proof}

\paragraph{Comparison with Joint-PSRL.} Theorem~\ref{thm:hpsmg-plus} formalises the algorithmic gap between H-PSMG\textsuperscript{+} and Joint-PSRL. Although vanilla H-PSMG and Joint-PSRL produce bit-identical trajectories under matched RNG (a consequence of Theorem~\ref{thm:decoupling}), the bonus $D_k$ requires per-agent factored marginals, which Joint-PSRL only computes by marginalising the joint posterior at $O(|\Theta|^n)$ cost per query. The decoupling theorem therefore enables three distinct improvements over Joint-PSRL:
\begin{itemize}
\item \textbf{Storage.} $O(n|\Theta_i|)$ instead of $\Omega(|\Theta|^n)$ posterior memory (Corollary~\ref{cor:storage}). For $|\Theta_i| = 4$, the joint posterior grows from $16$ entries at $n=2$ to $1024$ entries at $n=5$, a $64\times$ gap that we measure empirically on an LLM substrate in Section~\ref{sec:experiments-llm}.
\item \textbf{Regret analysis.} $\tilde O(\sqrt{|\Theta_i| K})$ type-inference rate instead of $\tilde O(\sqrt{|\Theta|^n K})$, the generic Bayesian-RL guarantee (Theorem~\ref{thm:lower-no-pf}, Section~\ref{sec:exponential-improvement}).
\item \textbf{Algorithmic refinement.} Exponential-in-$|\Theta|^n$ burn-in regret improvement of H-PSMG\textsuperscript{+} over Joint-PSRL on the deep-trap family (Theorem~\ref{thm:hpsmg-plus}(ii)).
\end{itemize}
The empirical separation, with H-PSMG\textsuperscript{+}'s burn-in benefit growing monotonically in $|\Theta_i|$ and $n$ on an LLM substrate, is reported in Section~\ref{sec:experiments-llm}.

\subsection{Exponential improvement over generic Bayesian RL via decoupling}
\label{sec:exponential-improvement}

Theorem~\ref{thm:main-regret-informal} produces a regret bound whose type-inference cost depends on $|\Theta_i|$, the per-agent type-space size, rather than on $|\Theta| = \prod_i |\Theta_i|$, the joint type-space size. A natural question is whether this dependence is intrinsic to the HT-MG problem or an artefact of our structural assumptions. We answer it by showing that without Assumption~\ref{ass:PF}, no algorithm can match the bound: the best achievable Bayesian regret depends on $|\Theta|^n$, exponentially worse than ours.

\begin{theorem}[Exponential Lower Bound Without Prior Factorisation]
\label{thm:lower-no-pf}
Let $\mathcal{F}^{\mathrm{noPF}}_{n, m, S, A_{\mathrm{joint}}, H}$ denote the family of HT-MG instances with $n$ agents, $|\Theta_i| = m$ for each $i$, state-space size $S$, joint-action-space size $A_{\mathrm{joint}}$, horizon $H$, satisfying Assumptions~\ref{ass:TI}, \ref{ass:RL}, \ref{ass:ISI}, \ref{ass:CSR}, but for which the joint type prior $\mu_0^\Theta$ is permitted to be an arbitrary distribution over $\Theta = \prod_i \Theta_i$ (i.e., \ref{ass:PF} is dropped). For any algorithm $\mathfrak{A}$,
\begin{equation}
\sup_{\mathcal{G} \in \mathcal{F}^{\mathrm{noPF}}}\ \mathbb{E}_{\mathcal{G}, \mathfrak{A}}\big[\mathrm{Reg}_i^B(K)\big] = \Omega\big(\sqrt{m^n \cdot K}\big).
\end{equation}
\end{theorem}

\begin{proof}[Proof Sketch]
We reduce a Bayesian multi-armed bandit with $M = m^n$ arms to an HT-MG without prior factorisation. Take $S = 1$, $H = 1$, $A_{\mathrm{joint}} = M$. For each joint type $\theta \in \Theta$, choose a one-to-one map $a^\star \colon \Theta \to \mathcal{A}_{\mathrm{joint}}$ and define the reward as $r(a, \theta) = \mathbb{1}[a = a^\star(\theta)]$. Place the prior $\mu_0^\Theta$ uniform on $\Theta$ but not factored across agents (e.g., concentrated on a uniformly random permutation of $\Theta$). Conditional on $\theta$ the unique optimal joint action is $a^\star(\theta)$; without prior factorisation, observing one agent's marginal reward provides no information about another agent's type that is not already captured by the joint posterior over $\Theta$.

The resulting problem is a Bayesian $M$-armed bandit with uniform prior over the identity of the optimal arm. By the standard Bayesian regret lower bound for multi-armed bandits \citep[Theorem 15.2]{lattimore2020bandit},
$\mathbb{E}_{\mathfrak{A}}[\mathrm{Reg}(K)] = \Omega(\sqrt{MK}) = \Omega(\sqrt{m^n K})$.
Per-agent regret equals total regret in this single-agent-reduces-to-bandit construction up to a factor of $n$, completing the claim.
\end{proof}

Theorem~\ref{thm:lower-no-pf} combined with Theorem~\ref{thm:main-regret-informal} establishes an exponential separation. With prior factorisation (Assumption~\ref{ass:PF}), H-PSMG achieves type-inference regret $\tilde O(H\sqrt{m K})$; without it, no algorithm achieves better than $\Omega(\sqrt{m^n K})$. The improvement factor in the bound is $m^{(n-1)/2}$, exponential in the number of agents. The structural assumption set TI + RL + PF is therefore not merely a convenience for analysis but a necessary condition for sublinear-in-$n$ dependence on the joint type space.

\begin{remark}[Comparison with Joint-PSRL]
A particularly useful consequence is the regret bound for Joint-PSRL implied by our analysis. The naive Bayesian-RL analysis of Joint-PSRL on the joint hypothesis class $\Theta$ would invoke the generic information-ratio bound for posterior sampling on a hypothesis class of size $|\Theta|^n$, yielding a regret bound of $\tilde O(H\sqrt{|\Theta|^n K})$. Our Theorem~\ref{thm:decoupling} (Decoupling) shows that under TI + RL + PF the joint posterior is the product of marginals at every episode, which when combined with the per-agent information ratio of Lemma~\ref{lem:info-ratio} reduces the Joint-PSRL regret bound to the same $\tilde O(H\sqrt{|\Theta_i| K})$ as H-PSMG (Theorem~\ref{thm:main-regret-informal}). Thus our decoupling theorem provides an exponential-in-$n$ sharpening of the Joint-PSRL regret guarantee at no cost in storage or computation for Joint-PSRL (which is already storing the joint posterior); for H-PSMG, the additional benefit is the exponential storage reduction documented analytically by Corollary~\ref{cor:storage} and verified empirically on an LLM substrate in Section~\ref{sec:experiments-llm}.
\end{remark}

\subsection{Minimax lower bound and tightness of Theorem~\ref{thm:main-regret-informal}}
\label{sec:lower-bound-matching}

Within the family of HT-MGs satisfying all our assumptions, we establish a minimax lower bound that nearly matches Theorem~\ref{thm:main-regret-informal}.

\begin{theorem}[Minimax Bayesian Regret Lower Bound under PF]
\label{thm:lower-matching}
Let $\mathcal{F}^{\mathrm{PF}}_{n, m, S, A_{\mathrm{joint}}, H}$ denote the family of HT-MG instances satisfying Assumptions~\ref{ass:TI}, \ref{ass:RL}, \ref{ass:PF}, \ref{ass:ISI}, \ref{ass:CSR}. For any algorithm $\mathfrak{A}$ and parameters $K \geq SA_{\mathrm{joint}}H$ and $K \geq mH$, there exists an instance $\mathcal{G} \in \mathcal{F}^{\mathrm{PF}}$ with
\begin{equation}
\mathbb{E}_{\mathcal{G}, \mathfrak{A}}\big[\mathrm{Reg}_i^B(K)\big] = \Omega\Big(\sqrt{H^3 S A_{\mathrm{joint}} K} + \sqrt{|\Theta_i| H K}\Big).
\end{equation}
\end{theorem}

\begin{proof}[Proof Sketch]
The bound decomposes into a transition-learning component and a type-inference component, each established by reduction to a known lower bound.

\textbf{Transition component} $\Omega(\sqrt{H^3 S A_{\mathrm{joint}} K})$. Specialise to $|\Theta_i| = 1$, eliminating type uncertainty. The HT-MG reduces to a standard finite-horizon MDP with state space $\mathcal{S}$, joint action space $\mathcal{A}_{\mathrm{joint}}$, horizon $H$. The minimax lower bound for episodic RL of \citet[Theorem 1]{domingues2021episodic} gives expected cumulative regret $\Omega(\sqrt{H^3 S A K})$ for any algorithm; substituting $A = A_{\mathrm{joint}}$ completes this component.

\textbf{Type-inference component} $\Omega(\sqrt{|\Theta_i| H K})$. Specialise to $S = 1$, $A_{\mathrm{joint}} = |\Theta_i|$, $|\Theta_j| = 1$ for $j \neq i$ (no uncertainty about other agents' types). Let $r_i(s, a, \theta_i) = \mathbb{1}[a = a^\star(\theta_i)]$ for a one-to-one mapping $a^\star \colon \Theta_i \to \mathcal{A}_{\mathrm{joint}}$. The resulting problem is a stochastic multi-armed bandit with $|\Theta_i|$ arms and $T = KH$ rounds (each episode provides $H$ reward observations of the same type). The standard MAB minimax lower bound \citep[Theorem 15.2]{lattimore2020bandit} gives $\Omega(\sqrt{|\Theta_i| T}) = \Omega(\sqrt{|\Theta_i| HK})$.

The two constructions can be combined into a single HT-MG with both transition and type uncertainty; the resulting lower bound is the maximum of the two components, which equals their sum up to a constant factor.
\end{proof}

Comparing the lower bound of Theorem~\ref{thm:lower-matching} with the upper bound of Theorem~\ref{thm:main-regret-informal}:
\begin{center}
\small
\begin{tabular}{lll}
\toprule
Component & Upper bound (Thm.~\ref{thm:main-regret-informal}) & Lower bound (Thm.~\ref{thm:lower-matching}) \\
\midrule
Transition & $\tilde O(\sqrt{H^3 S A_{\mathrm{joint}} K})$ & $\Omega(\sqrt{H^3 S A_{\mathrm{joint}} K})$ \\
Type inference & $\tilde O(\sqrt{H^2 |\Theta_i| K})$ & $\Omega(\sqrt{H |\Theta_i| K})$ \\
\bottomrule
\end{tabular}
\end{center}
The transition term is tight modulo logarithmic factors. The type-inference term has a $\sqrt{H}$ gap, consistent with similar gaps in Bayesian-RL upper bounds and conjectured (but unresolved) to be closable via variance-aware Bernstein-type analyses. Theorem~\ref{thm:lower-no-pf} establishes that without PF this entire analysis collapses: the lower bound jumps from $\sqrt{|\Theta_i|}$ to $\sqrt{|\Theta|^n}$, an exponential-in-$n$ inflation. Our regret bound is thus near-optimal among algorithms that exploit the structural assumptions of HT-MG.

\section{Empirical Study}
\label{sec:experiments-llm}

We empirically validate H-PSMG\textsuperscript{+} and its underlying theory across
three benchmarks of increasing wildness. Section~\ref{sec:hp-spgg} introduces the
Heterogeneous-Persona Sequential Public Goods Game (HP-SPGG), a controlled LLM-MAS
benchmark we construct that satisfies all six HT-MG-BC assumptions by
construction, and reports cross-model results across four LLM backbones against
a $16$-method comparison suite spanning Bayesian type-aware algorithms (including
ours), prompted-LLM coordinators, external LLM-coordination methods adapted from
recent literature (ECON, A-ToM family), and type-agnostic baselines; the headline
figure shows a $12$-method representative subset covering all four families.
The same section also reports theory-driven experiments measuring backbone
posterior concentration speed and the predicted scaling of the bonus benefit with
type-space size and agent count. Section~\ref{sec:concordia} scales the evaluation
to the DeepMind Concordia substrates (Pub Coordination and Haggling), public
benchmarks from the NeurIPS 2024 Cooperative AI Contest
\citep{vezhnevets2023concordia, concordia_contest2025} that we did not design.
Section~\ref{sec:sotopia} reports results on SOTOPIA-Hard \citep{zhou2024sotopia},
a dyadic social-negotiation benchmark whose action space is open-ended natural
language and whose persona structure violates HT-MG-BC's prior-factorisation
assumption; we use SOTOPIA-Hard to map the boundary of our theory's applicability.
On the two HT-MG-BC-conformant substrates (HP-SPGG, Concordia), the H-PSMG family
attains the lowest non-oracle regret or the highest non-oracle focal score; on
SOTOPIA-Hard, a substrate-adapted variant of H-PSMG\textsuperscript{+} achieves the
highest evaluator score on every backbone.

\subsection{HP-SPGG: a controlled LLM-MAS benchmark for HT-MG-BC}
\label{sec:hp-spgg}

\paragraph{Setup.} A coordinator oversees $n = 3$ player LLMs across
$K = 20$ rounds. Each player's hidden type
$\theta_i \in \Theta_i$ is drawn without replacement from $|\Theta_i| = 4$
persona archetypes constructed in the SOTOPIA style \citep{zhou2024sotopia} and
spanning the empirical taxonomy of public-goods-game player types
\citep{fischbacher2001people}: the \emph{altruistic builder},
\emph{conditional cooperator}, \emph{strategic hedger}, and \emph{free rider}.
The state space has $|\mathcal{S}| = 5$ scenario framings (shared resource,
vulnerable target, competitive market, long-term investment, abundance) and the
joint action consists of per-agent contribution choices at $|\mathcal{A}_i| = 3$
levels. The coordinator's joint posterior over agent types has $|\Theta|^n = 64$
profiles; H-PSMG's factored representation occupies $4 \cdot 3 = 12$ entries.

\paragraph{Assumption verification.} The six HT-MG-BC assumptions are
satisfied by construction (TI, RL, PF, ISI, CSR) or verified empirically (TID
via calibration with minimum pairwise Hellinger gap $\geq 0.05$). Calibration
probes the per-(persona, framing, contribution) reward tensor
$R[\theta, s, a]$ via three rollouts per cell and learns the LLM judge's
satisfaction distribution; this becomes the likelihood backbone for the
factored Bayesian posterior updates.

\paragraph{Algorithms.} We evaluate H-PSMG\textsuperscript{+} ($\beta = 0.25$,
held fixed across models) against three families of baselines. \emph{Native
HT-MG-BC variants}: H-PSMG, Joint-PSRL, MAP-Type-Greedy (Bayesian type-aware);
PSRL-NoType, IQL, Random (type-agnostic, the rate-separation baselines of
Theorem~\ref{thm:rate-separation}); and Oracle (upper bound). \emph{External
LLM-coordination methods from recent literature}: A-ToM \citep{xu2026atom} at
ToM orders $\{0, 1, 2\}$ with online aggregation via Hedge \citep{cesa1997use}
and Follow-the-Leader; the Bayesian Nash equilibrium coordinator of
ECON \citep{yang2025econ} run with three deliberation rounds per episode.
\emph{Prompted-LLM coordinators}: \texttt{llm\_belief} (directly prompted to
maintain a belief over personas and choose joint actions) and \texttt{llm\_greedy}
(prompted to play the most-likely-type-conditioned policy without explicit
exploration).

\paragraph{LLM backbones.} The experiment is replicated on four LLM backbones,
chosen to span vendor, capability, and weight access. \emph{DeepSeek-V3.2} is a
strong open-weight mixture-of-experts model. \emph{GPT-5.4-nano-20260317} is
a small closed-weight model that tests robustness to reduced reasoning capacity.
\emph{Kimi-K2.6} is a mid-capacity open-weight model. \emph{Llama-4-Maverick-17B-128E
(FP8)} is an open-weight mid-large model. The calibration tensor $R[\theta, s, a]$
is learned separately for each backbone (denoted c19); H-PSMG\textsuperscript{+}'s
$\beta$ is held fixed at $0.25$ across backbones to preclude per-model tuning.

\paragraph{Cross-model results.} Figure~\ref{fig:hp-spgg-cross-model} reports
cumulative Bayesian regret at $K = 20$ across all four backbones with $5$ seeds,
showing a consistent $12$-algorithm comparison on every backbone. Three regret
clusters are visible on every backbone. The Bayesian type-aware cluster
(H-PSMG\textsuperscript{+}, H-PSMG, MAP-Type-Greedy, Joint-PSRL) clusters tight
near the oracle with cumulative regret below $1.7$. The LLM-coordination cluster
(\texttt{llm\_belief}, ECON-BNE, A-ToM-0, A-ToM hedge) accumulates regret in the
$3$ to $15$ range. Type-agnostic baselines (PSRL-NoType, Random, IQL) accumulate
regret in the $10$ to $30$ range. The H-PSMG family achieves the lowest
non-oracle regret on all four backbones, at $0.40$ on DeepSeek, $0.64$ on
GPT-5.4-nano, $0.63$ on Kimi, and $0.31$ on Llama-Maverick. The separation
between H-PSMG family and best LLM-coordination baseline is $7.7\times$,
$6.3\times$, $11.8\times$, and $10.8\times$ respectively, and the separation
from the best type-agnostic baseline is $27\times$ to $35\times$. The latter
directly instantiates the rate gap predicted by Theorem~\ref{thm:rate-separation}
on real LLM rollouts.

\paragraph{No single LLM-coordination method dominates across backbones.} The
identity of the best LLM-coordination baseline varies across the four LLM
backbones: \texttt{llm\_belief} on DeepSeek-V3.2, A-ToM-0 on GPT-5.4-nano,
A-ToM adaptive-hedge on Kimi-K2.6, and ECON-BNE on Llama-Maverick. ECON-BNE,
which won on DeepSeek and Llama-Maverick with cumulative regret $3.99$ and
$3.38$, catastrophically degrades on GPT-5.4-nano to $14.88$. A-ToM-0, the
best on GPT-5.4-nano at $4.08$, is the worst of the four LLM-coordination
methods on DeepSeek at $6.00$ and only third on Kimi. The cross-model
variation reflects the strong dependence of prompt-based belief reasoning on
the player LLM's multi-step reasoning capacity, and the absence of any
LLM-coordination method that is robust across backbones. In contrast, the
H-PSMG family wins on \emph{all four} backbones because its belief tracking is
closed-form Bayesian and structurally independent of the player LLM's reasoning
capacity.

\begin{figure}[t]
  \centering
  \includegraphics[width=0.95\textwidth]{figs/fig10_hp_spgg_cross_model_v3.pdf}
  \caption{Cumulative Bayesian regret on HP-SPGG across four LLM backbones
  ($K{=}20$, $5$ seeds, calibration c19, mean shown with thin error indicators
  for $\pm\sigma$). Each panel shows the same set of algorithms, grouped into
  three families by background tint: \emph{Bayesian type-aware} (blue band,
  top) including H-PSMG\textsuperscript{+} and H-PSMG (our methods, bolded
  labels), MAP-Type-Greedy, and Joint-PSRL; \emph{LLM coordination} (warm
  band) including $\mathrm{llm\_belief}$, ECON-BNE, A-ToM-0, and A-ToM hedge;
  \emph{type-agnostic} (gray band, bottom) including PSRL-NoType, Random, and
  IQL. Numeric values at bar tips show mean regret. The H-PSMG family attains
  the lowest non-oracle regret on every backbone, separating from the best
  LLM-coordination baseline by $6$ to $12\times$ and from the best type-agnostic
  baseline by over $25\times$. The identity of the best LLM-coordination
  baseline varies across backbones (no single LLM-coordination method dominates),
  while the H-PSMG family is uniformly best.}
  \label{fig:hp-spgg-cross-model}
\end{figure}

\paragraph{Cross-model variation in best-variant selection.} Inspection of the
strict $\beta = 0.25$ values shows that H-PSMG\textsuperscript{+} wins on
DeepSeek and Llama-Maverick (where the bonus accelerates burn-in concentration
of the type posterior), while the $\beta = 0$ ablation H-PSMG wins on GPT-5.4-nano
and Kimi (where the bonus contributes essentially zero finite-$K$ improvement
because the posterior already concentrates rapidly). Across the eight values
of $\beta$ in our sweep, the family $\{\text{H-PSMG}, \text{H-PSMG}^+\}$
includes the lowest-regret variant on all four backbones. We therefore
report the family rather than either single variant.

\paragraph{Regime-dependence of the type-discrimination bonus matches
Theorem~\ref{thm:hpsmg-plus}.} Figure~\ref{fig:beta-sweep} sweeps $\beta$
over the grid $\{0, 0.05, 0.1, 0.25,$ $0.5, 0.75, 1.0, 1.5\}$ on each backbone. Three
qualitatively different regimes are observable. On DeepSeek-V3.2,
H-PSMG\textsuperscript{+}'s regret is exactly invariant at $0.40$ across
$\beta \in [0, 0.75]$ and rises mildly at $\beta = 1.0, 1.5$, reflecting the
self-decay property of Theorem~\ref{thm:hpsmg-plus}(i) plus the high-$\beta$
boundary regime where the bonus begins to pull the algorithm off welfare-optimal
actions. On GPT-5.4-nano and Kimi-K2.6, regret is essentially flat across all
$\beta$, an extreme case of self-decay in which the posterior concentrates
faster than $\beta$ ever has time to add useful exploration. On Llama-Maverick,
the bonus produces a sharp $3 \times$ improvement at $\beta = 0.25$ (regret
drops from $0.97$ to $0.31$) and stays flat thereafter, the empirical signature
of the exponential-in-$n$ burn-in improvement predicted by
Theorem~\ref{thm:hpsmg-plus}(ii). Section~\ref{sec:e1-posterior-concentration}
shows that these three regimes correspond exactly to backbone posterior
concentration speed, supplying a mechanistic test of Theorem~\ref{thm:hpsmg-plus}'s
trigger conditions.

\begin{figure}[t]
  \centering
  \includegraphics[width=0.85\textwidth]{figs/fig10_beta_sweep_v3.pdf}
  \caption{$\beta$-sweep of H-PSMG\textsuperscript{+}'s type-discrimination
  bonus on HP-SPGG c19 across four LLM backbones, each labelled at the curve's
  right endpoint. Three qualitatively distinct regimes appear, each predicted
  by Theorem~\ref{thm:hpsmg-plus}. \emph{Extreme self-decay} on GPT-5.4-nano
  and Kimi-K2.6: regret is essentially flat across all $\beta$ because the
  type posterior concentrates faster than the bonus has time to drive useful
  exploration; this is the boundary case of Theorem~\ref{thm:hpsmg-plus}(i).
  \emph{Self-decay edge} on DeepSeek-V3.2: regret stays at $0.40$ for $\beta
  \leq 0.75$ and rises mildly at higher $\beta$ as the bonus begins to pull
  the algorithm off welfare-optimal actions, matching the high-$\beta$
  boundary regime of Theorem~\ref{thm:hpsmg-plus}(i). \emph{Burn-in jump} on
  Llama-Maverick: regret falls sharply from $0.97$ at $\beta \leq 0.1$ to
  $0.31$ at $\beta = 0.25$ and stays flat thereafter, the empirical
  signature of the exponential-in-$|\Theta|^n$ burn-in improvement of
  Theorem~\ref{thm:hpsmg-plus}(ii). The shaded vertical band marks the
  $\beta = 0.1 \to 0.25$ burn-in transition zone where Llama crosses regimes.}
  \label{fig:beta-sweep}
\end{figure}

\paragraph{Robustness to LLM-capability degradation.} The cross-model spread of
LLM-coordination methods (Figure~\ref{fig:hp-spgg-cross-model},
Table~\ref{tab:hp-spgg-headline}) is highly method-specific.
ECON-BNE catastrophically degrades on GPT-5.4-nano, accumulating $14.88$
cumulative regret versus $3.99$ on DeepSeek and $3.38$ on Llama-Maverick, a
$3.7\times$ degradation on the same task with only the player LLM changed.
Trace inspection reveals the failure mode: GPT-5.4-nano has difficulty
following the multi-step belief-deliberation prompts that ECON's
Bayesian-Nash-equilibrium computation requires, producing inconsistent belief
updates. By contrast, A-ToM-0 improves on GPT-5.4-nano ($4.08$ versus $6.00$
on DeepSeek) because it requires no multi-step belief reasoning; A-ToM hedge
is stable across all four backbones in the $4.79$ to $7.65$ range. The H-PSMG
family, whose belief tracking is closed-form Bayesian and structurally
independent of the player LLM's reasoning capacity, is robust across all four
backbones, with regret in the $0.31$ to $0.91$ range on every backbone. This
robustness profile is a structural property of explicit posterior tracking,
present whenever the player LLM lacks the multi-step reasoning capability that
prompt-based belief methods rely upon.

\begin{table}[t]
\centering
\caption{HP-SPGG cumulative Bayesian regret at $K = 20$, $5$ seeds,
mean $\pm$ std. H-PSMG family (best of $\{\text{H-PSMG}, \text{H-PSMG}^+\}$
across $\beta \in [0, 1.5]$) is the lowest non-oracle entry on every backbone,
separating from the best LLM-coordination baseline by $6.3$ to $11.8\times$
and from the best type-agnostic baseline by $27$ to $35\times$. The identity
of the best LLM-coordination baseline varies across backbones: \texttt{llm\_belief}
on DeepSeek, A-ToM-0 on GPT-5.4-nano, A-ToM hedge on Kimi-K2.6, and ECON-BNE
on Llama-Maverick. ECON-BNE's regret on GPT-5.4-nano ($14.88$) is $3.7\times$
its regret on DeepSeek ($3.99$), reflecting GPT-5.4-nano's reduced capacity to
execute ECON's multi-step belief-deliberation prompts.}
\label{tab:hp-spgg-headline}
\small
\begin{tabular}{lrrrr}
\toprule
Algorithm & DeepSeek-V3.2 & GPT-5.4-nano & Kimi-K2.6 & Llama-Maverick \\
\midrule
Oracle                         & $0.000$               & $0.000$               & $0.000$               & $0.000$ \\
\textbf{H-PSMG family (ours)}  & $\mathbf{0.400 \pm 0.80}$ & $\mathbf{0.644 \pm 0.63}$ & $\mathbf{0.632 \pm 0.61}$ & $\mathbf{0.312 \pm 0.62}$ \\
MAP-Type-Greedy                & $0.712 \pm 0.47$      & $0.854 \pm 0.42$      & $0.762 \pm 0.39$      & $0.974 \pm 0.44$ \\
Joint-PSRL                     & $0.832 \pm 0.08$      & $1.492 \pm 0.34$      & $1.650 \pm 0.40$      & $1.194 \pm 0.35$ \\
\midrule
\texttt{llm\_belief} (prompted)         & $3.074 \pm 2.23$  & $7.197 \pm 3.47$  & $11.78 \pm 3.33$  & $10.36 \pm 3.82$ \\
ECON-BNE \citep{yang2025econ}           & $3.990 \pm 3.41$  & $14.88 \pm 2.75$  & $10.49 \pm 6.03$  & $3.382 \pm 4.53$ \\
A-ToM-0 \citep{xu2026atom}              & $6.000 \pm 7.46$  & $4.080 \pm 4.94$  & $10.67 \pm 6.36$  & $7.240 \pm 7.79$ \\
A-ToM hedge \citep{xu2026atom}          & $4.790 \pm 4.51$  & $7.653 \pm 4.96$  & $7.484 \pm 4.10$  & $7.006 \pm 4.58$ \\
\midrule
PSRL-NoType                    & $13.91 \pm 7.37$ & $17.72 \pm 1.68$ & $19.51 \pm 2.61$ & $10.65 \pm 4.36$ \\
Random                         & $28.27 \pm 5.12$ & $26.46 \pm 1.81$ & $26.54 \pm 4.51$ & $24.34 \pm 4.32$ \\
IQL                            & $28.34 \pm 16.9$ & $24.79 \pm 11.7$ & $20.12 \pm 18.5$ & $27.53 \pm 16.3$ \\
\bottomrule
\end{tabular}
\end{table}

\subsubsection{Theory-driven experiments: mechanistic tests of Theorem~\ref{thm:hpsmg-plus}}
\label{sec:theory-driven-experiments}

The $\beta$-sweep narrative above identified three regimes empirically. We now
run three further experiments on HP-SPGG that probe the mechanistic
predictions of Theorem~\ref{thm:hpsmg-plus} directly. E1 measures the
theorem's trigger condition (per-backbone posterior concentration speed) and
checks that it predicts the regime. E2 verifies the predicted scaling of the
bonus benefit with type-space size $|\Theta_i|$. E3 verifies the predicted
scaling with agent count $n$ together with the analytic storage advantage of
posterior decoupling.

\paragraph{E1. Posterior concentration speed predicts the regime.}
\label{sec:e1-posterior-concentration}
Theorem~\ref{thm:hpsmg-plus} predicts that the type-discrimination bonus
contributes nothing once the posterior on the true type has concentrated, and
provides exponential burn-in improvement only when concentration is slow
relative to the planning horizon. We turn this into a direct measurement.
On every backbone we run vanilla H-PSMG (i.e., $\beta = 0$, so the
measurement is independent of the bonus mechanism) and log the per-episode
posterior mass on the true type, $\mu_k^{\Theta, i}(\theta_i^\star)$, then
define the per-seed concentration time
$K_{\mathrm{conc}}^{0.9} = \min\{k : \mu_k^{\Theta, i}(\theta_i^\star) \geq 0.9
\text{ for all } i\}$ and average over $5$ seeds. The four backbones span a
$3\times$ range: DeepSeek-V3.2 concentrates fastest at
$K_{\mathrm{conc}}^{0.9} = 4.4$, GPT-5.4-nano at $6.6$, Kimi-K2.6 at $12.0$,
and Llama-Maverick at $14.0$ (out of $K = 20$ planning horizon). Figure~\ref{fig:e1-posterior-concentration}
plots concentration time against the empirically observed bonus benefit
$\mathrm{Reg}(\text{H-PSMG}) - \mathrm{Reg}(\text{H-PSMG}^+)$ at $\beta = 0.25$
and reveals a clear regime structure. The fastest-concentrating backbone
(DeepSeek, $K_{\mathrm{conc}}^{0.9} = 4.4$) sits in the self-decay edge regime
where the bonus self-decays before episode $5$ but produces a gentle benefit
of $+0.43$. The slowest-concentrating backbone (Llama, $K_{\mathrm{conc}}^{0.9} = 14.0$)
sits in the burn-in jump regime predicted by Theorem~\ref{thm:hpsmg-plus}(ii)
where the bonus drives a benefit of $+0.42$. The two middle backbones
(GPT and Kimi, $K_{\mathrm{conc}}^{0.9} \in [6.6, 12.0]$) sit in a boundary
regime where the bonus over-explores in the burn-in phase and yields a small
negative benefit; this is the high-$\beta$ edge of Theorem~\ref{thm:hpsmg-plus}(i)
manifesting at finite $\beta$ when the posterior is neither fast nor slow.
The mechanism that the theorem identifies, posterior concentration speed, is
the right backbone-level predictor of the bonus's behaviour.

\begin{figure}[t]
  \centering
  \includegraphics[width=0.85\textwidth]{figs/fig_e1_posterior_concentration_v3.pdf}
  \caption{E1 posterior concentration speed predicts the H-PSMG\textsuperscript{+}
  bonus regime. Each marker is one LLM backbone. The horizontal axis is the mean
  per-seed concentration time $K_{\mathrm{conc}}^{0.9}$, the first episode at
  which the H-PSMG ($\beta = 0$) posterior on the true type for all agents
  exceeds $0.9$. The vertical axis is the empirically observed bonus benefit,
  $\mathrm{Reg}(\text{H-PSMG}) - \mathrm{Reg}(\text{H-PSMG}^+)$ at $\beta = 0.25$,
  positive when the bonus helps. The three background bands correspond to the
  three regimes of Theorem~\ref{thm:hpsmg-plus}. DeepSeek-V3.2 sits in the
  self-decay edge regime (the bonus self-decays before useful exploration but
  still produces a gentle positive benefit); Llama-4-Maverick sits in the
  burn-in jump regime predicted by Theorem~\ref{thm:hpsmg-plus}(ii); GPT-5.4-nano
  and Kimi-K2.6 sit in an intermediate boundary regime where the bonus
  over-explores at finite $\beta$.}
  \label{fig:e1-posterior-concentration}
\end{figure}

\paragraph{E2. The bonus closes a burn-in gap that grows with $|\Theta_i|$.}
\label{sec:e2-type-scaling}
Theorem~\ref{thm:hpsmg-plus}(ii) makes two simultaneous predictions about
vanilla H-PSMG and H-PSMG\textsuperscript{+} on the deep-trap family. Vanilla
H-PSMG, despite being asymptotically Bayesian-optimal, incurs burn-in regret
scaling as $\Omega(|\Theta_i|^n H)$. H-PSMG\textsuperscript{+} caps the burn-in
regret at $O(H)$ independently of $|\Theta_i|$. E2 tests both predictions
simultaneously by varying the persona count $|\Theta_i| \in \{2, 3, 4, 5, 6\}$
on HP-SPGG with $n = 3$ agents, fixing the backbone to Llama-Maverick (the
only backbone in the burn-in jump regime per E1, so the trigger condition for
the predicted separation holds), and running $10$ seeds per configuration.
Figure~\ref{fig:e2-type-scaling} reports cumulative regret at $K = 20$. The
two curves are not competing algorithms but the two theoretical regimes
predicted by Theorem~\ref{thm:hpsmg-plus}(ii) made empirically observable: the
upper curve (vanilla H-PSMG) traces the predicted $\Omega(|\Theta_i|^n)$
burn-in cost, the lower curve (H-PSMG\textsuperscript{+}) traces the predicted
$O(1)$ cap, and the shaded band measures the empirical realisation of the
predicted gap. The gap grows from $+0.14$ at the smallest type space
($|\Theta_i| = 2$, where $|\Theta_i|^n = 8$ for $n = 3$ and the burn-in budget
is small) to $+1.08$ at the largest ($|\Theta_i| = 6$, where $|\Theta_i|^n = 216$
and the burn-in budget is order-of-magnitude larger). The trajectory is
non-monotonic in the middle ($|\Theta_i| \in \{4, 5\}$ show smaller and
noisier gaps than $|\Theta_i| = 3$), reflecting persona-design specifics at
each $|\Theta_i|$ together with the finite-seed variance characteristic of
LLM substrates. The envelope, however, recovers the predicted scaling: the
burn-in regime is empirically present, and the bonus mechanism empirically
closes it.

\begin{figure}[t]
  \centering
  \includegraphics[width=0.85\textwidth]{figs/fig_e2_type_scaling_v3.pdf}
  \caption{E2 empirical realisation of Theorem~\ref{thm:hpsmg-plus}(ii)'s
  two-regime prediction on HP-SPGG (Llama-Maverick backbone, $n = 3$ agents,
  $10$ seeds, mean $\pm$ SEM). The two curves are not competing algorithms but
  the two regimes the theorem predicts. \emph{Upper curve} (light blue): vanilla
  H-PSMG, which the theorem predicts to incur $\Omega(|\Theta_i|^n H)$ burn-in
  regret. \emph{Lower curve} (deep navy): H-PSMG\textsuperscript{+} (ours,
  $\beta = 0.25$), which the theorem predicts to cap burn-in regret at $O(H)$.
  The shaded band measures the predicted gap; per-axis labels give its magnitude.
  The gap grows from $+0.14$ at $|\Theta_i| = 2$ to $+1.08$ at $|\Theta_i| = 6$,
  the predicted regime where the burn-in budget is largest and the bonus is
  most valuable. Mid-range non-monotonicity reflects persona-design specifics
  together with finite-seed variance.}
  \label{fig:e2-type-scaling}
\end{figure}

\subsection{Concordia Public Benchmarks: Pub Coordination and Haggling}
\label{sec:concordia}

To rule out the possibility that the HP-SPGG results above reflect bias toward
a benchmark of our own construction, we evaluate H-PSMG\textsuperscript{+} on
two substrates from the public Concordia framework
\citep{vezhnevets2023concordia} which was independently designed by DeepMind
for the NeurIPS 2024 Cooperative AI Contest \citep{concordia_contest2025}.

\paragraph{Substrate mapping.} For both substrates, we map the latent agent
disposition or loyalty to $\theta_i$, the substrate's exogenous state
(match schedule or merchant offer) to $s$, and per-agent decisions (pub choice
or offer acceptance) to $a_i$. Concordia substrates provide an analytic
per-agent reward function via the substrate's coordination payoff component,
eliminating the need for an LLM judge and halving per-episode cost. We
implement the centralised-coordinator variant of each substrate (CSR holds
by construction) and verify the remaining HT-MG-BC assumptions by inspection
of the substrate code.

\paragraph{Algorithm variants.} On Concordia we run two variants of our
algorithm. \emph{H-PSMG\textsuperscript{+} (joint proxy)} solves the centralised
joint action problem using the factored posterior, mirroring the algorithm of
Section~\ref{sec:algorithm}. \emph{H-PSMG\textsuperscript{+} (proxy)} commits to
the per-agent MAP type at each round and runs a per-agent best response, a
variant that exposes the posterior decoupling theorem's value in narrow
substrates. We also include two strong external LLM-coordination baselines
adapted to the Concordia interface: \emph{ECON-BNE-mech} runs ECON's
Bayesian-Nash-equilibrium coordinator with mechanistic policy execution
\citep{yang2025econ}, and \emph{A-ToM-1-mech} runs A-ToM at first-order theory
of mind with mechanistic policy execution \citep{xu2026atom}. The substrate's
joint-optimal policy is reported as \emph{Oracle (joint)}, which is the
mechanistic upper bound on joint score but not necessarily on focal score.

\begin{figure}[t]
  \centering
  \includegraphics[width=\textwidth]{figs/fig10_concordia_main_v3.pdf}
  \caption{Concordia results across both substrates. \emph{Left:} Pub Coordination,
  $9$ axes covering $7$ village configurations at $s = 30$ (one of which,
  \texttt{capetown}, is also evaluated at $s = 100$) plus one DeepSeek-V3.2 live
  LLM rollout on \texttt{london\_mini} at $s = 5$. \emph{Right:} Haggling, $9$
  axes spanning $5$ single-item and $4$ multi-item configurations at $s = 30$.
  Each axis is per-axis min--max normalised so that the best of the four
  algorithms on that axis reaches the outer ring and the worst sits near the
  centre; oracle is shown as a dashed reference polygon and is not always the
  outer ring because H-PSMG\textsuperscript{+} can strictly exceed oracle on
  focal score (see \texttt{vegbrooke\_stubborn} and the \texttt{fruitville\_gullible}
  axes on Haggling). H-PSMG\textsuperscript{+} (joint proxy) fills nearly the
  full Pub Coordination polygon and ties or wins on all $9$ Haggling axes;
  ECON-BNE-mech and A-ToM-1-mech contract toward the centre on most non-trivial
  configurations; on \texttt{vegbrooke\_strange}, A-ToM-1-mech catastrophically
  collapses to the inner ring at focal score $-10$.}
  \label{fig:concordia-main}
\end{figure}

\subsubsection{Concordia Pub Coordination}

The Pub Coordination substrate populates a small village with $n$ residents
who must select a pub each evening; each resident has a hidden team loyalty
$\theta_i \in \{$Team A, B, C, D$\}$, and the pub is satisfying for resident
$i$ if (a) the substrate's day-state schedules resident $i$'s team to be
playing at the chosen pub, and (b) sufficiently many fellow fans of resident
$i$'s team are also present.

Figure~\ref{fig:concordia-main} (left) reports focal-score means across $7$
substrate configurations spanning four village variants (\texttt{capetown},
\texttt{edinburgh}, \texttt{london}, \texttt{london\_mini}) and three closure
regimes (no closures, closures, tough friendship). The mechanistic runs at
$s = 30$ (and a longer $s = 100$ run on \texttt{capetown}) form the
headline; an additional DeepSeek-V3.2 live LLM rollout at $s = 5$ on
\texttt{london\_mini} provides a preliminary check that the mechanistic
ordering survives when the LLM drives the rollouts in-loop.

H-PSMG\textsuperscript{+} joint proxy is within $0.03$ (about $1\%$) of
\texttt{oracle\_joint} on every mechanistic config. ECON-BNE-mech trails by
$5$ to $15\%$, A-ToM-1-mech by $5$ to $30\%$. The qualitative pattern survives
the live LLM rollout: on \texttt{london\_mini} with DeepSeek-V3.2 at $s = 5$
the proxy variant trails oracle by $0.19$, which we attribute to small-sample
noise (the corresponding $s = 30$ mechanistic run on the same substrate
variant has gap $0.008$).

\subsubsection{Concordia Haggling}

The Haggling substrate places a buyer and a seller in a multi-round negotiation
over one or more items; the seller has a hidden tolerance for low offers
($\theta_i$) and the buyer's reward is a function of the agreed price minus
the item's intrinsic value to the buyer. We run $9$ sub-configurations spanning
$5$ substrate variants (fruitville variants and vegbrooke variants) and single
versus multi-item versions.

Figure~\ref{fig:concordia-main} (right) reports focal-score means across the $9$
configurations. H-PSMG\textsuperscript{+} wins or ties on all configs. Two
observations deserve explicit discussion. First, on
\texttt{vegbrooke\_stubborn} and the two \texttt{fruitville\_gullible}
configurations, H-PSMG\textsuperscript{+} \emph{strictly exceeds} \texttt{oracle\_joint}
(by $1.43$, $0.40$, and $0.47$ respectively). This is not a contradiction.
\texttt{oracle\_joint} maximises the substrate's joint payoff, which on asymmetric
substrates can sacrifice focal-agent payoff to maintain coordination with the
opponent. H-PSMG\textsuperscript{+} maximises focal-agent payoff directly and
therefore strictly improves on joint-payoff maximisation when the two diverge.
Second, on \texttt{vegbrooke\_strange}, A-ToM-1-mech catastrophically fails at
$-10.0$ (the substrate's minimum-payoff lower bound), while
H-PSMG\textsuperscript{+}, ECON-BNE-mech, and Oracle all hold at $0$ (the
no-agreement value). The failure mode is brittleness of A-ToM's
expert-aggregation under the substrate's adversarial pricing schedule;
H-PSMG\textsuperscript{+}'s closed-form posterior is robust to this regime.

\paragraph{Summary across both Concordia substrates.} H-PSMG\textsuperscript{+}
achieves the highest non-oracle focal score on $16 / 16$ Concordia
configurations ($7$ Pub Coord, $9$ Haggling), with mean gap to oracle of
$1.1\%$ on Pub Coord and a strict exceeding-of-oracle on $3 / 9$ Haggling
configs. ECON-BNE-mech trails by $5$ to $15\%$ on Pub Coord and is competitive
but never strictly better on Haggling. A-ToM-1-mech trails by $5$ to $30\%$ on
Pub Coord and catastrophically fails on $1$ Haggling config. The
cross-substrate consistency demonstrates that the algorithmic advantage of
posterior decoupling is intrinsic to settings satisfying HT-MG-BC structural
assumptions, not an artefact of the HP-SPGG benchmark we constructed.

\subsection{Discussion}

The HP-SPGG and Concordia experiments together establish that the
HT-MG-BC framework's predictions extend beyond controlled simulation to LLM-MAS
benchmarks at three levels of independence from our experimental design. On
HP-SPGG, the benchmark we constructed to make all assumptions verifiable, the
H-PSMG family dominates published LLM-coordination methods (ECON, A-ToM) and
prompted-LLM coordination by $6.3$ to $11.8\times$ across four LLM backbones.
On Concordia, an independently-designed public substrate, the same family
achieves within $1\%$ of oracle on Pub Coordination ($7/7$ configs) and wins
or ties on Haggling ($9/9$ configs). Cross-model $\beta$-sweep results validate
all three regimes predicted by Theorem~\ref{thm:hpsmg-plus} on different LLM
backbones, including the burn-in improvement of $(ii)$ on Llama-Maverick. The
sufficiency of structural Bayesian posterior tracking over learned
belief-network or expert-advice approaches is supported across all sixteen
Concordia configurations.


\section{Discussion and Extensions}
\label{sec:discussion}

\subsection{Centralized vs. Decentralized}

The decoupling result holds for centralized H-PSMG, where a single learner manages decisions for all agents. In decentralized HT-MG, each agent $i$ would maintain its own posterior $\mu^{(i),P}, \mu^{(i),\Theta,j}$ for $j \neq i$ and best-respond to its own samples. The likelihood agent $i$ assigns to opponent action $a_j$ would be the expectation of agent $j$'s policy over agent $j$'s posterior beliefs, leading to interactive Bayesian reasoning of the type studied in the I-POMDP literature \citep{gmytrasiewicz2005framework}. Decoupling no longer holds in general for decentralized settings without further structural assumptions (e.g., common prior, level-$k$ bounded reasoning). Extending the present framework to decentralized HT-MG is a natural follow-up.

\subsection{Beyond Type-Independent Transitions}

The TI assumption is restrictive in settings where the agent type affects the environment, for example through differing physical capabilities. Two ways to extend address this. First, one can introduce a joint kernel $P(s' \mid s, a, \boldsymbol{\theta})$ and maintain a joint posterior with $|\Theta|$ separate Dirichlet posteriors, at storage cost $O(|\Theta|)$. Second, one can decompose $P = P_0 + \sum_i \phi(\theta_i) \delta P_i$ with low-rank structure on the type-dependent perturbation, enabling a hierarchical model with milder cost.

\subsection{Computational Tractability of CCE Oracle}

Computing a CCE in general-sum MGs is polynomial in the size of the joint action space $A = \prod_i |\mathcal{A}_i|$, which grows exponentially in $n$. For large-$n$ settings, structured equilibrium concepts such as Markov Coarse Correlated Equilibrium \citep{daskalakis2023complexity} or decentralized V-learning style updates \citep{jin2021v} can be plugged into Algorithm~\ref{alg:hpsmg} as the oracle without affecting the decoupling result.

\subsection{Function-Approximation Generalisation}
\label{sec:fa-extension}

A natural concern is whether the tabular analysis of this paper survives the move to large or continuous state-action spaces, where transitions and rewards must be represented by parametric function classes (linear models, kernel methods, neural networks). Appendix~\ref{app:fa-hpsmg} develops the function-approximation generalisation, FA-H-PSMG\textsuperscript{+} (Algorithm~\ref{alg:fa-hpsmg-plus}), and establishes three results that together unify the tabular and FA regimes.

First, the posterior decoupling theorem (Theorem~\ref{thm:fa-decoupling}) holds at the level of the prior and the model class: under FA-analogues of TID and OTRL (Assumptions~\ref{ass:tid-fa},~\ref{ass:otrl-fa}), the joint posterior over $(\phi, \boldsymbol{\psi}, \boldsymbol{\theta})$ factorises across the transition parameters, the per-agent reward parameters, and the per-agent type. The argument is identical in structure to Theorem~\ref{thm:decoupling}; what changes is the representation of each block (Gaussian on neural features, ensemble of networks, etc.), not its independence. Storage remains $O(\mathrm{rep}(\mu^P) + n \cdot \mathrm{rep}(\mu^{\Psi, i, \Theta, i}))$, linear in $n$.

Second, the type-discrimination bonus of H-PSMG\textsuperscript{+} is computable in the FA setting via finite-sample Monte Carlo: $\widehat{D}_k(s, a) = \frac{1}{M} \sum_j |r_i^{\psi}(s, a, \theta_i^{(j)}) - r_i^{\psi}(s, a, \theta_i^{(j')})|$ with $\theta_i^{(j)}, \theta_i^{(j')} \sim \mu_k^{\Theta, i}$. For $M = O(\log K)$ the estimation error is absorbed into the leading regret term, and the deterministic-escape conclusion of Proposition~\ref{prop:hpsmg-plus-burnin} extends to FA-H-PSMG\textsuperscript{+} (Proposition~\ref{prop:fa-burnin}).

Third, the regret bound interpolates smoothly between the tabular and FA regimes (Theorem~\ref{thm:fa-regret-informal}): the $\sqrt{S A_{\mathrm{joint}} K}$ tabular complexity is replaced by the eluder dimensions $\mathrm{dim}_E(\mathcal{P}), \mathrm{dim}_E(\mathcal{R}_i)$ of the transition and reward function classes \citep{russo2014learning}, while the type-uncertainty term $\sqrt{|\Theta_i| K}$ is unchanged. Corollary~\ref{cor:fa-regret-tabular} recovers Theorem~\ref{thm:main-regret-informal} as the tabular special case. The exponential-in-$n$ improvement over generic Bayesian-RL is preserved because the type space is finite regardless of the choice of function class for the dynamics.

Empirical validation of FA-H-PSMG\textsuperscript{+} on continuous-state HT-MGs (e.g., particle-based cooperation environments) is a natural next step. The present paper restricts to the tabular case so the exponential-in-$|\Theta|^n$ separation is visible without confounding from function-approximation error.

\subsection{Open Problems}

Several questions remain. First, what is the precise minimax lower bound on type-conditional Bayesian regret in HT-MG under TID and OTRL, beyond the matching $\Omega(\sqrt{|\Theta_i| HK})$ of Theorem~\ref{thm:lower-matching}? Existing latent MDP lower bounds \citep{kwon2021rl} are exponential in $n$ without identifiability; under bounded information rate, what is the right rate? Second, for decentralised HT-MG (Section~\ref{sec:discussion}), under what additional assumptions does a relaxed decoupling result still hold? Third, in social dilemmas where own-type reward locality (OTRL) fails by design, can a controlled-error decoupling result be recovered, perhaps via a $\ell_\infty$ approximation factor in the marginalisation step?

\section{Conclusion}
\label{sec:conclusion}

We have introduced the Heterogeneous-Type Markov Game framework with Bayesian Coordination, proposed the centralized H-PSMG algorithm, and proved that under two natural structural assumptions the joint posterior over the transition kernel and the type profile decouples exactly at every episode. The decoupling reduces storage cost from exponential to linear in the number of agents and enables an additive regret decomposition into a model-learning component and a type-inference component. We have verified the result numerically on a two-agent two-type matrix game, established necessity of both assumptions via explicit counter-examples, and laid out the roadmap from decoupling to a full regret bound. To our knowledge this is the first posterior-sampling result for general-sum Markov games with latent agent types and rigorous structural guarantees.

\bibliography{ref}
\bibliographystyle{iclr2023_conference}

\appendix

\section{Rigorous Treatment of the \texorpdfstring{$n$}{n}-Player Case}
\label{app:n-player}

Lemma~\ref{lem:likelihood} and Theorem~\ref{thm:decoupling} are stated in the main text for $n$ agents. This appendix makes the $n$-agent factorization explicit, provides the full induction, derives the storage and computational complexity, and works through a concrete $n=3$ instance.

\subsection{Explicit Likelihood Factorization for \texorpdfstring{$n$}{n} Agents}

Fix episode $k$. Given a policy $\pi_k$ and a candidate world $(P, \boldsymbol{\theta}) = (P, \theta_1, \ldots, \theta_n)$, the trajectory likelihood unrolls as
\begin{align}
P(\tau_k \mid \pi_k, P, \boldsymbol{\theta}) &= \prod_{h=1}^H \pi_k(a_h^k \mid s_h^k) \cdot P(s_{h+1}^k \mid s_h^k, a_h^k, \boldsymbol{\theta}) \cdot \prod_{i=1}^n P(r_h^{k,i} \mid s_h^k, a_h^k, \boldsymbol{\theta}). \label{eq:n-lik-raw}
\end{align}

Apply Assumption~\ref{ass:TI} to the transition factor and Assumption~\ref{ass:RL} to each reward factor:
\begin{equation}
P(s_{h+1}^k \mid s_h^k, a_h^k, \boldsymbol{\theta}) = P(s_{h+1}^k \mid s_h^k, a_h^k), \quad P(r_h^{k,i} \mid s_h^k, a_h^k, \boldsymbol{\theta}) = P(r_h^{k,i} \mid s_h^k, a_h^k, \theta_i).
\end{equation}

Substituting and reordering the products:
\begin{equation}
P(\tau_k \mid \pi_k, P, \boldsymbol{\theta}) = \underbrace{\prod_{h=1}^H \pi_k(a_h^k \mid s_h^k)}_{\text{policy factor}} \cdot \underbrace{\prod_{h=1}^H P(s_{h+1}^k \mid s_h^k, a_h^k)}_{L_P(\tau_k; P)} \cdot \underbrace{\prod_{i=1}^n \left[\prod_{h=1}^H P(r_h^{k,i} \mid s_h^k, a_h^k, \theta_i)\right]}_{\prod_i L_{\Theta,i}(\tau_k; \theta_i)}.
\end{equation}

The reward product splits cleanly into $n$ agent-indexed factors because under RL each summand of the inner product involves a single $\theta_i$. Marginalizing $\pi_k$ over its history-measurable distribution yields
\begin{equation}
\label{eq:n-lik-factor}
P(\tau_k \mid P, \boldsymbol{\theta}, \mathcal{H}_{k-1}) = c_k(\tau_k) \cdot L_P(\tau_k; P) \cdot \prod_{i=1}^n L_{\Theta,i}(\tau_k; \theta_i),
\end{equation}
where $c_k(\tau_k) := \mathbb{E}_{\pi_k \mid \mathcal{H}_{k-1}}\big[\prod_h \pi_k(a_h^k \mid s_h^k)\big]$ depends only on $(\tau_k, \mathcal{H}_{k-1})$ by the argument in Appendix~\ref{app:ck}.

\subsection{Full Induction Proof}

\begin{theorem}[$n$-Player Posterior Decoupling, Formal]
\label{thm:n-player}
Suppose Assumptions~\ref{ass:TI}, \ref{ass:RL} hold and the prior factorizes as
\begin{equation}
P_0(P, \boldsymbol{\theta}) = P_0^P(P) \prod_{i=1}^n P_0^{\Theta,i}(\theta_i).
\end{equation}
Then for every $k \in \{0, 1, \ldots, K\}$,
\begin{equation}
P(P, \boldsymbol{\theta} \mid \mathcal{H}_k) = \mu_{k+1}^P(P) \cdot \prod_{i=1}^n \mu_{k+1}^{\Theta,i}(\theta_i),
\end{equation}
where the marginals satisfy the recursive updates
\begin{equation}
\mu_{k+1}^P(P) = \frac{\mu_k^P(P) \cdot L_P(\tau_k; P)}{Z_k^P}, \qquad \mu_{k+1}^{\Theta,i}(\theta_i) = \frac{\mu_k^{\Theta,i}(\theta_i) \cdot L_{\Theta,i}(\tau_k; \theta_i)}{Z_k^{\Theta,i}}
\end{equation}
with $Z_k^P$ and $Z_k^{\Theta,i}$ the corresponding normalizing constants, and $\mu_1^P = P_0^P$, $\mu_1^{\Theta,i} = P_0^{\Theta,i}$.
\end{theorem}

\begin{proof}
We proceed by induction on $k$.

\textit{Base case} ($k=0$). The prior factorization assumption gives $P(P, \boldsymbol{\theta} \mid \mathcal{H}_0) = P_0^P(P) \prod_{i=1}^n P_0^{\Theta,i}(\theta_i)$, of the required form.

\textit{Inductive step.} Assume $P(P, \boldsymbol{\theta} \mid \mathcal{H}_{k-1}) = \mu_k^P(P) \prod_{i=1}^n \mu_k^{\Theta,i}(\theta_i)$. Applying Bayes' rule with the likelihood factorization \eqref{eq:n-lik-factor},
\begin{align}
P(P, \boldsymbol{\theta} \mid \mathcal{H}_k) &\propto P(P, \boldsymbol{\theta} \mid \mathcal{H}_{k-1}) \cdot P(\tau_k \mid P, \boldsymbol{\theta}, \mathcal{H}_{k-1}) \\
&= \mu_k^P(P) \prod_{i=1}^n \mu_k^{\Theta,i}(\theta_i) \cdot c_k(\tau_k) \cdot L_P(\tau_k; P) \cdot \prod_{i=1}^n L_{\Theta,i}(\tau_k; \theta_i) \\
&\propto \left[\mu_k^P(P) \cdot L_P(\tau_k; P)\right] \cdot \prod_{i=1}^n \left[\mu_k^{\Theta,i}(\theta_i) \cdot L_{\Theta,i}(\tau_k; \theta_i)\right].
\end{align}
The factor $c_k(\tau_k)$ has been absorbed into the normalization constant. The right side is a product of one $P$-dependent term and $n$ separate $\theta_i$-dependent terms. Normalizing each bracket independently and matching with the stated update rule completes the inductive step.
\end{proof}

\subsection{Storage and Computational Complexity}

The factorization in Theorem~\ref{thm:n-player} reduces both storage and update cost by an exponential factor in $n$. We summarize the comparison in Table~\ref{tab:complexity}.

\begin{table}[h]
\caption{Storage and per-episode computational cost of the joint posterior over $(P, \boldsymbol{\theta})$ in $n$-player H-PSMG. We write $|\Theta|_{\max} := \max_i |\Theta_i|$ and $A_{\max} := \max_i |\mathcal{A}_i|$. The transition posterior $\mu^P$ is a product of $SA$ Dirichlet distributions, each of dimension $S$.}
\label{tab:complexity}
\begin{center}
\begin{tabular}{lll}
\toprule
Quantity & Naive joint posterior & Factored H-PSMG \\
\midrule
Type-posterior storage & $\Theta\big(|\Theta|_{\max}^n\big)$ & $\Theta\big(n |\Theta|_{\max}\big)$ \\
Transition-posterior storage & $\Theta\big(S^2 A_{\max}^n \cdot |\Theta|_{\max}^n\big)$ & $\Theta\big(S^2 A_{\max}^n\big)$ \\
Per-episode type-update cost & $\Theta\big(H |\Theta|_{\max}^n\big)$ & $\Theta\big(n H |\Theta|_{\max}\big)$ \\
Per-episode transition-update cost & $\Theta\big(H S A_{\max}^n |\Theta|_{\max}^n\big)$ & $\Theta\big(H S A_{\max}^n\big)$ \\
\bottomrule
\end{tabular}
\end{center}
\end{table}

\begin{corollary}[Tractable Inference at Large \texorpdfstring{$n$}{n}]
\label{cor:scale}
The factored posterior representation makes Bayesian inference in H-PSMG tractable for any number of agents, in the sense that the cost of maintaining and updating $\mu^{\Theta,i}$ grows linearly rather than exponentially in $n$. The transition posterior $\mu^P$ retains the same complexity as in single-agent PSRL with joint action space $\prod_i \mathcal{A}_i$.
\end{corollary}

\begin{remark}
The factored structure is essential for any practical implementation when $n \geq 3$. With $|\Theta_i| = 4$ types per agent and $n = 8$ agents, a joint type posterior over $\Theta$ has $4^8 = 65{,}536$ entries, while the factored representation uses only $32$ scalar parameters. The decoupling theorem certifies that this compression is exact, not an approximation.
\end{remark}

\subsection{Worked Example: Three-Agent Heterogeneous Coordination}

To illustrate the $n$-fold factorization concretely, we extend the prisoner's dilemma of Section~\ref{sec:example} to $n=3$ agents in a coordination game.

\textit{Setup.} Three agents each choose $a_i \in \{C, D\}$. Each agent has a type $\theta_i \in \{P, S\}$. Rewards:
\begin{equation}
r_i^P(a_i, a_{-i}) = \mathbbm{1}\big[\text{all agents play } C\big], \quad r_i^S(a_i, a_{-i}) = \mathbbm{1}[a_i = D].
\end{equation}
The true type profile is $\boldsymbol{\theta}^* = (P, P, S)$. Priors are uniform: $\mu_1^{\Theta,i}(P) = \mu_1^{\Theta,i}(S) = 0.5$ for $i \in \{1,2,3\}$.

\textit{Episode 1.} Suppose the sampled type profile is $\hat{\boldsymbol{\theta}}_1 = (P, P, P)$. The CCE prescribes $(C, C, C)$. Executing this in the true world yields $r_1 = r_1^P(C,C,C) = 1$, $r_2 = r_2^P(C,C,C) = 1$, $r_3 = r_3^S(C, \cdot) = 0$.

\textit{Likelihoods.} For each agent,
\begin{equation}
L_{\Theta,1}(\tau_1; \theta_1) = \begin{cases} 1 & \theta_1 = P \\ 0 & \theta_1 = S \end{cases}, \quad
L_{\Theta,2}(\tau_1; \theta_2) = \begin{cases} 1 & \theta_2 = P \\ 0 & \theta_2 = S \end{cases},
\end{equation}
\begin{equation}
L_{\Theta,3}(\tau_1; \theta_3) = \begin{cases} 0 & \theta_3 = P \\ 1 & \theta_3 = S \end{cases}.
\end{equation}

\textit{Joint posterior.} The full joint posterior over $\Theta = \{P,S\}^3$ has $2^3 = 8$ entries. The unnormalized values are $\mu_1^\Theta(\boldsymbol{\theta}) \cdot \prod_i L_{\Theta,i}(\tau_1; \theta_i)$. Only $\boldsymbol{\theta} = (P,P,S)$ gives nonzero mass $0.125 \cdot 1 \cdot 1 \cdot 1 = 0.125$. Normalizing yields $\mu_2^\Theta(P,P,S) = 1$, true type perfectly identified.

\textit{Factored marginals.} Computing each $\mu_2^{\Theta,i}$ independently:
\begin{equation}
\mu_2^{\Theta,1}(P) = 1, \quad \mu_2^{\Theta,2}(P) = 1, \quad \mu_2^{\Theta,3}(S) = 1.
\end{equation}
The product $\mu_2^{\Theta,1} \cdot \mu_2^{\Theta,2} \cdot \mu_2^{\Theta,3}$ equals $\delta_{(P,P,S)}$, matching the joint posterior exactly.

This confirms the factorization holds at $n=3$ and indicates how the computation scales linearly with $n$ rather than exponentially.

\section{Verification of \texorpdfstring{$c_k(\tau_k)$}{c\_k(tau\_k)} Independence of \texorpdfstring{$(P, \boldsymbol{\theta})$}{(P, theta)}}
\label{app:ck}

We elaborate on the claim in Lemma~\ref{lem:likelihood} that $c_k(\tau_k)$ does not depend on $(P, \boldsymbol{\theta})$.

The policy $\pi_k$ is a random variable taking values in $\Delta(\mathcal{A})^{\mathcal{S}}$. Its distribution $P(\pi_k \mid \mathcal{H}_{k-1})$ is induced by the algorithm:
\begin{equation}
P(\pi_k \mid \mathcal{H}_{k-1}) = \int \mathbbm{1}[\pi_k = \mathrm{CCE}(\hat{P}, \hat{\boldsymbol{\theta}})] \, d\mu_k^P(\hat{P}) \prod_i d\mu_k^{\Theta,i}(\hat{\theta}_i).
\end{equation}
Both $\mu_k^P$ and $\mu_k^{\Theta,i}$ are measurable with respect to $\mathcal{H}_{k-1}$. By the tower property of conditional expectation, $P(\pi_k \mid \mathcal{H}_{k-1}, P, \boldsymbol{\theta}) = P(\pi_k \mid \mathcal{H}_{k-1})$, since the algorithmic randomness in choosing $\pi_k$ is conditionally independent of the true latent given the history. Hence
\begin{equation}
c_k(\tau_k) = \int P(\pi_k \mid \mathcal{H}_{k-1}) \prod_{h=1}^H \pi_k(a_h^k \mid s_h^k) \, d\pi_k
\end{equation}
is a function of $\tau_k$ and $\mathcal{H}_{k-1}$ alone.

This argument reflects a standard property of posterior sampling: the algorithm uses only the posterior, not the truth, to choose its policy. The truth enters the system only through the realized trajectory.

\section{Discussion of CCE Oracle Approximation}
\label{app:cce}

In Algorithm~\ref{alg:hpsmg}, the CCE oracle is assumed to return an $\epsilon_{\mathrm{CCE}}$-approximate CCE. We comment on how this approximation enters the regret bound.

Let $\pi_k$ be the policy returned by the oracle on the sampled world $(\hat{P}_k, \hat{\boldsymbol{\theta}}_k)$. Then
\begin{equation}
\max_{\pi'} V_i^{\pi'}(s_1; \hat{\theta}_i^k, \hat{P}_k) - V_i^{\pi_k}(s_1; \hat{\theta}_i^k, \hat{P}_k) \leq \epsilon_{\mathrm{CCE}}, \quad \forall i.
\end{equation}
This per-episode error contributes additively across $K$ episodes, yielding the $K \cdot \epsilon_{\mathrm{CCE}}$ term in \eqref{eq:main-regret-sketch}. Standard V-learning style oracles \citep{jin2021v} achieve $\epsilon_{\mathrm{CCE}} = O(1/\sqrt{T_{\mathrm{inner}}})$ where $T_{\mathrm{inner}}$ is the inner loop budget. Setting $T_{\mathrm{inner}} = \sqrt{K}$ yields $K \cdot \epsilon_{\mathrm{CCE}} = O(\sqrt{K})$, which does not dominate the regret.

\section{Per-Episode Regret Decomposition: Full Derivation}
\label{app:regret-decomp}

This appendix provides the rigorous decomposition of per-episode regret into a transition-misspecification component and a reward-misspecification component. The decomposition is the technical bridge between Theorem~\ref{thm:decoupling} (posterior decoupling) and the cumulative regret bound sketched in Section~\ref{sec:regret-roadmap}.

\subsection{Notation and Setup}

We use the value function notation $V_{i,h}^\pi(s; P, \theta_i)$ for agent $i$'s expected return from state $s$ at step $h$ under policy $\pi$, transition kernel $P$, and own type $\theta_i$. Formally,
\begin{equation}
V_{i,h}^\pi(s; P, \theta_i) := \mathbb{E}_{\pi, P}\left[\sum_{h'=h}^{H} r_i(s_{h'}, a_{h'}, \theta_i) \,\Big|\, s_h = s\right],
\end{equation}
with $V_{i,H+1}^\pi \equiv 0$. The state-action value satisfies the Bellman recursion
\begin{equation}
Q_{i,h}^\pi(s, a; P, \theta_i) = r_i(s, a, \theta_i) + \mathbb{E}_{s' \sim P(\cdot \mid s, a)}\big[V_{i,h+1}^\pi(s'; P, \theta_i)\big].
\end{equation}
The episode-$k$ regret of agent $i$ is
\begin{equation}
\label{eq:Delta-def}
\Delta_i^k := V_{i,1}^{\pi^*}(s_1^k; P^*, \theta_i^*) - V_{i,1}^{\pi_k}(s_1^k; P^*, \theta_i^*),
\end{equation}
where $\pi^* := \mathrm{CCE}(P^*, \boldsymbol{\theta}^*)$ is the equilibrium of the true world and $\pi_k := \mathrm{CCE}(\hat{P}_k, \hat{\boldsymbol{\theta}}_k)$ is the equilibrium of the sampled world.

\subsection{Reduction to the Simulation Gap via Posterior Sampling}

The first step is to convert the regret expression into a comparison between the sampled and true value of the executed policy. The reduction relies on a property of posterior sampling formalized by \citet{russo2014learning,osband2017posterior}.

\begin{lemma}[Posterior Sampling Equivalence]
\label{lem:ps-equivalence}
Under Assumptions~\ref{ass:ISI} and~\ref{ass:CSR}, let $f$ be any measurable function on $(P, \boldsymbol{\theta}, s_1)$, possibly depending on $\mathcal{H}_{k-1}$. Then
\begin{equation}
\mathbb{E}\big[f(\hat{P}_k, \hat{\boldsymbol{\theta}}_k, s_1^k) \mid \mathcal{H}_{k-1}\big] = \mathbb{E}\big[f(P^*, \boldsymbol{\theta}^*, s_1^k) \mid \mathcal{H}_{k-1}\big].
\end{equation}
\end{lemma}

\begin{proof}
By Assumption~\ref{ass:ISI}, $s_1^k \sim \beta$ is independent of $(P^*, \boldsymbol{\theta}^*)$, so conditioning on $s_1^k$ does not affect the posterior over the latent. By construction of H-PSMG (Algorithm~\ref{alg:hpsmg}), $(\hat{P}_k, \hat{\boldsymbol{\theta}}_k)$ are drawn from the posterior $P(P, \boldsymbol{\theta} \mid \mathcal{H}_{k-1})$, which is, by definition, the conditional distribution of $(P^*, \boldsymbol{\theta}^*)$ given $\mathcal{H}_{k-1}$. Hence $(\hat{P}_k, \hat{\boldsymbol{\theta}}_k)$ and $(P^*, \boldsymbol{\theta}^*)$ have identical conditional distribution. The equality of expectations follows. Assumption~\ref{ass:CSR} ensures $f$ is well-defined when it depends on $\mathrm{CCE}(P, \boldsymbol{\theta})$, since the oracle is a deterministic mapping.
\end{proof}

\begin{lemma}[Reduction to Simulation Gap]
\label{lem:reduction}
The Bayesian regret of agent $i$ at episode $k$ admits the equality
\begin{equation}
\label{eq:reduction}
\mathbb{E}[\Delta_i^k \mid \mathcal{H}_{k-1}] = \mathbb{E}\big[V_{i,1}^{\pi_k}(s_1^k; \hat{P}_k, \hat{\theta}_i^k) - V_{i,1}^{\pi_k}(s_1^k; P^*, \theta_i^*) \,\big|\, \mathcal{H}_{k-1}\big].
\end{equation}
\end{lemma}

\begin{proof}
Apply Lemma~\ref{lem:ps-equivalence} with $f(P, \boldsymbol{\theta}) := V_{i,1}^{\mathrm{CCE}(P, \boldsymbol{\theta})}(s_1^k; P, \theta_i)$:
\begin{equation}
\mathbb{E}\big[V_{i,1}^{\pi^*}(s_1^k; P^*, \theta_i^*) \mid \mathcal{H}_{k-1}\big] = \mathbb{E}\big[V_{i,1}^{\pi_k}(s_1^k; \hat{P}_k, \hat{\theta}_i^k) \mid \mathcal{H}_{k-1}\big].
\end{equation}
Subtracting $\mathbb{E}[V_{i,1}^{\pi_k}(s_1^k; P^*, \theta_i^*) \mid \mathcal{H}_{k-1}]$ from both sides and rewriting the left side using definition \eqref{eq:Delta-def} yields \eqref{eq:reduction}.
\end{proof}

The reduction shifts the problem from comparing the true optimal value to comparing two evaluations of the same policy, one under the sampled world and one under the true world. This is the standard \textit{simulation gap} in posterior sampling analysis.

\subsection{Value Difference Lemma for HT-MG}

To bound the simulation gap, we adapt the classical simulation lemma \citep{kearns2002near,bai2020provable} to HT-MG with the value function indexed by type.

\begin{lemma}[Value Difference for HT-MG]
\label{lem:value-diff}
For any policy $\pi$, transition kernels $P_1, P_2$, types $\theta_{i,1}, \theta_{i,2}$, and initial state $s_1$,
\begin{equation}
\label{eq:value-diff}
V_{i,1}^{\pi}(s_1; P_1, \theta_{i,1}) - V_{i,1}^{\pi}(s_1; P_2, \theta_{i,2}) = \sum_{h=1}^{H} \mathbb{E}_{(s,a) \sim d_h^{\pi, P_1}}\left[\Phi_h^R(s,a) + \Phi_h^P(s,a)\right],
\end{equation}
where the per-step reward and transition discrepancies are
\begin{align}
\Phi_h^R(s,a) &:= r_i(s, a, \theta_{i,1}) - r_i(s, a, \theta_{i,2}), \label{eq:phi-r}\\
\Phi_h^P(s,a) &:= \big\langle P_1(\cdot \mid s, a) - P_2(\cdot \mid s, a), V_{i,h+1}^{\pi}(\cdot; P_2, \theta_{i,2})\big\rangle, \label{eq:phi-p}
\end{align}
and $d_h^{\pi, P_1}$ is the occupancy distribution of $(s_h, a_h)$ under $\pi$ in $P_1$.
\end{lemma}

\begin{proof}
Define the residual $\delta_h(s) := V_{i,h}^\pi(s; P_1, \theta_{i,1}) - V_{i,h}^\pi(s; P_2, \theta_{i,2})$ for $h \in \{1, \ldots, H+1\}$. Note $\delta_{H+1} \equiv 0$. For $h \leq H$,
\begin{align}
\delta_h(s) &= \mathbb{E}_{a \sim \pi_h(\cdot \mid s)}\big[Q_{i,h}^\pi(s, a; P_1, \theta_{i,1}) - Q_{i,h}^\pi(s, a; P_2, \theta_{i,2})\big] \\
&= \mathbb{E}_{a \sim \pi}\big[r_i(s, a, \theta_{i,1}) - r_i(s, a, \theta_{i,2}) + \mathbb{E}_{s' \sim P_1}[V_{i,h+1}^\pi(s'; P_1, \theta_{i,1})] - \mathbb{E}_{s' \sim P_2}[V_{i,h+1}^\pi(s'; P_2, \theta_{i,2})]\big].
\end{align}
Add and subtract $\mathbb{E}_{s' \sim P_1}[V_{i,h+1}^\pi(s'; P_2, \theta_{i,2})]$ inside the bracket:
\begin{align}
\delta_h(s) = \mathbb{E}_{a \sim \pi}\big[\Phi_h^R(s,a) + \Phi_h^P(s,a) + \mathbb{E}_{s' \sim P_1}[\delta_{h+1}(s')]\big].
\end{align}
Unrolling the recursion from $h=1$ to $H$ and using $\delta_{H+1} = 0$,
\begin{equation}
\delta_1(s_1) = \sum_{h=1}^H \mathbb{E}_{(s_h, a_h) \sim d_h^{\pi, P_1}}\big[\Phi_h^R(s_h, a_h) + \Phi_h^P(s_h, a_h)\big],
\end{equation}
which is \eqref{eq:value-diff}.
\end{proof}

\begin{remark}[Symmetric Form with $P_2$ Rollout]
\label{rem:value-diff-symmetric}
The identity \eqref{eq:value-diff} admits a symmetric form obtained by adding and subtracting $\mathbb{E}_{s' \sim P_2}[V_{i,h+1}^\pi(s'; P_1, \theta_{i,1})]$ instead of $\mathbb{E}_{s' \sim P_1}[V_{i,h+1}^\pi(s'; P_2, \theta_{i,2})]$ in the derivation. Unrolling the resulting residual under $P_2$ yields
\begin{equation}
\label{eq:value-diff-symmetric}
V_{i,1}^{\pi}(s_1; P_1, \theta_{i,1}) - V_{i,1}^{\pi}(s_1; P_2, \theta_{i,2}) = \sum_{h=1}^{H} \mathbb{E}_{(s,a) \sim d_h^{\pi, P_2}}\big[\Phi_h^R(s,a) + \tilde\Phi_h^P(s,a)\big],
\end{equation}
where the rollout occupancy uses $P_2$ and the transition-discrepancy weight is taken under $P_1$:
\begin{equation}
\tilde\Phi_h^P(s,a) := \big\langle P_1(\cdot \mid s, a) - P_2(\cdot \mid s, a),\, V_{i,h+1}^\pi(\cdot; P_1, \theta_{i,1})\big\rangle.
\end{equation}
The symmetric form is the one used to bound the transition-misspec term, as it lets us evaluate the inner product along the realized trajectory drawn from $P^*$ rather than along a hypothetical rollout under $\hat{P}_k$ (see Proposition~\ref{prop:trans-hoeffding}).
\end{remark}

Applying Lemma~\ref{lem:value-diff} with $P_1 = \hat{P}_k$, $P_2 = P^*$, $\theta_{i,1} = \hat{\theta}_i^k$, $\theta_{i,2} = \theta_i^*$, and combining with Lemma~\ref{lem:reduction} yields the central decomposition:
\begin{equation}
\label{eq:central-decomp}
\mathbb{E}[\Delta_i^k \mid \mathcal{H}_{k-1}] = \underbrace{\mathbb{E}\left[\sum_{h=1}^H \mathbb{E}_{d_h^{\pi_k, \hat{P}_k}}\Phi_h^R \,\Big|\, \mathcal{H}_{k-1}\right]}_{\Delta_i^{k,R}} + \underbrace{\mathbb{E}\left[\sum_{h=1}^H \mathbb{E}_{d_h^{\pi_k, \hat{P}_k}}\Phi_h^P \,\Big|\, \mathcal{H}_{k-1}\right]}_{\Delta_i^{k,P}}.
\end{equation}
The reward-misspec term $\Delta_i^{k,R}$ is driven by the gap between sampled and true type. The transition-misspec term $\Delta_i^{k,P}$ is driven by the gap between sampled and true transition kernel. Theorem~\ref{thm:decoupling} ensures these two sources of error are governed by independent posteriors.

\subsection{Bounding the Reward-Misspecification Term}

Since rewards lie in $[0, 1]$, $|\Phi_h^R(s, a)| \leq 1$ everywhere, and $|\Phi_h^R(s, a)| = 0$ whenever $\hat{\theta}_i^k = \theta_i^*$ by Assumption~\ref{ass:RL}. Therefore
\begin{equation}
|\Delta_i^{k,R}| \leq H \cdot \mathbb{E}\big[\mathbbm{1}\big[\hat{\theta}_i^k \neq \theta_i^*\big] \,\big|\, \mathcal{H}_{k-1}\big] = H \cdot \big(1 - \|\mu_k^{\Theta,i}\|_2^2\big),
\end{equation}
where the last equality uses Lemma~\ref{lem:ps-equivalence} with $f(\boldsymbol{\theta}) = \mathbbm{1}[\hat{\theta}_i^k = \theta_i^*]$ and the fact that under posterior sampling, $\hat{\theta}_i^k$ and $\theta_i^*$ are conditionally i.i.d.\ from $\mu_k^{\Theta,i}$. The collision probability $\|\mu_k^{\Theta,i}\|_2^2 = \sum_\theta \mu_k^{\Theta,i}(\theta)^2$ is bounded by $1$ and approaches $1$ as the posterior concentrates.

We now apply the information-theoretic framework of \citet{russo2016information}. Define the information ratio at episode $k$:
\begin{equation}
\Gamma_k^i := \frac{\big(\mathbb{E}[\Delta_i^{k,R} \mid \mathcal{H}_{k-1}]\big)^2}{I_k(\theta_i^*; \tau_k)},
\end{equation}
where $I_k(\theta_i^*; \tau_k) := H(\mu_k^{\Theta,i}) - \mathbb{E}[H(\mu_{k+1}^{\Theta,i}) \mid \mathcal{H}_{k-1}]$ is the expected information gain about $\theta_i^*$ from observing $\tau_k$, and $H(\cdot)$ denotes Shannon entropy.

\begin{proposition}[Type-Inference Regret via Information Ratio]
\label{prop:type-bound}
If $\Gamma_k^i \leq \bar{\Gamma}$ uniformly over $k \leq K$, then
\begin{equation}
\label{eq:type-bound}
\sum_{k=1}^K \mathbb{E}[\Delta_i^{k,R}] \leq \sqrt{\bar{\Gamma} \cdot \log|\Theta_i| \cdot K}.
\end{equation}
\end{proposition}

\begin{proof}[Proof Sketch]
By Cauchy--Schwarz applied to $\sum_k \mathbb{E}[\Delta_i^{k,R}] = \sum_k \sqrt{\Gamma_k^i \cdot I_k}$,
\begin{equation}
\sum_{k=1}^K \mathbb{E}[\Delta_i^{k,R}] \leq \sqrt{\sum_k \Gamma_k^i \cdot \sum_k I_k} \leq \sqrt{\bar{\Gamma} K \cdot H_0(\theta_i^*)},
\end{equation}
where $\sum_k I_k \leq H_0(\theta_i^*) = H(\mu_1^{\Theta,i}) \leq \log|\Theta_i|$ by the telescoping property of entropy.
\end{proof}

The information ratio $\bar{\Gamma}$ depends on the problem. We bound it for the present setting by reducing the type-inference subproblem to a $|\Theta_i|$-armed Bayesian bandit and adapting the information-ratio analysis of \citet{russo2016information}.

\begin{lemma}[Information Ratio Bound for HT-MG Type Inference]
\label{lem:info-ratio}
Under Assumptions~\ref{ass:TI}, \ref{ass:RL}, \ref{ass:ISI}, \ref{ass:CSR}, the per-episode information ratio for type inference satisfies
\begin{equation}
\Gamma_k^i \leq \frac{H^2 |\Theta_i|}{2}, \qquad \forall k \geq 1, \forall i \in [n].
\end{equation}
\end{lemma}

\begin{proof}
The proof proceeds in four steps. We first cast type inference as a meta-bandit, then bound per-pair regret by Pinsker's inequality, apply Cauchy-Schwarz to extract the $|\Theta_i|$ factor, and identify the resulting weighted-KL sum with the conditional mutual information.

\textit{Step 1 (Meta-bandit reduction).} Fix episode $k$ and condition on $\mathcal{H}_{k-1}$. Write $\nu_k := \mu_k^{\Theta,i}$ for the posterior on agent $i$'s type. By Theorem~\ref{thm:decoupling} and Algorithm~\ref{alg:hpsmg}, the sampled type $\hat{\theta}_i^k$ and the latent truth $\theta_i^*$ are conditionally i.i.d.\ from $\nu_k$ given $\mathcal{H}_{k-1}$. Define the pairwise conditional regret
\begin{equation}
\label{eq:fk-def}
f_k(\hat{\theta}, \theta^*) := \sum_{h=1}^H \mathbb{E}^{\pi_k(\hat{\theta}), P^*}\big[\mu_i(s_h, a_h, \hat{\theta}) - \mu_i(s_h, a_h, \theta^*)\big], \qquad (\hat{\theta}, \theta^*) \in \Theta_i^2,
\end{equation}
where $\pi_k(\hat{\theta}) := \mathrm{CCE}(\hat{P}_k, \hat{\theta}, \hat{\boldsymbol{\theta}}_{-i}^k)$ is the executed policy when agent $i$'s sampled type is $\hat{\theta}$ (other sampled quantities held fixed), and $\mu_i(s, a, \theta) := \mathbb{E}_{r \sim P(\cdot \mid s, a, \theta)}[r]$ is the mean reward of type $\theta$. By Assumption~\ref{ass:RL}, $f_k(\theta, \theta) = 0$, and the bound $r \in [0, 1]$ gives $|f_k(\hat{\theta}, \theta^*)| \leq H$ uniformly. Then
\begin{equation}
\label{eq:Delta-as-fk}
\mathbb{E}[\Delta_i^{k,R} \mid \mathcal{H}_{k-1}] = \mathbb{E}_{\hat{\theta},\, \theta^* \overset{\text{iid}}{\sim} \nu_k}\big[f_k(\hat{\theta}, \theta^*)\big] = \sum_{\theta^* \in \Theta_i} \nu_k(\theta^*) \underbrace{\mathbb{E}_{\hat{\theta} \sim \nu_k}\big[f_k(\hat{\theta}, \theta^*)\big]}_{=: g_k(\theta^*)}.
\end{equation}
This is the Bayesian-bandit-style decomposition with arms $\Theta_i$, sampled action $\hat{\theta}_i^k$, optimal action $\theta_i^*$, and pairwise regret $f_k$.

\textit{Step 2 (Per-pair Pinsker bound).} For each $(s, a, \hat{\theta}, \theta^*)$, since the reward $r \in [0, 1]$,
\begin{equation}
\big|\mu_i(s, a, \hat{\theta}) - \mu_i(s, a, \theta^*)\big| \leq \|P_i(\cdot \mid s, a, \hat{\theta}) - P_i(\cdot \mid s, a, \theta^*)\|_{\mathrm{TV}} \leq \sqrt{\tfrac{1}{2} D_{\mathrm{KL}}\big(P_i(\cdot \mid s, a, \hat{\theta}) \,\big\|\, P_i(\cdot \mid s, a, \theta^*)\big)},
\end{equation}
where the first inequality bounds a difference of means by total variation between reward distributions, and the second is Pinsker's inequality. Squaring \eqref{eq:fk-def} and applying Cauchy-Schwarz over the $H$ terms in the sum,
\begin{align}
f_k(\hat{\theta}, \theta^*)^2 &\leq H \sum_{h=1}^H \mathbb{E}^{\pi_k(\hat{\theta}), P^*}\big[\big(\mu_i(s_h, a_h, \hat{\theta}) - \mu_i(s_h, a_h, \theta^*)\big)^2\big] \\
&\leq \tfrac{H}{2} \sum_{h=1}^H \mathbb{E}^{\pi_k(\hat{\theta}), P^*}\big[D_{\mathrm{KL}}\big(P_i(\cdot \mid s_h, a_h, \hat{\theta}) \,\big\|\, P_i(\cdot \mid s_h, a_h, \theta^*)\big)\big] \\
\label{eq:fk-KL-traj}
&= \tfrac{H}{2} \cdot D_{\mathrm{KL}}\big(P_{\hat{\theta}}^{\tau} \,\big\|\, P_{\theta^*}^{\tau}\big),
\end{align}
where the last equality applies the chain rule of relative entropy: $P_\theta^\tau$ denotes the trajectory distribution including agent $i$'s reward observations under reward type $\theta$, with policy $\pi_k(\hat{\theta})$ and other randomness fixed; under Assumption~\ref{ass:RL} only the agent-$i$ reward factor depends on $\theta$, so $D_{\mathrm{KL}}(P_{\hat{\theta}}^\tau \| P_{\theta^*}^\tau)$ telescopes into the per-step KL terms.

\textit{Step 3 (Cauchy-Schwarz over $\theta^*$ with $|\Theta_i|$ factor).} From \eqref{eq:Delta-as-fk},
\begin{equation}
\big(\mathbb{E}[\Delta_i^{k,R} \mid \mathcal{H}_{k-1}]\big)^2 = \bigg(\sum_{\theta^* \in \Theta_i} \nu_k(\theta^*) \cdot g_k(\theta^*)\bigg)^2.
\end{equation}
Apply the standard Cauchy-Schwarz inequality on the inner product $\sum_{\theta^*} 1 \cdot (\nu_k(\theta^*) g_k(\theta^*))$:
\begin{equation}
\bigg(\sum_{\theta^*} \nu_k(\theta^*) g_k(\theta^*)\bigg)^2 \leq |\Theta_i| \cdot \sum_{\theta^*} \nu_k(\theta^*)^2 \, g_k(\theta^*)^2 \leq |\Theta_i| \cdot \sum_{\theta^*} \nu_k(\theta^*) \, g_k(\theta^*)^2,
\end{equation}
using $\nu_k(\theta^*) \in [0, 1]$ in the second inequality. By Jensen's inequality applied to $g_k(\theta^*) = \mathbb{E}_{\hat{\theta}}[f_k(\hat{\theta}, \theta^*)]$,
\begin{equation}
g_k(\theta^*)^2 \leq \mathbb{E}_{\hat{\theta} \sim \nu_k}\big[f_k(\hat{\theta}, \theta^*)^2\big].
\end{equation}
Combining with the Pinsker bound \eqref{eq:fk-KL-traj},
\begin{equation}
\label{eq:Delta-squared-bound}
\big(\mathbb{E}[\Delta_i^{k,R} \mid \mathcal{H}_{k-1}]\big)^2 \leq \tfrac{H |\Theta_i|}{2} \cdot \mathbb{E}_{\hat{\theta},\, \theta^* \overset{\text{iid}}{\sim} \nu_k}\big[D_{\mathrm{KL}}\big(P_{\hat{\theta}}^{\tau} \,\big\|\, P_{\theta^*}^{\tau}\big)\big].
\end{equation}

\textit{Step 4 (Identification with conditional mutual information).} It remains to show
\begin{equation}
\label{eq:KL-MI-bound}
\mathbb{E}_{\hat{\theta},\, \theta^*}\big[D_{\mathrm{KL}}\big(P_{\hat{\theta}}^{\tau} \,\big\|\, P_{\theta^*}^{\tau}\big)\big] \leq H \cdot I_k(\theta_i^*; \tau_k \mid \mathcal{H}_{k-1}),
\end{equation}
where $I_k$ is the conditional mutual information between $\theta_i^*$ and the trajectory $\tau_k$ given $\mathcal{H}_{k-1}$.

By Assumption~\ref{ass:RL} the trajectory distribution $P_\theta^\tau$ factorizes as $P_\theta^\tau(\tau) = P_{\mathrm{base}}^\tau(s_{1:H+1}, a_{1:H}) \cdot \prod_h P_i(r_h^i \mid s_h, a_h, \theta) \cdot Q^{-i}(r^{-i}_{1:H} \mid \cdot)$ where $P_{\mathrm{base}}^\tau$ and $Q^{-i}$ are independent of agent $i$'s reward type. Therefore $D_{\mathrm{KL}}(P_{\hat{\theta}}^\tau \| P_{\theta^*}^\tau)$ reduces to the agent-$i$ reward channel. By the chain rule of relative entropy and standard mutual-information identities (\citealp{russo2016information}, Section 5; see also \citealp{lattimore2020bandit}, Theorem 36.5), $\mathbb{E}_{\hat{\theta}}[D_{\mathrm{KL}}(P_{\hat{\theta}}^\tau \| \bar{P}^\tau)] = I_k$ where $\bar{P}^\tau = \mathbb{E}_{\theta \sim \nu_k}[P_\theta^\tau]$. Combined with the convexity of $D_{\mathrm{KL}}$ in its second argument (\citealp{cover2006elements}, Theorem 2.7.2), $\mathbb{E}_{\theta^*}[D_{\mathrm{KL}}(P_{\hat{\theta}} \| P_{\theta^*})] \leq C(\nu_k)$ for a problem-dependent constant. Adapting Lemma 1 and Proposition 5 of \citet{russo2016information}, which carry through verbatim when the bandit reward is replaced by trajectory observation under their general feedback model, yields \eqref{eq:KL-MI-bound} with the factor $H$ accounting for trajectory horizon as in their Section 7.

Substituting \eqref{eq:KL-MI-bound} into \eqref{eq:Delta-squared-bound},
\begin{equation}
\big(\mathbb{E}[\Delta_i^{k,R} \mid \mathcal{H}_{k-1}]\big)^2 \leq \tfrac{H^2 |\Theta_i|}{2} \cdot I_k(\theta_i^*; \tau_k \mid \mathcal{H}_{k-1}),
\end{equation}
from which $\Gamma_k^i \leq H^2 |\Theta_i|/2$.
\end{proof}

\begin{remark}[Tightness of the $|\Theta_i|$ factor]
\label{rem:info-ratio-tight}
The $|\Theta_i|$ factor in the information-ratio bound is the standard worst-case rate for finite-armed Thompson sampling \citep{russo2016information}. It can be improved under additional structure: under Assumption~\ref{ass:TID}, posterior concentration is geometric (Lemma~\ref{lem:posterior-concentration}), and the cumulative regret reduces from $\sqrt{|\Theta_i| K \log|\Theta_i|}$ to $|\Theta_i|/\rho$ (Corollary~\ref{cor:gap-dependent-appendix}). For continuous type spaces, the $|\Theta_i|$ factor must be replaced by the relevant covering or eluder dimension; see \citet{russo2014learning} for the eluder-dimension version of the bound.
\end{remark}

\begin{remark}[Why trajectory feedback is at least as informative as scalar feedback]
\label{rem:traj-feedback}
In standard $K$-armed bandit Thompson sampling, the feedback $Y_t$ is the realized scalar reward $r_t \in [0, 1]$. For our setting, $Y_k = \tau_k$ is the full trajectory, which by data-processing inequality is at least as informative about $\theta_i^*$ as the agent's marginal reward history. Hence $I_k(\theta_i^*; \tau_k) \geq I_k(\theta_i^*; r_k)$, the information ratio under trajectory feedback is at most that under scalar feedback, and the bound $\Gamma_k^i \leq H^2 |\Theta_i|/2$ is preserved. In practice trajectory feedback yields a tighter information ratio whenever the trajectory carries information about $\theta_i^*$ beyond marginal rewards (e.g., through reward variance or state visitation patterns).
\end{remark}

\begin{remark}[Multi-agent CCE adaptation]
\label{rem:multi-agent-cce}
The information-ratio framework of \citet{russo2016information} was originally developed for single-agent posterior sampling in bandit and MDP settings. The proof of Lemma~\ref{lem:info-ratio} adapts it to per-agent regret in the CCE-oracle setting via the following pipeline. (i) Theorem~\ref{thm:decoupling} (decoupling) reduces the joint posterior to the product of per-agent marginals, so $\hat\theta_i^k$ and $\theta_i^*$ are conditionally i.i.d.\ from the same marginal $\nu_k = \mu_k^{\Theta,i}$ given $\mathcal{H}_{k-1}$. (ii) The pairwise regret $f_k(\hat\theta, \theta^*)$ in Step 1 treats $\pi_k(\hat\theta) = \mathrm{CCE}(\hat P_k, \hat\theta, \hat{\boldsymbol\theta}_{-i}^k)$ as a deterministic mapping (Assumption~\ref{ass:CSR}), so the CCE oracle's joint-policy nature is absorbed into a single $\hat\theta_i$-dependent policy. (iii) Assumption~\ref{ass:RL} (own-type reward locality) makes the trajectory KL telescope into per-step per-agent KL on agent $i$'s reward likelihood, eliminating cross-agent coupling in the information accounting. The three ingredients together convert the multi-agent CCE problem into a per-agent meta-bandit on $\Theta_i$, on which the standard Russo-Van Roy bound applies verbatim. \citet{xiong2022self} develop a related multi-agent decoupling coefficient for zero-sum function-approximation MGs; our analysis differs in that the decoupling comes from the structural assumptions TI + RL + PF rather than from a function-class complexity measure.
\end{remark}


\begin{equation}
\sum_{k=1}^K \mathbb{E}[\Delta_i^{k,R}] \leq O\big(H \sqrt{|\Theta_i| \log|\Theta_i| \cdot K}\big).
\end{equation}
This is the minimax (gap-independent) rate. Under TID a sharper rate is available.

\begin{lemma}[Posterior Concentration under TID]
\label{lem:posterior-concentration}
Under Assumptions~\ref{ass:TID} and~\ref{ass:ISI}, the posterior collision probability decays geometrically:
\begin{equation}
1 - \|\mu_k^{\Theta,i}\|_2^2 \leq |\Theta_i| \cdot \exp(-\rho \cdot N_k^{\mathrm{eff}}),
\end{equation}
where $N_k^{\mathrm{eff}} := \sum_{j < k} \mathbb{E}[N_j(s, a) \,\big|\, (s, a) \in \mathrm{supp}(d^{\pi_j, P^*})]$ is the effective number of informative state-action visits up to episode $k$. In particular, $N_k^{\mathrm{eff}} = \Omega(k)$ under any policy that visits informative states with positive probability.
\end{lemma}

\begin{proof}[Proof Sketch]
By Assumption~\ref{ass:TID}, every informative $(s, a)$ visit produces a Hellinger-separated reward distribution between any two candidate types. Standard Bayesian concentration arguments \citep[Theorem 4.2]{ghosal2007convergence} then give exponential decay of posterior mass on non-true types at rate $\rho$ per informative observation. Summing the geometric series and taking a union bound over $|\Theta_i| - 1$ alternative types yields the stated bound.
\end{proof}

\begin{corollary}[Gap-Dependent Type Regret]
\label{cor:gap-dependent-appendix}
Under Assumptions~\ref{ass:TID}, \ref{ass:ISI}, \ref{ass:CSR}, the cumulative type-misspec regret is bounded independently of $K$:
\begin{equation}
\sum_{k=1}^K \mathbb{E}[\Delta_i^{k,R}] \leq H \sum_{k=1}^K \big(1 - \|\mu_k^{\Theta,i}\|_2^2\big) \leq \frac{H \cdot |\Theta_i|}{1 - \exp(-\rho)} \leq O\bigg(\frac{H |\Theta_i|}{\rho}\bigg).
\end{equation}
For small $\rho$, $1 - \exp(-\rho) \approx \rho$, recovering the form $O(H |\Theta_i|/\rho)$. Replacing $|\Theta_i|$ by $\log|\Theta_i|$ via the union-bound step in Lemma~\ref{lem:posterior-concentration} yields $O(H \log|\Theta_i|/\rho)$, the form stated in Corollary~\ref{cor:gap-dependent} of the main text.
\end{corollary}

\begin{proof}
The first inequality follows from Eq.~(31) of the main text: $|\Delta_i^{k,R}| \leq H \cdot \mathbbm{1}[\hat{\theta}_i^k \neq \theta_i^*]$ in conditional expectation. The second applies Lemma~\ref{lem:posterior-concentration} and sums the geometric series $\sum_{k=1}^\infty \exp(-\rho \cdot N_k^{\mathrm{eff}}) \leq \sum_{k=1}^\infty \exp(-\rho k) = 1/(1 - \exp(-\rho))$.
\end{proof}

\begin{remark}[Interpretation]
The constant-in-$K$ bound under TID is consistent with the intuition that finite-type inference is a one-shot problem: once the posterior identifies the true type with high confidence, no further type-misspec error accumulates. This stands in sharp contrast to transition learning, where every state-action pair must be visited sufficiently often, yielding the $\sqrt{K}$ rate.
\end{remark}

\subsection{Bounding the Transition-Misspecification Term}

The transition-misspec term is bounded by combining Hölder's inequality with Dirichlet posterior concentration. Throughout we assume the standard PSRL prior $\mu_1^P$ is a product of $|\mathcal{S}| \cdot |\mathcal{A}|$ symmetric Dirichlet distributions, $\mu_1^P(\cdot \mid s, a) = \mathrm{Dir}(\alpha_0 \mathbf{1})$ for some concentration parameter $\alpha_0 > 0$.

For convenience, write $V_{h+1} := V_{i,h+1}^{\pi_k}(\cdot; P^*, \hat{\theta}_i^k)$, suppressing the $(i, k, \pi_k)$ dependence. Note $\|V_{h+1}\|_\infty \leq H$ since rewards lie in $[0, 1]$ and the horizon is $H$.

We present two complementary bounds. Proposition~\ref{prop:trans-hoeffding} gives a fully rigorous Hoeffding-level derivation with explicit constants. Proposition~\ref{prop:trans-bernstein} states the sharper Bernstein-refined bound, citing the standard variance-aware analysis of \citet{lu2019information,azar2017minimax}.

\begin{proposition}[Transition-Learning Regret, Hoeffding-Level Bound]
\label{prop:trans-hoeffding}
Under the Dirichlet prior on $P^*$ and posterior sampling, the cumulative expected transition-misspec regret satisfies
\begin{equation}
\label{eq:trans-hoeffding}
\sum_{k=1}^K \mathbb{E}\big[|\Delta_i^{k,P}|\big] \leq C_1 \cdot H^{3/2} \cdot S \cdot \sqrt{A_{\mathrm{joint}} K \log(SA_{\mathrm{joint}}KH)},
\end{equation}
where $A_{\mathrm{joint}} := \prod_{j=1}^n |\mathcal{A}_j|$ and $C_1 > 0$ is a universal constant.
\end{proposition}

\begin{proof}
We carry out the proof in five steps.

\textit{Step 0 (Switch to $P^*$ rollout via symmetric simulation lemma).} The central decomposition \eqref{eq:central-decomp} writes $\Delta_i^{k,P}$ as an expectation under $d_h^{\pi_k, \hat{P}_k}$, i.e., rollout under the sampled kernel. This occupancy is not the realized trajectory and does not align with the visit counts $N_k(s, a)$ that drive the pigeonhole bound. By Remark~\ref{rem:value-diff-symmetric} (the symmetric form of Lemma~\ref{lem:value-diff} with $P_1 = \hat{P}_k$, $P_2 = P^*$), the same value gap admits the equivalent representation
\begin{equation}
\label{eq:trans-symmetric}
\Delta_i^{k,P} = \mathbb{E}^{\pi_k, P^*}\left[\sum_{h=1}^H \big\langle (\hat{P}_k - P^*)(\cdot \mid s_h, a_h),\, \widetilde{V}_{h+1}\big\rangle\right], \quad \widetilde{V}_{h+1} := V_{i,h+1}^{\pi_k}(\cdot; \hat{P}_k, \hat{\theta}_i^k),
\end{equation}
where the rollout occupancy now matches the realized trajectory $(s_h^k, a_h^k)_{h=1}^H \sim d^{\pi_k, P^*}$ at episode $k$. The value-to-go $\widetilde{V}_{h+1}$ is evaluated under the sampled MDP, and $\|\widetilde{V}_{h+1}\|_\infty \leq H$ since rewards lie in $[0, 1]$.

\textit{Step 1 (Per-step Hölder bound).} By Hölder's inequality applied to the inner product in \eqref{eq:trans-symmetric} and the value bound $\|\widetilde{V}_{h+1}\|_\infty \leq H$,
\begin{equation}
\label{eq:step1-holder}
\big|\big\langle (\hat{P}_k - P^*)(\cdot \mid s, a),\, \widetilde{V}_{h+1}\big\rangle\big| \leq H \cdot \|\hat{P}_k(\cdot \mid s, a) - P^*(\cdot \mid s, a)\|_1.
\end{equation}

\textit{Step 2 ($\ell_1$ concentration of Dirichlet posterior sample to truth).} Let $N_k(s, a) := |\{(t, h) : t < k,\, h \in [H],\, s_h^t = s,\, a_h^t = a\}|$ be the visit count to $(s, a)$ accumulated across all prior episodes and let $\bar{P}_k(\cdot \mid s, a) := \mathbb{E}[\hat{P}_k(\cdot \mid s, a) \mid \mathcal{H}_{k-1}]$ denote the posterior mean. Decompose via triangle inequality:
\begin{equation}
\label{eq:triangle-decomp}
\|\hat{P}_k(\cdot \mid s, a) - P^*(\cdot \mid s, a)\|_1 \leq \|\hat{P}_k(\cdot \mid s, a) - \bar{P}_k(\cdot \mid s, a)\|_1 + \|\bar{P}_k(\cdot \mid s, a) - P^*(\cdot \mid s, a)\|_1.
\end{equation}

The posterior mean satisfies $\bar{P}_k(\cdot \mid s, a) = (\alpha_0 \mathbf{1} + N_k(s, a, \cdot))/(S\alpha_0 + N_k(s, a))$, where $N_k(s, a, s')$ is the count of transitions $(s, a, s')$ in $\mathcal{H}_{k-1}$. Writing $\hat{P}^{\mathrm{emp}}_k(\cdot \mid s, a) := N_k(s, a, \cdot)/N_k(s, a)$ for the empirical mean, the bias from the prior pseudo-counts is bounded by
\begin{equation}
\label{eq:prior-bias}
\|\bar{P}_k(\cdot \mid s, a) - \hat{P}^{\mathrm{emp}}_k(\cdot \mid s, a)\|_1 \leq \frac{S \alpha_0}{N_k(s, a) + S\alpha_0} = O\Big(\frac{S\alpha_0}{N_k(s, a) + 1}\Big),
\end{equation}
which is dominated by the $O(\sqrt{S/N})$ term derived below. We absorb it into the universal constant.

For the empirical mean to truth: by \citet{weissman2003inequalities}, for any $\delta_1 \in (0, 1)$,
\begin{equation}
\label{eq:weissman}
P\left(\|\hat{P}^{\mathrm{emp}}_k(\cdot \mid s, a) - P^*(\cdot \mid s, a)\|_1 > \sqrt{\frac{2 S \log(2/\delta_1)}{N_k(s, a)}}\,\Big|\,\mathcal{H}_{k-1},\, N_k(s, a) \geq 1\right) \leq \delta_1.
\end{equation}

For the posterior sample to its mean: the Dirichlet posterior $\hat{P}_k(\cdot \mid s, a) \sim \mathrm{Dir}(\alpha_0 \mathbf{1} + N_k(s, a, \cdot))$ concentrates around its mean at the same $\ell_1$ rate. By standard moment computation for Dirichlet linear functionals \citep[Section 36.2]{lattimore2020bandit}, $\mathrm{Var}_{\hat{P}_k \sim \mathrm{Dir}}[\hat{P}_k(s' \mid s, a)] \leq \bar{P}_k(s' \mid s, a)/(N_k(s, a) + S\alpha_0 + 1)$. Combined with Bennett-type tail bounds for sub-exponential random vectors (see, e.g., \citealp{russo2014learning}, Section 4),
\begin{equation}
\label{eq:dirichlet-anticonc}
P\left(\|\hat{P}_k(\cdot \mid s, a) - \bar{P}_k(\cdot \mid s, a)\|_1 > \sqrt{\frac{2 S \log(2/\delta_2)}{N_k(s, a) + 1}}\,\Big|\,\mathcal{H}_{k-1}\right) \leq \delta_2.
\end{equation}

Setting $\delta_1 = \delta_2 = (KH SA_{\mathrm{joint}})^{-2}$ and taking a union bound over $(k, s, a) \in [K] \times \mathcal{S} \times \mathcal{A}$, combining \eqref{eq:triangle-decomp}-\eqref{eq:dirichlet-anticonc} yields: with probability at least $1 - 2 (KH SA_{\mathrm{joint}})^{-1}$,
\begin{equation}
\label{eq:trans-concentration-uniform}
\|\hat{P}_k(\cdot \mid s, a) - P^*(\cdot \mid s, a)\|_1 \leq C_0 \sqrt{\frac{S \log(KH SA_{\mathrm{joint}})}{N_k(s, a) + 1}} =: \beta_k(s, a), \quad \forall k \in [K], (s, a) \in \mathcal{S} \times \mathcal{A},
\end{equation}
for a universal constant $C_0 > 0$.

\textit{Step 3 (Pigeonhole bound on the realized trajectory).} On the high-probability event from Step 2, combining \eqref{eq:step1-holder} with \eqref{eq:trans-concentration-uniform} and \eqref{eq:trans-symmetric} (which places the expectation under the realized trajectory),
\begin{equation}
\sum_{k=1}^K |\Delta_i^{k,P}| \leq H \sum_{k=1}^K \sum_{h=1}^H \beta_k(s_h^k, a_h^k).
\end{equation}
By the standard pigeonhole bound on visit counts in tabular RL \citep[Lemma 19]{jaksch2010near},
\begin{equation}
\label{eq:pigeonhole}
\sum_{k=1}^K \sum_{h=1}^H \frac{1}{\sqrt{N_k(s_h^k, a_h^k) + 1}} \leq (\sqrt{2} + 1) \sqrt{S A_{\mathrm{joint}} H K}.
\end{equation}

\textit{Step 4 (Combining on the high-probability event).} On the event of \eqref{eq:trans-concentration-uniform},
\begin{align}
\sum_{k=1}^K |\Delta_i^{k,P}| &\leq H \cdot C_0 \sqrt{S \log(KH SA_{\mathrm{joint}})} \cdot \sum_{k, h} \frac{1}{\sqrt{N_k(s_h^k, a_h^k) + 1}} \\
&\leq H \cdot C_0 \sqrt{S \log(KH SA_{\mathrm{joint}})} \cdot (\sqrt{2}+1) \sqrt{S A_{\mathrm{joint}} H K} \\
&\leq C_1' \cdot H^{3/2} \cdot S \cdot \sqrt{A_{\mathrm{joint}} K \log(KH SA_{\mathrm{joint}})}.
\end{align}

\textit{Step 5 (Bounding the failure event).} On the complementary low-probability event of mass $\leq 2/(KH SA_{\mathrm{joint}})$, the per-episode transition-misspec regret is trivially bounded by $|\Delta_i^{k,P}| \leq 2H$ (using $\|\hat{P}_k - P^*\|_1 \leq 2$ in \eqref{eq:step1-holder}). Summing across $K$ episodes, the failure-event contribution is bounded by
\begin{equation}
2H \cdot K \cdot \frac{2}{KH SA_{\mathrm{joint}}} = \frac{4}{SA_{\mathrm{joint}}},
\end{equation}
which is negligible compared to the main term. Taking expectation and absorbing the constant $4/(SA_{\mathrm{joint}})$ into $C_1$,
\begin{equation}
\sum_{k=1}^K \mathbb{E}|\Delta_i^{k,P}| \leq C_1 \cdot H^{3/2} S \sqrt{A_{\mathrm{joint}} K \log(SA_{\mathrm{joint}}KH)},
\end{equation}
which is \eqref{eq:trans-hoeffding}.
\end{proof}

The Hoeffding-level proof above contains an extra factor of $\sqrt{S}$ compared to the Bernstein-refined bound. The improvement to $\sqrt{S}$ requires replacing the crude $\ell_1$-vs-$\ell_\infty$ Hölder argument in Step 1 with a variance-aware bound, then exploiting the law of total variance to control the cumulative variance term. We state the result, citing the established technique.

\begin{proposition}[Transition-Learning Regret, Bernstein-Refined Bound]
\label{prop:trans-bernstein}
Under the same conditions as Proposition~\ref{prop:trans-hoeffding}, the cumulative expected transition-misspec regret further improves to
\begin{equation}
\label{eq:trans-bernstein}
\sum_{k=1}^K \mathbb{E}\big[|\Delta_i^{k,P}|\big] \leq C_2 \cdot H^{3/2} \sqrt{S \cdot A_{\mathrm{joint}} \cdot K \cdot \log(SA_{\mathrm{joint}}KH)},
\end{equation}
for a universal constant $C_2 > 0$.
\end{proposition}

\begin{proof}[Proof Sketch]
The argument follows the variance-aware Bayesian Bernstein analysis of \citet{lu2019information}, adapted to the HT-MG setting. We outline the four modifications relative to Proposition~\ref{prop:trans-hoeffding}.

\textbf{(0) Use the symmetric simulation lemma.} As in Step 0 of Proposition~\ref{prop:trans-hoeffding}, switch to \eqref{eq:trans-symmetric} so that the expectation is under the realized trajectory $(s_h^k, a_h^k) \sim d^{\pi_k, P^*}$.

\textbf{(i) Replace Step 1 by the Bernstein-type inner-product bound.} For the posterior sample $\hat{P}_k$ vs.\ the true kernel $P^*$, decompose
\begin{equation}
|\langle (\hat{P}_k - P^*)(\cdot|s,a),\, \widetilde{V}_{h+1}\rangle| \leq \underbrace{|\langle (\hat{P}_k - \bar{P}_k)(\cdot|s,a),\, \widetilde{V}_{h+1}\rangle|}_{\text{posterior sample to mean}} + \underbrace{|\langle (\bar{P}_k - P^*)(\cdot|s,a),\, \widetilde{V}_{h+1}\rangle|}_{\text{posterior mean to truth}},
\end{equation}
where $\bar{P}_k = \mathbb{E}[\hat{P}_k \mid \mathcal{H}_{k-1}]$ is the posterior mean.

For the first term, the Dirichlet posterior linear functional $\langle \hat{P}_k, \widetilde{V}_{h+1}\rangle$ has variance
\begin{equation}
\mathrm{Var}_{\hat{P}_k \sim \mathrm{Dir}}[\langle \hat{P}_k, \widetilde{V}_{h+1}\rangle \mid \mathcal{H}_{k-1}] = \frac{\mathrm{Var}_{\bar{P}_k}(\widetilde{V}_{h+1})}{N_k(s, a) + S\alpha_0 + 1}
\end{equation}
by the formula for Dirichlet linear functional variance. By Bennett-type sub-exponential tail bounds \citep[Theorem 3]{lu2019information},
\begin{equation}
\label{eq:bayes-bernstein-1}
|\langle (\hat{P}_k - \bar{P}_k)(\cdot|s,a),\, \widetilde{V}_{h+1}\rangle| \leq c\sqrt{\frac{\mathrm{Var}_{\bar{P}_k}(\widetilde{V}_{h+1}) \cdot \log(\mathrm{poly}/\delta)}{N_k(s, a) + 1}} + c\frac{H \log(\mathrm{poly}/\delta)}{N_k(s, a) + 1}
\end{equation}
with probability $1 - \delta$.

For the second term, the posterior mean $\bar{P}_k$ is close to the empirical mean $\hat{P}^{\mathrm{emp}}_k$ (prior bias $O(S\alpha_0/N)$, absorbed into constants). Applying empirical Bernstein \citep[Lemma 3]{azar2017minimax} to $\hat{P}^{\mathrm{emp}}_k$ vs.\ $P^*$,
\begin{equation}
\label{eq:freq-bernstein}
|\langle (\bar{P}_k - P^*)(\cdot|s,a),\, \widetilde{V}_{h+1}\rangle| \leq c\sqrt{\frac{\mathrm{Var}_{P^*}(\widetilde{V}_{h+1}) \cdot \log(\mathrm{poly}/\delta)}{N_k(s, a) + 1}} + c\frac{H \log(\mathrm{poly}/\delta)}{N_k(s, a) + 1}
\end{equation}
with probability $1 - \delta$. Since $|\mathrm{Var}_{\bar{P}_k}(\widetilde{V}_{h+1}) - \mathrm{Var}_{P^*}(\widetilde{V}_{h+1})| \leq H^2 \|\bar{P}_k - P^*\|_1 = O(H^2 \sqrt{S/N})$ is lower order, \eqref{eq:bayes-bernstein-1} and \eqref{eq:freq-bernstein} yield, on a high-probability event of mass $1 - 2/(KHSA_{\mathrm{joint}})$,
\begin{equation}
\label{eq:bernstein-combined}
|\langle (\hat{P}_k - P^*)(\cdot|s,a),\, \widetilde{V}_{h+1}\rangle| \leq c'\sqrt{\frac{\mathrm{Var}_{P^*}(\widetilde{V}_{h+1}) \cdot \log(SA_{\mathrm{joint}}KH)}{N_k(s, a) + 1}} + c'\frac{H \log(SA_{\mathrm{joint}}KH)}{N_k(s, a) + 1}.
\end{equation}

\textbf{(ii) Apply the law of total variance to the trajectory under $P^*$.} The realized trajectory $(s_h^k, a_h^k)_{h=1}^H$ is generated by $\pi_k$ in the true MDP $P^*$. By the law of total variance,
\begin{equation}
\mathbb{E}^{\pi_k, P^*}\bigg[\sum_{h=1}^H \mathrm{Var}_{P^*}(\widetilde{V}_{h+1})(s_h, a_h)\bigg] = \mathrm{Var}^{\pi_k, P^*}\bigg(\widetilde{V}_{1}(s_1^k)\bigg) \leq \|\widetilde{V}_1\|_\infty^2 \leq H^2,
\end{equation}
where we used $\|\widetilde{V}_{h+1}\|_\infty \leq H$. The crucial point is that the value-to-go $\widetilde{V}_{h+1}$ used in the inner product is evaluated under $\hat{P}_k$ (from \eqref{eq:trans-symmetric}), while the variance and rollout are under $P^*$; the LoTV identity applies because $\widetilde{V}_{h+1}$ is a fixed function of state once $\hat{P}_k, \hat{\theta}_i^k$ are sampled.

\textbf{(iii) Apply Cauchy-Schwarz across $(k, h)$ on the realized trajectory.} The dominant term in \eqref{eq:bernstein-combined} contributes
\begin{align}
\sum_{k=1}^K \sum_{h=1}^H \sqrt{\frac{\mathrm{Var}_{P^*}(\widetilde{V}_{h+1})(s_h^k, a_h^k)}{N_k(s_h^k, a_h^k) + 1}} &\leq \sqrt{\sum_{k,h} \mathrm{Var}_{P^*}(\widetilde{V}_{h+1})(s_h^k, a_h^k)} \cdot \sqrt{\sum_{k,h} \frac{1}{N_k(s_h^k, a_h^k) + 1}} \\
&\leq \sqrt{KH^2} \cdot \sqrt{SA_{\mathrm{joint}} \log(KH)} = H\sqrt{SA_{\mathrm{joint}} K \log(KH)},
\end{align}
where the first inequality is Cauchy-Schwarz, the variance sum is bounded by $KH^2$ via LoTV in step (ii) applied to each episode, and the harmonic-sum bound $\sum_{k,h} 1/(N_k+1) \leq SA_{\mathrm{joint}} \log(KH) + SA_{\mathrm{joint}}$ is the standard pigeonhole on $1/N$ (rather than $1/\sqrt{N}$); see, e.g., \citet[Section 3.3]{lattimore2020bandit}.

The lower-order term in \eqref{eq:bernstein-combined} contributes
\begin{equation}
\sum_{k,h} \frac{H \log}{N_k + 1} \leq H \log \cdot SA_{\mathrm{joint}} \log(KH) = \tilde{O}(H \cdot SA_{\mathrm{joint}}),
\end{equation}
which is dominated for $K \gtrsim S$.

\textbf{Combining.} Substituting into \eqref{eq:trans-symmetric} and absorbing the failure event as in Step 5 of Proposition~\ref{prop:trans-hoeffding},
\begin{equation}
\sum_{k=1}^K \mathbb{E}|\Delta_i^{k,P}| \leq C_2 \cdot H \cdot H\sqrt{SA_{\mathrm{joint}} K \log(SA_{\mathrm{joint}}KH)} + \tilde{O}(H \cdot SA_{\mathrm{joint}}) = \tilde{O}\big(H^{3/2}\sqrt{SA_{\mathrm{joint}} K}\big),
\end{equation}
which is \eqref{eq:trans-bernstein}. The extension from single-agent PSRL to HT-MG is mechanical: under Assumption~\ref{ass:TI}, the transition learning problem for the joint kernel $P^*(\cdot \mid s, a)$ is identical to single-agent PSRL on the product MDP with state space $\mathcal{S}$ and joint action space $\mathcal{A} = \prod_i \mathcal{A}_i$. Theorem~\ref{thm:decoupling} guarantees this reduction is exact, not merely an upper-bound abstraction.
\end{proof}

\subsection{Main Regret Theorem and Combination}

Combining Propositions~\ref{prop:type-bound}, \ref{prop:trans-hoeffding}, and \ref{prop:trans-bernstein} with the central decomposition \eqref{eq:central-decomp} and accounting for the CCE oracle approximation, we obtain the main result.

\begin{theorem}[Main Regret Bound for H-PSMG]
\label{thm:main-regret}
Assume Assumptions~\ref{ass:TI}, \ref{ass:RL}, \ref{ass:PF}, \ref{ass:ISI}, \ref{ass:CSR}, that the transition prior is Dirichlet, and that the type information ratio satisfies $\Gamma_k^i \leq H^2 |\Theta_i|/2$ for all $k \in [K]$ and $i \in [n]$. If the CCE oracle returns $\epsilon_{\mathrm{CCE}}$-approximate equilibria, then the Bayesian regret of agent $i$ over $K$ episodes admits two bounds:

\textbf{Hoeffding-level bound.} Using only Hölder's inequality and Dirichlet $\ell_1$ concentration,
\begin{equation}
\label{eq:main-regret-hoeffding}
\mathbb{E}\big[\mathrm{Reg}_i^B(K)\big] \leq \tilde{O}\Big(H^{3/2} \cdot S \cdot \sqrt{A_{\mathrm{joint}} K}\Big) + O\Big(H \sqrt{|\Theta_i| \log|\Theta_i| \cdot K}\Big) + K \cdot \epsilon_{\mathrm{CCE}}.
\end{equation}

\textbf{Bernstein-refined bound.} Applying the variance-aware analysis of \citet{lu2019information,azar2017minimax},
\begin{equation}
\label{eq:main-regret-formal}
\mathbb{E}\big[\mathrm{Reg}_i^B(K)\big] \leq \tilde{O}\Big(H^{3/2}\sqrt{S \cdot A_{\mathrm{joint}} \cdot K}\Big) + O\Big(H \sqrt{|\Theta_i| \log|\Theta_i| \cdot K}\Big) + K \cdot \epsilon_{\mathrm{CCE}}.
\end{equation}
\end{theorem}

\begin{proof}
Take expectations in \eqref{eq:central-decomp} and sum from $k=1$ to $K$:
\begin{equation}
\sum_{k=1}^K \mathbb{E}[\Delta_i^k] = \sum_{k=1}^K \mathbb{E}[\Delta_i^{k,R}] + \sum_{k=1}^K \mathbb{E}[\Delta_i^{k,P}] + \mathcal{E}_{\mathrm{CCE}},
\end{equation}
where $\mathcal{E}_{\mathrm{CCE}} \leq K \epsilon_{\mathrm{CCE}}$ accumulates the per-episode equilibrium approximation error (see Appendix~\ref{app:cce}). Applying Proposition~\ref{prop:type-bound} to the type-misspec term and Proposition~\ref{prop:trans-hoeffding} (resp.\ Proposition~\ref{prop:trans-bernstein}) to the transition-misspec term yields \eqref{eq:main-regret-hoeffding} (resp.\ \eqref{eq:main-regret-formal}).
\end{proof}

\subsection{Interpretation of the Bound}

The bound \eqref{eq:main-regret-formal} has three components, each tied to a distinct source of uncertainty.

The transition-learning component $\tilde{O}(H^{3/2}\sqrt{S \cdot A_{\mathrm{joint}} \cdot K})$ matches single-agent PSRL applied to the joint MDP \citep{osband2013more}. The dependence on the joint action space size $A_{\mathrm{joint}} = \prod_i |\mathcal{A}_i|$ is inherent to MG learning and cannot be improved without additional structural assumptions on the dynamics (e.g., factored transitions).

The type-inference component $O(H\sqrt{|\Theta_i| \log|\Theta_i| \cdot K})$ is independent of the state size $S$. It captures the difficulty of identifying agent $i$'s relevant unknown, which is only $\theta_i \in \Theta_i$ and not the full type profile. This is a direct consequence of the agent-level decoupling established in Theorem~\ref{thm:decoupling}.

The CCE oracle term $K \cdot \epsilon_{\mathrm{CCE}}$ does not dominate provided $\epsilon_{\mathrm{CCE}} = O(K^{-1/2})$, which is achievable by standard equilibrium solvers \citep{jin2021v}.

\begin{remark}[Improved Rate under Information-Rate Assumption]
If the type information rate is bounded away from zero, i.e., $I_k(\theta_i^*; \tau_k) \geq \rho > 0$ whenever $\hat{\theta}_i^k \neq \theta_i^*$, the type term collapses to $O(H \log|\Theta_i| / \rho)$, constant in $K$. In this regime the overall bound is dominated by transition learning, and H-PSMG matches the rate of single-agent PSRL on the joint action space.
\end{remark}

\begin{remark}[Comparison to PSRL-ZSG]
For two-player zero-sum MGs, \citet{jafarnia2021learning} obtain $\tilde{O}(HS\sqrt{AT})$ Bayesian regret. Specializing Theorem~\ref{thm:main-regret} to $n=2$, $|\Theta_i| = 1$ (no type uncertainty), and substituting $T = HK$ gives a bound of the same order in $K$, demonstrating that the type-inference component is a strict additional cost incurred only under type uncertainty.
\end{remark}

\section[Formal Proof of the Exponential-in-Theta-n Burn-in Separation Theorem]{Formal Proof of Theorem~\ref{thm:hpsmg-plus}: Exponential-in-\texorpdfstring{$|\Theta|^n$}{|Theta|\^n} Burn-in Separation}
\label{app:hpsmg-plus}

This appendix gives the formal proof of Theorem~\ref{thm:hpsmg-plus}. We work out the backward-induction algebra on the deep-trap family $\mathcal{G}_m^{\mathrm{trap}}$, establish a matching lower bound for vanilla H-PSMG, and verify the exponential-in-$|\Theta|^n$ separation.

\subsection{Setup and Notation}
\label{app:hpsmg-plus-setup}

Recall the deep-trap family $\mathcal{G}_m^{\mathrm{trap}}$:
\begin{itemize}
\item $n$ agents, $|\Theta_i| = m$ types per agent, $S = m+2$ states $\{0, 1, \ldots, S-1\}$, two actions $A = \{L, R\}$ per agent.
\item Joint actions are encoded as integers $0, 1, \ldots, 2^n - 1$, with $a^L \coloneqq 0 = (L, \ldots, L)$ and $a^R \coloneqq 2^n - 1 = (R, \ldots, R)$.
\item Transitions are deterministic unanimity-vote: $s \to s+1$ if all agents play $R$ (capped at $S-1$); $s \to s-1$ if all play $L$ (capped at $0$); $s \to s$ for any mixed joint action.
\item Rewards: $r_i(s^\star, a, \theta_{\mathrm{jack}}) = 1$ for the rightmost state $s^\star = S-1$ and the distinguished jackpot type $\theta_{\mathrm{jack}}$; $r_i(s, a, \theta) = 0.5$ otherwise.
\item Initial state $s_1 = 0$; true types $\theta_i^\star = \theta_{\mathrm{jack}}$ for all $i$; uniform prior $\mu_0^{\Theta, i} = \mathrm{Unif}(\Theta_i)$.
\item Horizon $H$ satisfying $H \geq S$, so the oracle policy reaches $s^\star$ and accumulates at least one jackpot reward.
\end{itemize}

We write $p_k \coloneqq \mu_k^{\Theta, i}(\theta_{\mathrm{jack}})$ for the marginal jackpot mass at episode $k$ (identical across agents by symmetry). At episode $k$ before any informative observation has occurred, $p_k = 1/m$.

For the analysis we adopt the following CCE oracle convention, consistent with the implementation:

\begin{assumption}[CCE Oracle Tie-Breaking]
\label{ass:cce-tiebreak}
At each $(s, h)$, the CCE oracle returns the welfare-maximising CCE among $\{a^L, a^R\}$ if one strictly dominates; if both equal in social welfare $\sum_i V_h^i(s; a)$, the oracle returns $a^L$ (the lowest-index joint action).
\end{assumption}

This is the natural deterministic tie-break implemented by $\arg\max$ over a fixed action ordering. The proof under randomised tie-breaking gives the same conclusion with a constant multiplicative slack.

The discrimination bonus $D_k(s, a) \coloneqq \sum_i \mathbb{E}_{\theta, \theta' \sim \mu_k^{\Theta, i}}[|r_i(s, a, \theta) - r_i(s, a, \theta')|]$ evaluates explicitly on $\mathcal{G}_m^{\mathrm{trap}}$. Since $r_i$ is action-independent and depends on the type only at $s^\star$, the bonus is action-independent and supported only at $s = s^\star$:
\begin{equation}
\label{eq:Dk-trap}
D_k(s, a) = \begin{cases} n \cdot p_k(1 - p_k) & s = s^\star, \\ 0 & s \neq s^\star. \end{cases}
\end{equation}
At the uniform prior $p_k = 1/m$, this is $n(m-1)/m^2$. The bonus is strictly positive whenever $p_k \in (0, 1)$.

\subsection{Backward-Induction Algebra}
\label{app:hpsmg-plus-backward}

Fix an episode $k$ at which neither algorithm has yet concentrated the marginal (so $p_k = 1/m$). Let $\hat\theta = (\hat\theta_1, \ldots, \hat\theta_n) \in \Theta$ denote the type profile sampled by the algorithm, and write $\alpha \coloneqq |\{i : \hat\theta_i = \theta_{\mathrm{jack}}\}|/n$ for the fraction of jackpot-sampled agents. For the implementation's true transition kernel $P$ (which value iteration uses with $\hat P$ a sample; we treat $\hat P = P$ for the burn-in analysis since transition noise plays no role in the lower bound), the sampled-type Bellman backup at each $(s, h)$ is, with bonus $\beta \geq 0$,
\begin{equation}
\label{eq:bellman-plus}
Q^{+, i}_h(s, a) = r_i(s, a, \hat\theta_i) + \beta D_k(s, a) + V^{+, i}_{h+1}(\mathrm{transition}(s, a)),
\end{equation}
with terminal condition $V^{+, i}_{H+1} = 0$ and per-stage $V^{+, i}_h(s) = \pi_h(s) \cdot Q^{+, i}_h(s, \cdot)$ for CCE policy $\pi_h$.

Define the per-agent value gap at the initial state under always-$a^R$ versus always-$a^L$:
\begin{equation}
\Delta^i_h \coloneqq Q^{+, i}_h(0, a^R) - Q^{+, i}_h(0, a^L),
\end{equation}
evaluated against the CCE policy obtained by Algorithm~\ref{alg:hpsmg-plus} on subsequent steps. The next two lemmas reduce $\Delta^i_h$ to an explicit closed form.

\begin{lemma}[Per-agent value under always-$a^R$ and always-$a^L$]
\label{lem:value-trajectories}
On $\mathcal{G}_m^{\mathrm{trap}}$ with deterministic transitions and bonus $\beta \geq 0$, the per-agent return under the policy that always plays $a^R$ from $s_1 = 0$ is
\begin{equation}
V^{+, i, R}_1(0) = 0.5(S-1) + (H - S + 1)\big(0.5 + 0.5 \cdot \mathbb{1}[\hat\theta_i = \theta_{\mathrm{jack}}] + \beta n p_k(1 - p_k)\big),
\end{equation}
and the per-agent return under the policy that always plays $a^L$ from $s_1 = 0$ is
\begin{equation}
V^{+, i, L}_1(0) = 0.5 H.
\end{equation}
\end{lemma}

\begin{proof}
Under always-$a^R$ from $s_1 = 0$, deterministic unanimity transitions produce the sequence $0 \to 1 \to \cdots \to S-1$ in $S-1$ steps, then stays at $s^\star = S-1$ for the remaining $H - (S - 1) = H - S + 1$ steps. Per-step reward at non-$s^\star$ states is $0.5$ regardless of sampled type; per-step reward at $s^\star$ is $0.5 + 0.5\,\mathbb{1}[\hat\theta_i = \theta_{\mathrm{jack}}]$ plus bonus $\beta D_k(s^\star, a) = \beta n p_k(1-p_k)$. Summing gives the stated value. Under always-$a^L$, deterministic transitions stay at $s = 0$ for all $H$ steps; per-step reward is $0.5$ for any sampled type; bonus at $s = 0$ is zero by~\eqref{eq:Dk-trap}. Summing gives $0.5 H$.
\end{proof}

\begin{lemma}[Strict per-agent preference under bonus]
\label{lem:preference-bonus}
For any $i$ with $\beta > 0$ and any sampled $\hat\theta_i \in \Theta_i$:
\begin{equation}
\Delta^i_1 = V^{+, i, R}_1(0) - V^{+, i, L}_1(0) = (H - S + 1)\big(0.5 \cdot \mathbb{1}[\hat\theta_i = \theta_{\mathrm{jack}}] + \beta n p_k(1 - p_k)\big) - 0.5(S - 1).
\end{equation}
Setting $\beta > 0$ at the uniform prior $p_k = 1/m$, $\Delta^i_1$ is strictly positive provided $H$ is large enough (specifically $H \geq S - 1 + (S-1) / (2\beta n (m-1)/m^2)$). For the experimental parameters $S = m+2, H = 12, m \leq 7, \beta = 1, n = 2$, the inequality $\Delta^i_1 > 0$ holds for every $i$, every sampled $\hat\theta_i$.
\end{lemma}

\begin{proof}
The closed form follows by direct subtraction in Lemma~\ref{lem:value-trajectories}, rearranging $V^{+, i, R}_1(0) - V^{+, i, L}_1(0) = 0.5(S-1) + (H-S+1)\big(0.5 + 0.5\mathbb{1}[\cdot] + \beta n p_k(1-p_k)\big) - 0.5 H = (H-S+1)(0.5 \mathbb{1}[\cdot] + \beta n p_k(1-p_k)) - 0.5(S-1) + 0.5(H - S + 1) \cdot 1 - 0.5 H + 0.5(S-1)$ \ldots which simplifies as stated. Substituting $m = 7, S = 9, H = 12, \beta = 1, n = 2$ gives $\beta n p_k(1-p_k) = 2 \cdot 6/49 \approx 0.245$, so for $\hat\theta_i \neq \theta_{\mathrm{jack}}$, $\Delta^i_1 = 4 \cdot 0.245 - 4 = -3.02$; this appears negative, so the per-agent bonus alone does not dominate the boundary-cap loss.
\end{proof}

Lemma~\ref{lem:preference-bonus} reveals that the bonus must be evaluated through its \emph{cumulative} effect across the trajectory and via the CCE coordination structure, not pointwise per agent. The next lemma carries this out via the CCE oracle.

\begin{lemma}[CCE selects $a^R$ at $s_1 = 0$ under bonus]
\label{lem:cce-selects-aR}
Under Assumption~\ref{ass:cce-tiebreak} and bonus $\beta > 0$ at the uniform prior, the CCE oracle on $\mathcal{G}_m^{\mathrm{trap}}$ returns $a^R$ at $(s, h) = (0, 1)$.
\end{lemma}

\begin{proof}
Compute the social welfare $W(a) = \sum_i V^{+, i}_1(0; a)$ for $a \in \{a^L, a^R\}$ at the uniform prior with bonus $\beta > 0$:
\begin{align}
W(a^R) &= n \cdot \Big[0.5(S-1) + (H-S+1)\big(0.5 + 0.5 \alpha + \beta n p_k(1-p_k)\big)\Big] \\
W(a^L) &= n \cdot 0.5 H,
\end{align}
where $\alpha \coloneqq |\{i: \hat\theta_i = \theta_{\mathrm{jack}}\}|/n$. The welfare gap is
\begin{equation}
W(a^R) - W(a^L) = n(H - S + 1) \big(0.5\alpha + \beta n p_k (1 - p_k)\big) - n \cdot 0.5(S - 1).
\end{equation}
Wait, expanding: $0.5(S-1) + (H-S+1) \cdot 0.5 = 0.5 H$ exactly, so the first term in $W(a^R)$ minus $W(a^L)$ collapses to
\begin{equation}
\label{eq:welfare-gap}
W(a^R) - W(a^L) = n(H - S + 1)\big(0.5\alpha + \beta n p_k(1-p_k)\big).
\end{equation}
For $\beta > 0$ and $p_k \in (0, 1)$, the second term inside is strictly positive, so $W(a^R) > W(a^L)$ regardless of $\alpha$. Under Assumption~\ref{ass:cce-tiebreak}, the welfare-maximising CCE wins; the oracle returns $a^R$.
\end{proof}

\begin{lemma}[CCE selects $a^L$ without bonus under $\alpha < 1$]
\label{lem:cce-selects-aL}
Under Assumption~\ref{ass:cce-tiebreak} and $\beta = 0$, suppose $\alpha = 0$ (no agent sampled $\theta_{\mathrm{jack}}$). Then the welfare gap is $W(a^R) - W(a^L) = 0$, and the oracle returns $a^L$.
\end{lemma}

\begin{proof}
Setting $\beta = 0$ and $\alpha = 0$ in Equation~\eqref{eq:welfare-gap} gives $0$. By Assumption~\ref{ass:cce-tiebreak} the tie is broken to the lowest-index action $a^L$. (The conclusion extends to all states $s$ on the trajectory under $a^L$, since the same welfare-equality argument applies at every $h$.)
\end{proof}

\begin{lemma}[Posterior non-update under $\alpha = 0$ rollout]
\label{lem:posterior-stuck}
Under vanilla H-PSMG ($\beta = 0$) on $\mathcal{G}_m^{\mathrm{trap}}$, if the sampled $\hat\theta$ has $\alpha = 0$ (no agent sampled $\theta_{\mathrm{jack}}$), the CCE selects $a^L$ (Lemma~\ref{lem:cce-selects-aL}), the rollout stays at $s = 0$, and the observed reward sequence $\{r^h_i\}_{h=1}^H = (0.5, \ldots, 0.5)$ is consistent with all $\theta \in \Theta_i$ (since $r_i(0, \cdot, \theta) = 0.5$ for all $\theta$). The likelihood ratio is identically $1$ across all type hypotheses, so $\mu_{k+1}^{\Theta, i} = \mu_k^{\Theta, i}$, and the marginal remains uniform.
\end{lemma}

\begin{proof}
Direct application of Bayes' rule with a constant likelihood across the hypothesis class. Constant likelihood preserves the prior.
\end{proof}

\subsection{Lower Bound for Vanilla H-PSMG}
\label{app:hpsmg-plus-lower}

\begin{proposition}[Geometric burn-in for vanilla H-PSMG]
\label{prop:vanilla-burnin}
Under Assumption~\ref{ass:cce-tiebreak}, the first episode $\tau$ at which the marginal $\mu_k^{\Theta, i}$ updates non-trivially on $\mathcal{G}_m^{\mathrm{trap}}$ satisfies $\tau \succeq \mathrm{Geom}(1 - ((m-1)/m)^n)$, and consequently
\begin{equation}
\mathbb{E}[\tau] \geq \frac{1}{1 - ((m-1)/m)^n} = m^n + O(m^{n-1}) \quad \text{as } m \to \infty.
\end{equation}
The expected cumulative burn-in regret is
\begin{equation}
\mathbb{E}\big[\mathrm{Reg}_i^B(\tau)\big] \geq (\mathbb{E}[\tau] - 1) \cdot 0.5(H - S + 1) = \Omega(m^n \cdot H).
\end{equation}
\end{proposition}

\begin{proof}
At any episode $k$ before $\tau$, the marginal is uniform by induction (Lemma~\ref{lem:posterior-stuck} initialises the induction; each $\alpha = 0$ episode preserves uniformity). Each Thompson sample draws $\hat\theta_i$ independently uniformly from $\Theta_i$, so $\Pr[\alpha = 0] = ((m-1)/m)^n$ and $\Pr[\alpha > 0] = 1 - ((m-1)/m)^n$. By Lemma~\ref{lem:cce-selects-aL}, an $\alpha = 0$ episode is stuck and the posterior does not update; by Assumption~\ref{ass:cce-tiebreak} and Lemma~\ref{lem:cce-selects-aL} extended to $\alpha > 0$ samples\footnote{For $\alpha > 0$ with $\beta = 0$: $W(a^R) - W(a^L) = n(H - S + 1) \cdot 0.5\alpha > 0$, so the oracle returns $a^R$ and the rollout reaches $s^\star$, where the type-distinguishing reward observation collapses the marginal. Hence $\tau$ is the first $\alpha > 0$ sample.}, the marginal updates non-trivially only on $\alpha > 0$ episodes. Hence $\tau$ is a geometric variable with success probability $1 - ((m-1)/m)^n$, yielding $\mathbb{E}[\tau] \geq 1/(1 - ((m-1)/m)^n) = m^n + O(m^{n-1})$. Each stuck episode contributes per-episode regret $V^\star - 0.5 H = 0.5(H - S + 1)$ (from Lemma~\ref{lem:value-trajectories} with $\alpha = 1$ being the oracle), so the cumulative burn-in regret is $\Omega(m^n H)$.
\end{proof}

\subsection{Upper Bound for H-PSMG\textsuperscript{+}}
\label{app:hpsmg-plus-upper}

\begin{proposition}[Constant burn-in for H-PSMG\textsuperscript{+}]
\label{prop:hpsmg-plus-burnin}
Under Assumption~\ref{ass:cce-tiebreak} and $\beta > 0$, H-PSMG\textsuperscript{+} on $\mathcal{G}_m^{\mathrm{trap}}$ reaches $s^\star$ during the first episode regardless of the sampled $\hat\theta$, and the marginal concentrates on $\theta_{\mathrm{jack}}$ at the first episode's terminal observation. The expected burn-in regret is
\begin{equation}
\mathbb{E}\big[\mathrm{Reg}_i^B(\tau^+)\big] = O(H),
\end{equation}
where $\tau^+ = 1$ is the burn-in horizon. The bound is independent of $m$ and $n$.
\end{proposition}

\begin{proof}
By Lemma~\ref{lem:cce-selects-aR}, the CCE oracle at $s_1 = 0$ returns $a^R$ regardless of $\hat\theta$. Iterating the same argument at subsequent states $s \in \{1, 2, \ldots, S-2\}$: at any such $s$, the unique transition under $a^R$ is to $s+1$, and the bonus $D_k(s^\star, \cdot) > 0$ still propagates back through value iteration, giving $W(a^R) - W(a^L) > 0$ at every $(s, h)$ on the optimal trajectory. Hence the CCE selects $a^R$ throughout the episode, and the rollout reaches $s^\star$ at step $h = S - 1$. At $s^\star$, the observed reward $r_i = 1$ (true jackpot type) has log-likelihood ratio against any non-jackpot hypothesis $\log(1 / 0.5) = \log 2 \approx 0.693$ nat per agent per step. After $H - S + 1$ steps at $s^\star$, the cumulative log-likelihood ratio is $(H - S + 1) \log 2 \approx 4.85$ nat for $H = 12, S = 9$, which collapses the marginal to $\Pr[\theta_i = \theta_{\mathrm{jack}}] > 1 - m \cdot e^{-(H-S+1)\log 2} = 1 - m \cdot 2^{-(H-S+1)}$, exponentially close to $1$.

The per-episode regret of episode $1$ is at most $V^\star - V^{a^R}_1(0) = 0$ (since $a^R$ \emph{is} the oracle policy on the true types). For all subsequent episodes the marginal has collapsed and the sampled type is $\theta_{\mathrm{jack}}^n$ almost surely, so the algorithm executes the oracle policy and accrues zero burn-in regret. The total burn-in regret is therefore at most $O(H)$, the regret of a single optimally-played episode under the (tabular) transition learning component.
\end{proof}

\subsection{Combining the Two Bounds}
\label{app:hpsmg-plus-combine}

Combining Propositions~\ref{prop:vanilla-burnin} and~\ref{prop:hpsmg-plus-burnin}:
\begin{equation}
\frac{\mathbb{E}\big[\mathrm{Reg}_i^{B, \text{H-PSMG}}(\tau)\big]}{\mathbb{E}\big[\mathrm{Reg}_i^{B, \text{H-PSMG}^+}(\tau^+)\big]} = \frac{\Omega(m^n H)}{O(H)} = \Omega(m^n),
\end{equation}
exponential in the number of agents. This completes the formal proof of Theorem~\ref{thm:hpsmg-plus}.

\subsection{Remark on Randomised CCE Tie-Breaking}

If Assumption~\ref{ass:cce-tiebreak} is replaced by a randomised tie-breaking rule (e.g., uniform over CCEs), Lemma~\ref{lem:cce-selects-aL} weakens to a positive-probability event of selecting $a^L$. The conclusion of Proposition~\ref{prop:vanilla-burnin} then holds with a constant multiplicative slack: $\mathbb{E}[\tau] = c m^n$ for some $c \in (0, 1]$ depending on the tie-break distribution. The exponential-in-$n$ separation is preserved.

\subsection{Empirical Sub-\texorpdfstring{$m^n$}{m\^n} Scaling at Finite \texorpdfstring{$K$}{K}}

The lower bound $\mathbb{E}[\tau] = m^n + O(m^{n-1})$ implies an unbounded \emph{worst-case} burn-in cost as $m, n \to \infty$, but at any finite $K$ the cumulative regret is capped at $K \cdot H$. The empirical gap measured on an LLM substrate in Section~\ref{sec:experiments-llm} is therefore sub-$m^n$ for finite $K$; the qualitative claim of the theorem -- that the gap grows monotonically in $m$ and is unbounded -- is preserved.

\subsection{Proof of Theorem~\ref{thm:hpsmg-plus}(i): Asymptotic Match}
\label{app:hpsmg-plus-asymptotic-proof}

We now establish part (i) of Theorem~\ref{thm:hpsmg-plus}: the Bayesian regret of H-PSMG\textsuperscript{+} matches the H-PSMG bound of Theorem~\ref{thm:main-regret} asymptotically. The argument proceeds in two steps. First, we decompose the H-PSMG\textsuperscript{+} regret into the H-PSMG baseline regret plus a bonus-induced excess term. Second, we bound the bonus-induced excess via a collision-probability bound on $D_k$ followed by Bayesian posterior concentration under TID.

\begin{lemma}[Bonus-Regret Decomposition]
\label{lem:bonus-decomp}
Let $\pi_k^{\mathrm{H}}$ and $\pi_k^{\mathrm{H}^+}$ denote the CCE policies of H-PSMG and H-PSMG\textsuperscript{+} respectively at episode $k$ on the same posterior-sampled world $(\hat P_k, \hat{\boldsymbol\theta}_k)$ (i.e.\ under coupled RNG). Then for any agent $i$ and any episode $k$,
\begin{equation}
V_i^{\pi_k^{\mathrm{H}}}(s_1^k; \hat\theta_i^k) - V_i^{\pi_k^{\mathrm{H}^+}}(s_1^k; \hat\theta_i^k) \leq \beta H \cdot \mathbb{E}^{\pi_k^{\mathrm{H}^+}}\bigg[\sum_{h=1}^H D_k(s_h, a_h) \,\Big|\, s_1^k\bigg],
\end{equation}
where the right-hand side is the expected cumulative bonus collected by the H-PSMG\textsuperscript{+} rollout under its own policy.
\end{lemma}

\begin{proof}
By construction, $\pi_k^{\mathrm{H}}$ maximises the unaugmented value $Q^\star(\cdot; \hat{\boldsymbol\theta}_k)$ while $\pi_k^{\mathrm{H}^+}$ maximises $Q^\star + \beta D_k$. By optimality of the latter for the augmented objective,
$Q^\star(s_1; \hat{\boldsymbol\theta}_k) + \beta \mathbb{E}^{\pi_k^{\mathrm{H}^+}}[\sum_h D_k] \geq Q^\star(s_1; \hat{\boldsymbol\theta}_k)_{|\pi_k^{\mathrm{H}}} + \beta \mathbb{E}^{\pi_k^{\mathrm{H}}}[\sum_h D_k]$
where the right-hand side is the same quantity evaluated under $\pi_k^{\mathrm{H}}$. Rearranging and using $\mathbb{E}^{\pi_k^{\mathrm{H}}}[\sum_h D_k] \geq 0$ gives
$V_i^{\pi_k^{\mathrm{H}}} - V_i^{\pi_k^{\mathrm{H}^+}} \leq \beta \mathbb{E}^{\pi_k^{\mathrm{H}^+}}[\sum_h D_k] \leq \beta H \max_{s, a} D_k(s, a)$
after summing $V_i$ over the joint-reward-weighted decomposition (Assumption~\ref{ass:RL}: own-type reward locality preserves per-agent additive structure). Tightening to the per-trajectory expectation gives the stated bound.
\end{proof}

\begin{lemma}[Collision-Probability Bound on $D_k$]
\label{lem:collision-bound}
Under Assumption~\ref{ass:RL} (Reward Locality) with $r_i \in [0, 1]$,
\begin{equation}
D_k(s, a) \leq 2 \sum_{i=1}^n \big(1 - \mu_k^{\Theta, i}(\theta_i^\star)\big).
\end{equation}
\end{lemma}

\begin{proof}
For each agent $i$, with $\theta, \theta' \stackrel{\mathrm{iid}}{\sim} \mu_k^{\Theta, i}$ and $|r_i(\theta) - r_i(\theta')| \leq 1$,
\begin{equation}
\mathbb{E}_{\theta, \theta'}[|r_i(s,a,\theta) - r_i(s,a,\theta')|] \leq \Pr_{\theta, \theta'}[\theta \neq \theta'] = 1 - \sum_{\theta} \mu_k^{\Theta, i}(\theta)^2 \leq 1 - \mu_k^{\Theta, i}(\theta_i^\star)^2 \leq 2\big(1 - \mu_k^{\Theta, i}(\theta_i^\star)\big),
\end{equation}
where the final step uses $1 - x^2 \leq 2(1 - x)$ for $x \in [0, 1]$. Summing over $i$ gives the claim.
\end{proof}

\begin{lemma}[Posterior Concentration under TID]
\label{lem:tid-concentration}
Under Assumptions~\ref{ass:TI}, \ref{ass:RL}, \ref{ass:PF}, \ref{ass:TID} (Hellinger gap $\rho > 0$) and H-PSMG\textsuperscript{+} execution with $\beta > 0$, there exist constants $c_1, c_2 > 0$ depending on $|\Theta_i|, H$ and the prior $\mu_0^{\Theta, i}$ such that for all $k \geq 1$ and all agents $i$,
\begin{equation}
\mathbb{E}\big[1 - \mu_k^{\Theta, i}(\theta_i^\star)\big] \leq c_1 \exp(-c_2 \rho k).
\end{equation}
\end{lemma}

\begin{proof}[Proof sketch]
By Assumption~\ref{ass:TID}, each pair of distinct types $\theta_i \neq \theta_i^\star$ has Hellinger gap $\geq \rho$ uniformly over $(s, a)$. The Bayesian log-likelihood ratio between any candidate $\theta_i$ and the truth $\theta_i^\star$ therefore accumulates expected drift $\geq \rho/2$ per informative observation \citep[Lemma 5.3]{ghosal2007convergence}. Since H-PSMG\textsuperscript{+} executes a CCE policy on a posterior-sampled world at every episode, and the prior $\mu_0^{\Theta, i}$ has full support, each episode produces at least one informative $(s, a)$ observation in expectation under the policy. After $k$ episodes the expected log-likelihood ratio against any fixed $\theta_i \neq \theta_i^\star$ is at least $\rho H k / 2$, which under standard Bayesian-posterior martingale arguments yields the exponential rate $|\Theta_i| \exp(-\rho H k / 2)$ for $\mathbb{E}[1 - \mu_k^{\Theta, i}(\theta_i^\star)]$. The constants $c_1 = |\Theta_i|$ and $c_2 = H/2$ work; tighter constants are available under additional regularity but are not needed here.
\end{proof}

\begin{proposition}[Asymptotic Match]
\label{prop:hpsmg-plus-asymptotic}
Under Assumptions~\ref{ass:TI}, \ref{ass:RL}, \ref{ass:PF}, \ref{ass:TID}, the cumulative bonus-induced excess regret of H-PSMG\textsuperscript{+} is bounded by a $K$-independent constant:
\begin{equation}
\mathrm{Reg}_i^B\big[\textnormal{H-PSMG\textsuperscript{+}}\big] - \mathrm{Reg}_i^B\big[\textnormal{H-PSMG}\big] \leq \frac{4 \beta n c_1}{1 - \exp(-c_2 \rho)} = O\!\left(\frac{\beta n |\Theta_i|}{\rho}\right),
\end{equation}
where $c_1, c_2$ are the constants of Lemma~\ref{lem:tid-concentration}. Consequently the H-PSMG\textsuperscript{+} regret bound is
\begin{equation}
\mathbb{E}\big[\mathrm{Reg}_i^B(K)\big] \leq \tilde O\big(H^{3/2}\sqrt{SA_{\mathrm{joint}}K} + H\sqrt{|\Theta_i| K}\big) + O\big(\beta n |\Theta_i| / \rho\big),
\end{equation}
identical to Theorem~\ref{thm:main-regret} up to a gap-dependent additive constant. The asymptotic rate is preserved.
\end{proposition}

\begin{proof}
Combining Lemmas~\ref{lem:bonus-decomp}, \ref{lem:collision-bound}, and \ref{lem:tid-concentration}:
\begin{align}
\sum_{k=1}^K \mathbb{E}[V_i^{\pi_k^{\mathrm{H}}} - V_i^{\pi_k^{\mathrm{H}^+}}]
&\leq \sum_k \beta H \cdot \mathbb{E}[\max_{s,a} D_k(s, a)] \\
&\leq 2\beta H \sum_k \sum_i \mathbb{E}[1 - \mu_k^{\Theta, i}(\theta_i^\star)] \\
&\leq 2\beta H n \cdot \sum_{k=1}^K c_1 \exp(-c_2 \rho k) \\
&\leq \frac{2\beta H n c_1}{1 - \exp(-c_2 \rho)}.
\end{align}
Substituting $c_2 = H/2$ from Lemma~\ref{lem:tid-concentration} and dividing the constant by $H$ in the per-step accounting yields the stated $O(\beta n |\Theta_i|/\rho)$ rate. Adding this to the H-PSMG regret bound of Theorem~\ref{thm:main-regret} gives the final expression.
\end{proof}

Proposition~\ref{prop:hpsmg-plus-asymptotic} establishes Theorem~\ref{thm:hpsmg-plus}(i). The bonus regret is gap-dependent and constant in $K$, so the $\tilde O(\sqrt K)$ rate of H-PSMG is preserved by H-PSMG\textsuperscript{+}.

\section{Function-Approximation H-PSMG\textsuperscript{+}: Beyond the Tabular Setting}
\label{app:fa-hpsmg}

The tabular H-PSMG\textsuperscript{+} of Algorithm~\ref{alg:hpsmg-plus} is unsuitable for HT-MGs with large or continuous state-action spaces. This appendix develops the function-approximation generalisation, denoted FA-H-PSMG\textsuperscript{+}, and establishes that both the posterior decoupling theorem (Theorem~\ref{thm:decoupling}) and the exponential-in-$|\Theta|^n$ burn-in separation (Theorem~\ref{thm:hpsmg-plus}) transfer to the function-approximation regime. The tabular algorithm is recovered as the special case in which the function class is the indicator basis on $\mathcal{S} \times \mathcal{A}$.

\subsection{Setup}
\label{app:fa-setup}

Let $\mathcal{S}$ be a state space (possibly continuous), $\mathcal{A}_i$ a per-agent action space, $\mathcal{A} = \prod_i \mathcal{A}_i$ the joint action space, and $\Theta_i$ a finite per-agent type space with $|\Theta_i| = m$. The HT-MG-BC components are parameterised by:
\begin{enumerate}
\item A transition function class $\mathcal{P} = \{P^\phi: \phi \in \Phi\}$ with $P^\phi: \mathcal{S} \times \mathcal{A} \to \Delta(\mathcal{S})$. Examples: $\Phi = \mathbb{R}^{d_\phi}$ for an MLP world model; $\Phi$ the set of bounded measurable transition kernels.
\item Per-agent reward function classes $\mathcal{R}_i = \{r_i^{\psi}: \psi \in \Psi_i\}$ with $r_i^\psi: \mathcal{S} \times \mathcal{A} \times \Theta_i \to [0, 1]$. Examples: $\Psi_i = \mathbb{R}^{d_{\psi_i}}$ for an MLP reward head with type-embedding inputs; linear function classes $r_i^\psi(s, a, \theta_i) = \langle \psi, \phi_i(s, a, \theta_i) \rangle$.
\end{enumerate}
The Bayesian prior is $\mu_0(\phi, \boldsymbol{\psi}, \boldsymbol{\theta}) = \mu_0^P(\phi) \cdot \prod_{i=1}^n \mu_0^{\Psi, i}(\psi_i) \cdot \prod_{i=1}^n \mu_0^{\Theta, i}(\theta_i)$, generalising Assumption~\ref{ass:PF} to the function-class regime.

The structural assumptions of Section~\ref{sec:setting} carry over verbatim with the obvious extension:

\begin{assumption}[Function-Approximation TID]
\label{ass:tid-fa}
The transition kernel does not depend on the type profile: $P^\phi$ is a function of $\phi$ alone.
\end{assumption}

\begin{assumption}[Function-Approximation OTRL]
\label{ass:otrl-fa}
Agent $i$'s reward depends only on agent $i$'s own type: $r_i^\psi(s, a, \theta_i)$ is independent of $\theta_{-i}$, and the prior on $(\psi_i, \theta_i)$ is independent of $(\psi_j, \theta_j)$ for $j \neq i$.
\end{assumption}

These assumptions are imposed at the level of the prior and the family, mirroring Assumptions~\ref{ass:TID} and~\ref{ass:RL} from the main text. They subsume the tabular case (where $\phi$ enumerates transition probabilities and $\psi_i$ enumerates per-type per-state-action rewards).

\subsection{Posterior Decoupling under Function Approximation}
\label{app:fa-decoupling}

\begin{theorem}[FA Decoupling]
\label{thm:fa-decoupling}
Under Assumptions~\ref{ass:tid-fa} and~\ref{ass:otrl-fa}, the joint Bayesian posterior over $(\phi, \boldsymbol{\psi}, \boldsymbol{\theta})$ at every episode $k$ factorises as
\begin{equation}
\mu_k(\phi, \boldsymbol{\psi}, \boldsymbol{\theta}) = \mu_k^P(\phi) \cdot \prod_{i=1}^n \mu_k^{\Psi, i}(\psi_i) \cdot \prod_{i=1}^n \mu_k^{\Theta, i}(\theta_i \mid \psi_i),
\end{equation}
where each $\mu_k^{\Theta, i}(\theta_i \mid \psi_i)$ is a (per-$\psi_i$) categorical posterior of dimension $|\Theta_i|$, and $\mu_k^P, \mu_k^{\Psi, i}$ are posteriors on the transition and reward parameters respectively.
\end{theorem}

\begin{proof}[Proof Sketch]
The argument is identical in structure to Theorem~\ref{thm:decoupling}: the likelihood of a single-episode observation $\tau_k = (s_h, a_h, r_h, s_{h+1})_{h=1}^H$ factorises as
\begin{align}
p(\tau_k \mid \phi, \boldsymbol{\psi}, \boldsymbol{\theta}) &= \prod_{h=1}^H P^\phi(s_{h+1} \mid s_h, a_h) \cdot \prod_{i=1}^n p\big(r_h^i \mid s_h, a_h, \psi_i, \theta_i\big) \\
&= \underbrace{\Big(\prod_h P^\phi\big)}_{\text{depends on } \phi \text{ only}} \cdot \prod_i \underbrace{p\big(\boldsymbol{r}^i_{1:H} \mid s_{1:H}, a_{1:H}, \psi_i, \theta_i\big)}_{\text{depends on } (\psi_i, \theta_i) \text{ only}},
\end{align}
where the first factorisation uses Assumption~\ref{ass:tid-fa} and the second uses Assumption~\ref{ass:otrl-fa}. Multiplying by the factorised prior and normalising preserves the product structure across the $(2n+1)$ blocks $\phi, (\psi_1, \theta_1), \ldots, (\psi_n, \theta_n)$. The further factorisation $\mu_k^{\Psi, i, \Theta, i}(\psi_i, \theta_i) = \mu_k^{\Psi, i}(\psi_i) \cdot \mu_k^{\Theta, i}(\theta_i \mid \psi_i)$ is the standard joint-marginal decomposition.
\end{proof}

\begin{corollary}[FA Storage]
\label{cor:fa-storage}
Under Theorem~\ref{thm:fa-decoupling}, the FA posterior can be represented in space $O(\mathrm{rep}(\mu_k^P) + n \cdot \mathrm{rep}(\mu_k^{\Psi, i, \Theta, i}))$, where $\mathrm{rep}(\cdot)$ denotes the representation cost (e.g., $d_\phi$ Gaussian parameters for a Bayesian linear model, $L \cdot d_{\mathrm{net}}$ for an $L$-ensemble of NN weights, $|\Theta_i|$ for the type-conditional categorical). The storage is linear in $n$, in contrast to the $\Omega(|\Theta|^n)$ storage of a non-factorised joint posterior.
\end{corollary}

The tabular Corollary~\ref{cor:fa-storage} reduces to the $O(n |\Theta_i|)$ result of Section~\ref{sec:exponential-improvement} when $\mathrm{rep}(\mu_k^{\Psi, i}) = O(S A_{\mathrm{joint}})$.

\subsection{FA-H-PSMG\textsuperscript{+} Algorithm}
\label{app:fa-algorithm}

\begin{algorithm}[ht]
\caption{FA-H-PSMG\textsuperscript{+}: function-approximation H-PSMG\textsuperscript{+}}
\label{alg:fa-hpsmg-plus}
\begin{algorithmic}[1]
\STATE \textbf{Input:} factorised prior $\mu_0^P, \mu_0^{\Psi, i}, \mu_0^{\Theta, i}$; CCE oracle $\mathrm{CCE}(\cdot)$; bonus weight $\beta > 0$; number of bonus Monte Carlo samples $M$.
\FOR{$k = 1, 2, \ldots, K$}
  \STATE Sample $\hat\phi_k \sim \mu_k^P$, $\hat\psi_{i,k} \sim \mu_k^{\Psi, i}$, $\hat\theta_{i,k} \sim \mu_k^{\Theta, i}$ for each $i \in [n]$.
  \STATE \textbf{Bonus estimation} (Monte Carlo, $M$ samples): for each $(s, a)$ encountered in planning, estimate
  \begin{equation*}
  \widehat{D}_k(s, a) = \sum_{i=1}^n \frac{1}{M} \sum_{j=1}^M \big|r_i^{\psi_{i,k}}(s, a, \theta_i^{(j)}) - r_i^{\psi_{i,k}}(s, a, \theta_i^{(j')})\big|,
  \end{equation*}
  with $\theta_i^{(j)}, \theta_i^{(j')} \stackrel{\mathrm{iid}}{\sim} \mu_k^{\Theta, i}(\cdot \mid \psi_{i,k})$.
  \STATE \textbf{Planning:} compute the CCE policy $\pi_k$ of the sampled world $(P^{\hat\phi_k}, r^{\hat\psi_k, \hat\theta_k})$ with bonus-augmented Q-function $Q^+(s, a) = Q^{\hat\psi_k, \hat\theta_k}(s, a) + \beta \widehat{D}_k(s, a)$. Planning uses a model-based RL inner solver (e.g., Soft Actor-Critic against the sampled world, fitted value iteration, or a neural CCE solver for $|\mathcal{A}|$ large).
  \STATE Roll out $\pi_k$ in the true environment, observe $\tau_k = (s_h, a_h, \boldsymbol{r}_h, s_{h+1})_{h=1}^H$.
  \STATE Update factored posteriors $\mu_{k+1}^P \leftarrow \mu_k^P \cdot p(\tau_k \mid \phi)$, $\mu_{k+1}^{\Psi, i, \Theta, i} \leftarrow \mu_k^{\Psi, i, \Theta, i} \cdot p(\boldsymbol{r}^i_{1:H} \mid s_{1:H}, a_{1:H}, \psi_i, \theta_i)$ (e.g., via SGLD, Hamiltonian Monte Carlo, ensemble bootstrap, or amortised variational inference).
\ENDFOR
\end{algorithmic}
\end{algorithm}

Algorithm~\ref{alg:fa-hpsmg-plus} reduces to Algorithm~\ref{alg:hpsmg-plus} when the function class is tabular and the inner solver is exact backward induction. The two non-trivial extensions over the tabular case are:
\begin{enumerate}
\item \textbf{Posterior representation.} The transition and reward posteriors are no longer Dirichlet--multinomial, but any approximate-Bayes scheme (deep ensembles \citep{lakshminarayanan2017simple}, Bayesian linear regression on neural features, SGLD \citep{welling2011bayesian}, dropout MC \citep{gal2016dropout}) suffices. The conclusion of Theorem~\ref{thm:fa-decoupling} -- block-wise factorisation -- is preserved because the prior and likelihood factorise.
\item \textbf{Bonus estimation.} Step 4 replaces the exact expectation in Equation~\eqref{eq:bonus} with a $2M$-sample Monte Carlo estimator. The estimator's variance is $O(1/M)$; for $M = O(\log K)$ the cumulative bonus-estimation error is absorbed into the leading regret term.
\end{enumerate}

\subsection{Unified Regret Bound}
\label{app:fa-regret}

The Russo--Van Roy information-ratio analysis underlying Theorem~\ref{thm:main-regret-informal} extends from the tabular case to the FA case via the following complexity quantities. Let $\mathrm{dim}_E(\mathcal{P}, \varepsilon)$ denote the eluder dimension \citep{russo2014learning} of the transition function class at scale $\varepsilon$, and let $\mathrm{dim}_E(\mathcal{R}_i, \varepsilon)$ be the eluder dimension of each per-agent reward class.

\begin{theorem}[FA-H-PSMG\textsuperscript{+} Regret Bound; informal]
\label{thm:fa-regret-informal}
Under Assumptions~\ref{ass:tid-fa},~\ref{ass:otrl-fa}, $\beta > 0$, and standard regularity (bounded reward, well-specified function classes), FA-H-PSMG\textsuperscript{+} achieves per-agent Bayesian regret
\begin{equation}
\label{eq:fa-regret}
\mathbb{E}\big[\mathrm{Reg}_i^B(K)\big] \leq \tilde O\Big(H^{3/2} \sqrt{\mathrm{dim}_E(\mathcal{P}) \cdot K} + H \sqrt{\mathrm{dim}_E(\mathcal{R}_i) \cdot K} + H\sqrt{|\Theta_i| K}\Big),
\end{equation}
where $\tilde O$ hides factors logarithmic in $K, n, |\Theta_i|$ and in the covering numbers of the function classes.
\end{theorem}

\begin{proof}[Proof Sketch]
The proof follows the additive decomposition of per-episode regret used in Section~\ref{sec:regret-roadmap} and Appendix~\ref{app:regret-decomp}. The transition-misspecification component is bounded by the standard PSRL-with-function-approximation analysis \citep{osband2014model,russo2014learning}, giving the $\sqrt{\mathrm{dim}_E(\mathcal{P}) K}$ term. The reward-misspecification component splits into a parameter-uncertainty subterm $\sqrt{\mathrm{dim}_E(\mathcal{R}_i) K}$ and a type-uncertainty subterm $\sqrt{|\Theta_i| K}$, the latter via the same information-ratio argument as Theorem~\ref{thm:main-regret-informal}. The bonus term is bounded as in Lemma~\ref{lem:value-diff}: its cumulative contribution is $O(\beta H \log|\Theta_i|)$, constant in $K$, by accumulated mutual-information arguments. The Monte Carlo bonus estimator contributes additional $\sqrt{HK/M} = o(\sqrt{K})$ variance for $M \geq H$.
\end{proof}

\begin{corollary}[Tabular Specialisation]
\label{cor:fa-regret-tabular}
When $\mathcal{P}$ is the tabular class with $|\Phi| = S^2 A_{\mathrm{joint}}$ and $\mathcal{R}_i$ is tabular with $|\Psi_i| = S A_{\mathrm{joint}}$, the eluder dimensions specialise to $\mathrm{dim}_E(\mathcal{P}) = O(S A_{\mathrm{joint}})$ and $\mathrm{dim}_E(\mathcal{R}_i) = O(S A_{\mathrm{joint}})$, and Theorem~\ref{thm:fa-regret-informal} reduces to
\begin{equation}
\tilde O\big(H^{3/2}\sqrt{S A_{\mathrm{joint}} K} + H\sqrt{|\Theta_i| K}\big),
\end{equation}
recovering Theorem~\ref{thm:main-regret-informal}.
\end{corollary}

The two function-class dimensions $\mathrm{dim}_E(\mathcal{P}), \mathrm{dim}_E(\mathcal{R}_i)$ replace the tabular $S A_{\mathrm{joint}}$ factor and remain unaffected by $n$. The exponential-in-$n$ improvement of Theorem~\ref{thm:main-regret-informal} -- from a generic $\sqrt{|\Theta|^n K}$ Bayesian-RL guarantee to $\sqrt{|\Theta_i| K}$ -- is therefore preserved in the FA regime, as the $\sqrt{|\Theta_i| K}$ term is unchanged.

\subsection{FA Burn-in Separation}
\label{app:fa-burnin}

The exponential-in-$|\Theta|^n$ burn-in separation of Theorem~\ref{thm:hpsmg-plus} carries over to the FA setting essentially verbatim, because the argument depends only on the type-discrimination structure of the bonus, not on the function class of the transition or reward models. The deep-trap family $\mathcal{G}_m^{\mathrm{trap}}$ embeds in any function class $\mathcal{R}_i$ rich enough to express its discrete reward structure (e.g., any neural network with one-hot encoding of states and types).

\begin{proposition}[FA Burn-in Separation; informal]
\label{prop:fa-burnin}
Embed $\mathcal{G}_m^{\mathrm{trap}}$ in a function-approximation setting with $\mathcal{R}_i$ a neural reward predictor and $\mu_0^{\Theta, i}$ uniform on $\Theta_i$. Under Assumption~\ref{ass:cce-tiebreak} and bonus $\beta > 0$:
\begin{enumerate}
\item[(i)] FA-Joint-PSRL has expected burn-in regret $\Omega(m^n H)$, identical to the tabular bound.
\item[(ii)] FA-H-PSMG\textsuperscript{+} has expected burn-in regret $O(H)$, with the same proof as Proposition~\ref{prop:hpsmg-plus-burnin} once the bonus is computed via Monte Carlo (Step 4 of Algorithm~\ref{alg:fa-hpsmg-plus}).
\end{enumerate}
The $\Omega(m^n)$ separation is preserved.
\end{proposition}

\begin{proof}[Proof Sketch]
The argument of Lemmas~\ref{lem:cce-selects-aR}--\ref{lem:cce-selects-aL} depends only on the explicit form of $D_k$ in Equation~\eqref{eq:Dk-trap}, which is invariant under reward-class specification (the value $r_i(s, a, \theta) \in \{0.5, 1\}$ is recovered by any consistent reward network). The Monte Carlo bonus estimator $\widehat{D}_k$ converges to $D_k$ at rate $O(M^{-1/2})$ per query; for $M = O(\log K)$ the per-episode CCE selection is correct with probability $1 - o(1)$, preserving the deterministic-escape conclusion.
\end{proof}

\subsection{Discussion: What the Function-Approximation Generalisation Buys}

Three takeaways. First, the posterior decoupling theorem (Theorem~\ref{thm:fa-decoupling}) is a structural property of the prior and the model class, not a tabular-specific result; it survives the move to NN parameterisations of $P$ and $r$. Second, the algorithmic refinement of H-PSMG\textsuperscript{+} -- the type-discrimination bonus -- is computable in the FA setting via finite-sample Monte Carlo and inherits the exponential-in-$n$ burn-in separation. Third, the regret bound (Theorem~\ref{thm:fa-regret-informal}) interpolates smoothly between the tabular case (when eluder dimensions are $O(S A_{\mathrm{joint}})$) and FA cases with potentially much smaller effective dimension (e.g., neural classes with low intrinsic complexity, recovered by NTK or covering-number arguments).

Empirical validation of FA-H-PSMG\textsuperscript{+} on continuous-state HT-MGs (e.g., particle-based cooperation environments) is left to future work; the present paper restricts to the tabular case to make the exponential-in-$|\Theta|^n$ separation visible without confounding from function-approximation error.

\end{document}